%File: formatting-instructions-latex-2025.tex
%release 2025.0
\documentclass[letterpaper]{article} % DO NOT CHANGE THIS
\usepackage{aaai25}  % DO NOT CHANGE THIS
\usepackage{times}  % DO NOT CHANGE THIS
\usepackage{helvet}  % DO NOT CHANGE THIS
\usepackage{courier}  % DO NOT CHANGE THIS
\usepackage[hyphens]{url}  % DO NOT CHANGE THIS
\usepackage{graphicx} % DO NOT CHANGE THIS
\urlstyle{rm} % DO NOT CHANGE THIS
\def\UrlFont{\rm}  % DO NOT CHANGE THIS
\usepackage{natbib}  % DO NOT CHANGE THIS AND DO NOT ADD ANY OPTIONS TO IT
\usepackage{caption} % DO NOT CHANGE THIS AND DO NOT ADD ANY OPTIONS TO IT
\frenchspacing  % DO NOT CHANGE THIS
\setlength{\pdfpagewidth}{8.5in}  % DO NOT CHANGE THIS
\setlength{\pdfpageheight}{11in}  % DO NOT CHANGE THIS
%

% OUR PACKAGES VV
\usepackage{afterpage}

\usepackage{amsmath}
\usepackage{amsthm}
\usepackage{pgfplots}
\usepackage{tikz}

\usepackage{subcaption}
\usepackage{lipsum}

\usepackage{xcolor}

%\usepackage{hyperref}

\usetikzlibrary{calc, angles, quotes}

%OUR PACKAGES ^^
%OUR DEFINITIONS VV
\newcommand{\sumgraph}{\underset{n \in [N]}{\sum}}

\def \sumgraphnotu {\underset{v \neq u }{ \sum}}

\def \sumgraphu {\underset{n\in [N]}{ \sum}}

\def \sumneighbors {\underset {m\in{\mathcal{N}(n)}}{\sum}}

\def \sumstrangers{\underset {m\notin{\mathcal{N}(n)}}{\sum}}

\newcommand{\flf}{\mathbf{f}_n^\top   \mathbf{L} \mathbf{f}_m}

\def \szemerdi{Szemerédi}
\def \PieClam \textsc{\large P\small I\small E\large C\small L\small A\small M}
\def \pclam \textsc{\large P\large C\small L\small A\small M}
\newtheorem{theorem}{Theorem}
\newtheorem{lemma}[theorem]{Lemma}
\newtheorem{claim}[theorem]{Claim}
\newtheorem{definition}[theorem]{Definition}
\newtheorem{remark}[theorem]{Remark}

% OUR DEFINITIONS ^^
% These are recommended to typeset algorithms but not required. See the subsubsection on algorithms. Remove them if you don't have algorithms in your paper.
\usepackage{algorithm}
\usepackage{algorithmic}

%
% These are are recommended to typeset listings but not required. See the subsubsection on listing. Remove this block if you don't have listings in your paper.
\usepackage{newfloat}
\usepackage{listings}
\usepackage{color}
\DeclareCaptionStyle{ruled}{labelfont=normalfont,labelsep=colon,strut=off} % DO NOT CHANGE THIS
\lstset{%
	basicstyle={\footnotesize\ttfamily},% footnotesize acceptable for monospace
	numbers=left,numberstyle=\footnotesize,xleftmargin=2em,% show line numbers, remove this entire line if you don't want the numbers.
	aboveskip=0pt,belowskip=0pt,%
	showstringspaces=false,tabsize=2,breaklines=true}
\floatstyle{ruled}
\newfloat{listing}{tb}{lst}{}
\floatname{listing}{Listing}
%
% Keep the \pdfinfo as shown here. There's no need
% for you to add the /Title and /Author tags.
\pdfinfo{
/TemplateVersion (2025.1)
}

\usepackage{amsfonts}


\newcommand{\ip}[2]{\left\langle#1,#2\right\rangle}


\newcommand\norm[1]{\left\lVert#1\right\rVert}


\newcommand{\ron}[1]{{{\textcolor{red}{\textbf{[Ron:} {#1}\textbf{]}}}}}
\newcommand{\dan}[1]{{{\textcolor{blue}{ {#1}\textbf{}}}}}

% DISALLOWED PACKAGES
% \usepackage{authblk} -- This package is specifically forbidden
% \usepackage{balance} -- This package is specifically forbidden
% \usepackage{color (if used in text)
% \usepackage{CJK} -- This package is specifically forbidden
% \usepackage{float} -- This package is specifically forbidden
% \usepackage{flushend} -- This package is specifically forbidden
% \usepackage{fontenc} -- This package is specifically forbidden
% \usepackage{fullpage} -- This package is specifically forbidden
% \usepackage{geometry} -- This package is specifically forbidden
% \usepackage{grffile} -- This package is specifically forbidden
% \usepackage{hyperref} -- This package is specifically forbidden
% \usepackage{navigator} -- This package is specifically forbidden
% (or any other package that embeds links such as navigator or hyperref)
% \indentfirst} -- This package is specifically forbidden
% \layout} -- This package is specifically forbidden
% \multicol} -- This package is specifically forbidden
% \nameref} -- This package is specifically forbidden
% \usepackage{savetrees} -- This package is specifically forbidden
% \usepackage{setspace} -- This package is specifically forbidden
% \usepackage{stfloats} -- This package is specifically forbidden
% \usepackage{tabu} -- This package is specifically forbidden
% \usepackage{titlesec} -- This package is specifically forbidden
% \usepackage{tocbibind} -- This package is specifically forbidden
% \usepackage{ulem} -- This package is specifically forbidden
% \usepackage{wrapfig} -- This package is specifically forbidden
% DISALLOWED COMMANDS
% \nocopyright -- Your paper will not be published if you use this command
% \addtolength -- This command may not be used
% \balance -- This command may not be used
% \baselinestretch -- Your paper will not be published if you use this command
% \clearpage -- No page breaks of any kind may be used for the final version of your paper
% \columnsep -- This command may not be used
% \newpage -- No page breaks of any kind may be used for the final version of your paper
% \pagebreak -- No page breaks of any kind may be used for the final version of your paperr
% \pagestyle -- This command may not be used
% \tiny -- This is not an acceptable font size.
% \vspace{- -- No negative value may be used in proximity of a caption, figure, table, section, subsection, subsubsection, or reference
% \vskip{- -- No negative value may be used to alter spacing above or below a caption, figure, table, section, subsection, subsubsection, or reference

\setcounter{secnumdepth}{2} %May be changed to 1 or 2 if section numbers are desired.

% The file aaai25.sty is the style file for AAAI Press
% proceedings, working notes, and technical reports.
%

% Title

% Your title must be in mixed case, not sentence case.
% That means all verbs (including short verbs like be, is, using,and go),
% nouns, adverbs, adjectives should be capitalized, including both words in hyphenated terms, while
% articles, conjunctions, and prepositions are lower case unless they
% directly follow a colon or long dash
%\title{PieClam: Prior Inclusive-Exclusive Community Affiliation Model \\ With Application to Graph Anomaly Detection}
\title{PieClam: Supplementary Material (After the main paper)}
\author{
    %Authors
    % All authors must be in the same font size and format.
    Anonymous Authors
}

%\author{
%    Written by AAAI Press Staff\textsuperscript{\rm 1}\thanks{With help from the AAAI Publications Committee.}\\
%    AAAI Style Contributions by Pater Patel Schneider,
%    Sunil Issar,\\
%    J. Scott Penberthy,
%    George Ferguson,
%    Hans Guesgen,
%    Francisco Cruz\equalcontrib,
%    Marc Pujol-Gonzalez\equalcontrib
%}

\affiliations{
    %Afiliations
    \textsuperscript{\rm 1}Association for the Advancement of Artificial Intelligence\\
    % If you have multiple authors and multiple affiliations
    % use superscripts in text and roman font to identify them.
    % For example,

    % Sunil Issar\textsuperscript{\rm 2}, 
    % J. Scott Penberthy\textsuperscript{\rm 3}, 
    % George Ferguson\textsuperscript{\rm 4},
    % Hans Guesgen\textsuperscript{\rm 5}
    % Note that the comma should be placed after the superscript

    1101 Pennsylvania Ave, NW Suite 300\\
    Washington, DC 20004 USA\\
    % email address must be in roman text type, not monospace or sans serif
    proceedings-questions@aaai.org
%
% See more examples next
}

%Example, Single Author, ->> remove \iffalse,\fi and place them surrounding AAAI title to use it
\iffalse
\title{My Publication Title --- Single Author}
\author {
    Author Name
}
\affiliations{
    Affiliation\\
    Affiliation Line 2\\
    name@example.com
}
\fi

\iffalse
%Example, Multiple Authors, ->> remove \iffalse,\fi and place them surrounding AAAI title to use it
\title{My Publication Title --- Multiple Authors}
\author {
    % Authors
    First Author Name\textsuperscript{\rm 1,\rm 2},
    Second Author Name\textsuperscript{\rm 2},
    Third Author Name\textsuperscript{\rm 1}
}
\affiliations {
    % Affiliations
    \textsuperscript{\rm 1}Affiliation 1\\
    \textsuperscript{\rm 2}Affiliation 2\\
    firstAuthor@affiliation1.com, secondAuthor@affilation2.com, thirdAuthor@affiliation1.com
}
\fi


% REMOVE THIS: bibentry
% This is only needed to show inline citations in the guidelines document. You should not need it and can safely delete it.
\usepackage{bibentry}
% END REMOVE bibentry

\begin{document}

\maketitle

\begin{abstract}


%\textcolor{red}{Technical contribution to focus on: - Minkowski inner produce decoder - a simple extension of inner produce decoder  - that is able to represent bipartite parts of graphs, and leads to universality of the  - graph model. - The addition of prior to Clam, so it becomes a generative model. - Efficiency of the method for sparse graph - order of the number of edges. - Focus on the usefulness for anomaly detection}

% i copied so i can change, the original is down there
% We propose PieClam (Prior Inclusive Exclusive Cluster Affiliation Model): a probabilistic graph model for representing any graph as overlapping  generalized communities.  Our method can be interpreted as a graph autoencoder: nodes are embedded into a code space by an algorithm that maximizes the log-likelihood of the decoded graph, given the input graph. PieClam is a community affiliation model that extends well-known methods like BigClam in two main manners. First, instead of the decoder being defined via pairwise interactions between the nodes in the code space, we also incorporate a learned prior on the distribution of nodes in the code space, turning our method into a graph generative model. Secondly, we generalize the notion of communities by allowing not only sets of nodes with strong connectivity, which we call inclusive communities, but also sets of nodes with strong disconnection, which we call exclusive communities.  To model both types of communities, we propose a new type of decoder based the Lorentz inner product, which we prove to be much more expressive than standard decoders based on standard inner products or norm distances, making PieClam a universal autoencoder. Our method is shown to obtain competitive performance in graph anomaly detection benchmarks.



We propose PieClam (Prior Inclusive Exclusive Cluster Affiliation Model): a probabilistic graph model for representing any graph as overlapping  generalized communities.  Our method can be interpreted as a graph autoencoder: nodes are embedded into a code space by an algorithm that maximizes the log-likelihood of the decoded graph, given the input graph. PieClam is a community affiliation model that extends well-known methods like BigClam in two main manners. First, instead of the decoder being defined via pairwise interactions between the nodes in the code space, we also incorporate a learned prior on the distribution of nodes in the code space, turning our method into a graph generative model. Secondly, we generalize the notion of communities by allowing not only sets of nodes with strong connectivity, which we call inclusive communities, but also sets of nodes with strong disconnection, which we call exclusive communities.  To model both types of communities, we propose a new type of decoder based the Lorentz inner product, which we prove to be much more expressive than standard decoders based on standard inner products or norm distances. By introducing a new graph similarity measure, that we call the log cut distance, we show that PieClam is a universal autoencoder, able to uniformly approximately reconstruct any graph. Our method is shown to obtain competitive performance in graph anomaly detection benchmarks.


%\textcolor{cyan}{
%Inner product based decoders are ubiquitous in all kinds of graph representation models and have been shown to approximate and set performance records on a wide range of real world and synthetic graph related problems. Despite this unquestionable success, inner product based decoding is provably and evidently limited in it's expressivity in that it is unable to recreate heterophilic graph structures.
%We propose PieClam (Prior Inclusive Exclusive Cluster Affiliation Model): a probabilistic graph model for representing any graph as overlapping generalized communities which can be interpreted as a graph autoencoder. PieClam is a community affiliation model that extends well-known methods like BigClam, embedding nodes into a code space by maximizing the log likelihood of the decoded graph.
%We generalize the notion of communities by allowing not only sets of nodes with strong connectivity, which we call inclusive communities, but also sets of nodes with strong disconnection, which we call exclusive communities. To model both types of communities, we propose a new type of decoder based on the Lorentz inner product, which we prove to be universally expressive and therefore superior to standard decoders based on inner products or norm distances. Furthermore, instead of the decoder being defined only via pairwise interactions between the nodes in the code space, we also incorporate a learned prior over the encoded feature distribution, turning our method into a graph generative model with the ability to also incorporate external node attributes if they exist. 
%Finally, our method is shown to obtain competitive performance in graph anomaly detection benchmarks.} 


%Large graphs are of much interest in many fields of science and engineering, as the non trivial structure combined with massive size required extra processing to extract useful information.  Examples can be found in any field of network science from social networks to drug manufacturing.

%Generative graph models are a solution that interprets graphs as being generated from some underlying phenomenon distribution. Specifically, feature embedding models encode a graph into a feature vector space that can also be thought of as a continuous version of overlapping group affiliations where similar node features imply similar node connectivity and thus the graph topology is replaced with the metric topology of the feature space. 

%The BigClam algorithm is a community affiliation model which maximizes the probability of the of the graph conditioned on the community embedding, and has proved very efficient in finding community affiliations for graphs containing billions of connections in $\mathcal{O}(E)$ time complexity.
%Building up on the proven foundation of BigClam, we created a generative graph model which we prove to be a universal approximator. Though effective for data mining, BigClam is lacking in expressivity for as we prove in the following, it cannot approximate even a simple bipartite graph in a non trivial way. Also, it is weak as a probabilistic model, as it relies on maximizing a conditional probability and therefore if we think of the joint probability of features and edges, the marginal probability of the features is uniform. 

%To achieve universality, we define the code space to be a generalized for of the Minkowski space known as the Pseudo-Metric, and prove that in this space, the BigClam can achieve universal approximation using an embedding dimension of  only $\mathcal{O}(\frac{1}{\epsilon^2})$ for a precision level of $\epsilon$ in cut distance.
%\textcolor{cyan}{in the course asymptotic methods we were taught that $\log({\epsilon})$ can be treated as $\mathcal{O}(1)$ in a case like this so instead of $\frac{log(\epsilon)}{\epsilon^2}$ i wrote this}

%Furthermore, we endow the probabilistic model with a prior distribution on the feature space which improves the model's generative expressivity and provides regularization in the fitting process. 

%Lastly, we demonstrate the power and versatility of our model with experiments. First, we generate synthetic graphs from given esoteric prior distributions and off diagonal SBM models, and relearn these distributions using only the graph structure. Second, we use our model on an unsupervised anomaly detection task with several datasets and show that we are competitive with SOTA on both attributed and non attributed graphs.

\end{abstract}

% Uncomment the following to link to your code, datasets, an extended version or similar.
%
% \begin{links}
%     \link{Code}{https://aaai.org/example/code}
%     \link{Datasets}{https://aaai.org/example/datasets}
%     \link{Extended version}{https://aaai.org/example/extended-version}
% \end{links}


\section{Introduction}

In recent years, considerable research has concentrated on graph representation learning, aiming to develop vector representations for graph entities, including nodes, edges, and subgraphs \cite{hamilton2017representation,GRL_survet2020}. In graph autoencoders, e.g.,  \cite{kipf2016variational,grover2019graphite,JMLR:v21:19-671,mehta2019stochastic}, the vertices of a graph are embedded in a code space, where edges are inferred from the locations of the vertices in this space. Encoding graphs into a standard space has a number of advantages. While different graphs can have different sizes and topology, the code space is fixed with a fixed dimension. This helps when  learning downstream tasks, where the representation of the graph in the code space can be processed. For example, in link prediction, one infers unknown edges by defining \cite{kipf2016variational} or learning \cite{kumar2020linkreview} a function that takes pairs of nodes in the code space and predicts if there is an edge between them. In anomaly detection, one defines \cite{ding2019dominant,fan2020anomalydae} or learns \cite{chen2020gaan,ding2021aegis},  a function that takes the representation of a node and its neighborhood and predicts if this node is normal or an anomaly. In graph and node classification, the graph is represented in the code space, and one learns a model that predicts from this representation the classes of the nodes or of the graph \cite{GCN,MPNN}. %\ron{references. Basically MPNNs,GNNs. Cite surveys and books.}. 

\paragraph{Our Contribution.} In this paper, we derive a new graph autoencoder from a statistical model of graphs. In our model, similarly to SBM \cite{nowicki2001estimation,lee2019review} or community-based statistical models  \cite{airoldi2008mixed,yang2013BigClam}, graphs are generated from a combination of intersecting communities/cliques. Namely, each node belongs to a different subset of a predefined set of communities, and the affiliations of the nodes to the different communities determine the probabilities of the edges of the graph. Here, the estimation of the community affiliations is seen as an encoder of graphs to a  \emph{community affiliation space}, and the computation of the corresponding edge probabilities is seen as a decoder. As opposed to past works, we consider two types of \emph{generalized communities}. First, standard \emph{inclusive communities}, where any two nodes in the same community are likely to be connected. Second, we propose \emph{exclusive communities}, where belonging to the same community reduces the probability of nodes being connected. For illustration, consider a social network of employers/recruiters and employees/job-seekers. Such a network has roughly bipartite components, where job-seekers do not tend to connect to other seekers, and employers do not connect to other employers, but employers and seekers of the same subsector tend to connect. Hence, inclusive communities for such a graph can correspond to job titles or subsectors, and exclusive communities can correspond to sets of job-seekers and sets of employees from the same sector. %\ron{Can we find any reference related to this? Read a survey on social networks data analysis}

To formalize the above ideas, we propose the \emph{Prior Inclusive Exclusive Cluster Affiliation Model (PieClam)}. This model represents graphs as overlapping inclusive and exclusive communities. Moreover, instead of embedding the nodes of a given graph in the community code space, our model also learns a prior on the code space, so new graphs can be generated from the learned model. This makes PieClam a graph generative model, like, e.g.,  \cite{kipf2016variational,grover2019graphite,JMLR:v21:19-671,mehta2019stochastic,sun2019vgraph}. 

To model both types of communities, we propose a new type of decoder based on the \emph{Lorentz inner product}, in which case the code space is typically called a \emph{pseudo Euclidean space} \cite{greub1963pseudo}. The addition of exclusive communities in our model is not merely aimed at improving it heuristically for special graphs like social networks. Rather, we prove in Theorems \ref{thm:IEUniversality} and \ref{thm:IEUniversality2} that using exclusive communities (via the Lorentz inner product) makes our model universal, namely, able to approximate with a fixed budget of parameters any graph. This is in contrast to standard decoders based on only inclusive communities, which we show are unable to represent many graphs. %In contrast, we prove that ou PieClam, based on the Lorentz inner product, can approximately represent any graph, of any size and topology, with a fixed limited budget of parameters.  

To formalize this universality property, we propose a new similarity measure between graphs with edge probabilities, that we call the \emph{log cut distance}. We formalize the universality of our model as follows: one can choose the dimension of the code space (the number of communities)  a priori, and guarantee that any graph of any size and topology can be approximated up to small log cut distance by decoding points from this fixed space. We show that other related decoders do not satisfy this universality property. 

In Section \ref{Experiments} we support our construction with experiments.
We first conduct some toy experiments that illustrate the merits of our method.
We then use PieClam to perform graph anomaly detection. Here, PieClam learns a probabilistic model corresponding to a given graph, and this model can be directly used for inspecting the probabilities of different nodes in this graph. Nodes with low probabilities are deemed to be anomalies. Our models achieve competitive performance with respect to state of the art anomaly detection methods.


%
%Our method is shown to obtain competitive performance in graph anomaly detection benchmarks.


Appendix \ref{Extended Related Work} has 
%offers
an extended discussion on related work.

\section{Community Affiliation Models With Prior}

Our model PieClam is best understood as an extension of the BigClam model \cite{Overlapping2012}. Hence, after introducing basic notations, we start with a detailed exposition of BigClam, before introducing our novel constructions.


\subsection{Notations}
\label{def:notations}
We denote by $\mathbb{R}$ the real line, and by $\mathbb{R}_+$ the non-negative real numbers. We denote ``and'' by $\wedge$.
We denote matrices as boldface capital letter $\mathbf{B}=\{b_{n,m}\}_{n,m}$, vectors as boldface lowercase letters $\mathbf{b}=\{b_n\}_n$,  and scalars as regular lowercase letters $b$. Vectors $\mathbf{b}\in\mathbb{R}^N$ are always column vectors, and their transpose $\mathbf{b}^{\top }$ row vectors. The rows of a matrix $\mathbf{B}\in\mathbb{R}^{N\times C}$ are denoted by the same letter in lowercase $\mathbf{b}_n^{\top }$, where $\mathbf{b}_n\in\mathbb{R}^C$, where we write in short $\mathbf{b}_n^{\top }\in\mathbb{R}^C$. A diagonal matrix $\mathbf{D}\in\mathbb{R}^{N\times N}$ with diagonal entries $\mathbf{d}$ is denoted by ${\rm diag}(\mathbf{d})={\rm diag}(d_1,\ldots, d_N)$.  The $\ell^2$ norm of a vector $\mathbf{b} \in \mathbb{R}^N$ is defined to be $\|\mathbf{b}\| = (\sum_{n=1}^Nb_n^2)^{1/2}$. 

A graph is denoted by $G=([N],E,\mathbf{A})$, where $[N]=\{1,\ldots, N\}$ is the set of $N$ nodes,  $E\subseteq [N]\times[N]$ is the set of edges, and $\mathbf{A}\in \{0,1\}^{N\times N}$ is the adjacency matrix. For weighted graphs, $\mathbf{A}\in[0,1]^{N\times N}$, where $a_{n,m}$ is the edge weight of $(n,m)$.
In this work we focus on undirected graphs, for which $(m,n)\in E \Leftrightarrow (n,m)\in E$, and $\mathbf{A}=\mathbf{A}^{\top }$.  
%
Any pair $(n, m)\in [N]\times[N]$ is called a \emph{dyad}.
The neighborhood of a node $n\in[N]$ is  $\mathcal{N}(n)  = \{m \in [N] \mid (m,n) \in E\}$.
 A graph-signal is a graph with node features  $G=([N],E,\mathbf{A},\mathbf{X})$ where $\mathbf{X}=\{\mathbf{x}_n^{\top }\}_{n=1}^N\in \mathbb{R}^{N\times D}$, $\mathbf{x}_n\in\mathbb{R}^D$, and $D$ is called the feature dimension. 
A random graph is a graph-valued random variable. Given a random graph with nodes $[N]$, 
 we denote the event in which  $(n,m)\in E$ by  $n\sim m$, and the event $(n,m)\notin E$ by $\neg (n\sim m)$.  
 %
%
A graph is bipartite if its vertex set $[N]$ can be partitioned into two disjoint sets $\mathcal{U}$ and $\mathcal{V}$, with  $\mathcal{U}\cup\mathcal{V}=[N]$, such that every edge has one endpoint in $\mathcal{U}$ and the other in $\mathcal{V}$. 
%In a complete bipartite graph, all of the nodes in $U$ are connected to all of the nodes in $W$ and there are no connections within the sets.

%In other words, a bipartite graph can be divided into two sets of nodes where each set connects to the other set but not within itself. 


\subsection{BigClam}
\label{sec:BigClam}
%\textcolor{red}{Explain the method. Introduce terms that we will use: affiliation space, dyad probability, etc (or maybe some of these belong in the previous chapter)}

%Since our model extends BigClam \cite{Overlapping2012}, let us start with introducing this model in detail. 
 The code space in BigClam is $\mathbb{R}_+^C$, where each axis is interpreted as a community.  Each entry $f^c$ of a point $\mathbf{f}=(f^1,\ldots,f^C)\in\mathbb{R}_+^C$ is interpreted as how much the point belongs to community $c$, where $0$ means ``not included in $c$'' and $1$ ``included.'' In this paper we  call $\mathbb{R}_+^C$ \emph{the affiliation space (AS)}, and call any point in the affiliation space an \emph{affiliation feature (AF)}.

BigClam is a model where a simple random graph is decoded from a sequence of AFs $\mathbf{F}= \{\mathbf{f}_n=(f_n^1,\ldots,f_n^C)\}_{n=1}^N$ in the AS. For each pair of nodes $n,m\in[N]$, the probability of the event $\neg (n\sim m)$ (no edge between $n$ and $m$), given their amount of membership in the same community $c$, is defined to be
\[
%\label{eqn:poisson}
P(\neg (n\sim m)|f^c_n, f^c_m) = e^{-f^c_n f^c_m}.
\]
BigClam makes two assumptions of independence. First, the membership in one community does not affect the membership in another community. Hence,
\begin{equation}
\label{eq:ClamEdge}
    P(\neg (n\sim m))|\mathbf{F})=P(\neg (n\sim m))|\mathbf{f}_n, \mathbf{f}_m)=e^{-\mathbf{f}_n^{\top }\mathbf{f}_m},
\end{equation}
\[P(n\sim m)|\mathbf{F})=P(n\sim m|\mathbf{f}_n, \mathbf{f}_m)=1-e^{-\mathbf{f}_n^{\top }\mathbf{f}_m}.\]
Due to this formula, BigClam is a so called \emph{Bernoulli-Poisson model} 
(see Appendix \ref{Stochastic block models}  for more details).
The second assumption is that the events of having and edge between different pairs of nodes are independent.  As a result, the probability of the entire graph $([N],E)$ conditioned on the AFs $\mathbf{F}$ is
%\textcolor{cyan}{this formula is already given in the previous section,}
\begingroup
\begin{equation}
\footnotesize
\begin{aligned}
   &  P(E|\mathbf{F}) =\\
   &  \prod_{n \in [N]}\Bigg(
        \sqrt{\prod_{m \in \mathcal{N}(n)}P( n\sim m|\mathbf{F})\prod_{m \notin \mathcal{N}(n)} P(\neg (n\sim m)|\mathbf{F})}\Bigg)
\end{aligned}
\label{eqn:clam_prob}
\end{equation}
\endgroup
Here, the square root is taken because the product considers each edge twice. 

BigClam consist of a decoder, decoding from the AFs $\mathbf{F}$ the random graph $G(\mathbf{F})$ with node set $[N]$ and independent Bernoulli distributed edges with probabilities $\mathbf{P}=\{P(n\sim m|\mathbf{f}_n, \mathbf{f}_m)\}_{n,m=1}^N$.

From the encoding side, given a simple graph $([N],E)$, BigClam encodes the graph into the affiliation space by maximizing the log likelihood with respect to the AFs $\mathbf{F}$
\begingroup
\footnotesize
\begin{align}
\label{eqn:naiveloss}
     l(\mathbf{F}) &= \frac{1}{2}\sumgraph \Big(\sumneighbors{}{\log(1-e^{-\mathbf{f}_n^{\top } \mathbf{f}_m})} - 
 \sumstrangers{\mathbf{f}_n^{\top } \mathbf{f}_m}\Big).
 \end{align}
 \endgroup
BigClam is optimized by gradient descent, with update at iteration $i$
    \begin{align}
    \label{eqn:vanilla_optimization}
    \mathbf{F}^{(i+1)} = \mathbf{F}^{(i)} + \delta\nabla_{\mathbf{F}}l({\mathbf{F}}^{(i)}),
    \end{align}
for some learning rate $\delta>0$.
In order to implement the above iteration with $O(|E|)$ operations at each step, instead of $O(N^2)$, the loss can be rearranged as
\begingroup
\begin{equation}
\footnotesize
\begin{aligned}
\label{eqn:comp_loss}
   &  2l(\mathbf{F})=\\
   & \sumgraphu \Big(\sumneighbors{\log(e^{\mathbf{f}_n^{\top } \mathbf{f}_n} - 1)}  - \mathbf{f}_n^\top {\sumgraph{\mathbf{f}_m}} + \|\mathbf{f}_n\|^2\Big).
\end{aligned}
\end{equation}
\endgroup
The gradient of the loss is now
\begin{align}
\label{eqn:BCupdate}
    % \nabla_{F_u}l = 
    % \sumneighbors{F_v \frac{e^{\inr[F_u]{F_v}}}{e^{\inr[F_u]{F_v}-1}}} - \sumgraph{F_v} + F_u \\
    \nabla_{\mathbf{f}_n}l = 
    \sumneighbors{\mathbf{f}_m \big(1-e^{-\mathbf{f}_n^\top  \mathbf{f}_m}
    \big)^{-1}} - \sumgraph{\mathbf{f}_m} + \mathbf{f}_n .
\end{align}
Since the global term needs to be calculated only once, the nu
mber of operations is  $O(|E|)$ instead of $O(N^2)$.  

We observe that the optimization process is a message passing scheme.
Looking at the dynamics of the optimization process, we see that every node is pushed in the direction of a weighted average of all of its neighbors, and pushed in the opposite direction by an average of all of the nodes. The sum of both forces tends to drive communities towards the axes in the optimization dynamics.


\subsection{Inclusive-Exclusive Cluster Affiliation Model}


%\textcolor{red}{Motivate the construction of bipartites. The need to find a decoder that satisfy what is needed of bipartites: two sets well connected to each other but not to themselves.} \textcolor{cyan}{i defined it in "notations"}
%\textcolor{cyan}{not sure how deep i should go about this... tell me if showing for a full bipartite is not good enough}
%To see the limited expressivity of a decoder that is monotonic with the inner product between nodes (like BigClam), let us consider what an embedding of a bipartite (\ref{def:notations}\textcolor{cyan}{problem referencing to sections as there is no numbering in this format}) graph should look like for it to be decoded in the best way. The nodes of $U$ will have to be perpendicular to one another, each being aligned with all of the nodes in $W$, which must also be perpendicular to one another.\textcolor{cyan}{don't know how to prove this yet. it's easy to show if the nodes are in a cluster}    
%\textcolor{cyan}{...}
%\textcolor{red}{Define the IeClam model.}



%\[\]

The BigClam decoder has a limitation due to a  ``triangle inequality type'' behavior. Namely, suppose that we would like to construct two features $\mathbf{f}_1$ and $\mathbf{f_2}$ with strong connectivity to a third feature $\mathbf{f}_3$, and  we are limited by a fixed affiliation feature dimension $C$. A naive approach to achieve this would be to put $\mathbf{f}_1$ and $\mathbf{f_2}$ close to $\mathbf{f}_3$ so they have large inner products $\mathbf{f}_1^{\top }\mathbf{f}_3$ and $\mathbf{f}_2^{\top }\mathbf{f}_3$ . This would mean that $\mathbf{f}_1^{\top }\mathbf{f}_2$ would also be large, so $\mathbf{f}_1$ would be strongly connected to  $\mathbf{f}_2$ under the BigClam model. However, some graphs, like bipartite graphs, do not exhibit this triangle inequality type behavior. For a rigorous treatment, see Section \ref{BigClam is Not Universal} and Appendix \ref{Proof That BigClam Is Not Universal}.  Next, we build the IeClam decoder, that allows decoding any graph, that may have bipartite components without being limited by a  triangle inequality-type behavior. 

Our \emph{Inclusive Exclusive Cluster Affiliation Model (IeClam)} can be extended from BigClam by replacing the inner product in the non-edge probability (\ref{eq:ClamEdge}) by the more expressive \emph{Lorentz inner product}, which does not enforce a triangle inequality-type behavior. 
For that, we extend the affiliation space (AS) to be $\mathbb{R}^{2C}$, with two types of communities. The first $C$ axes are called the \emph{inclusive communities}, and their corresponding features are called \emph{inclusive affiliation features (IAF)}, denoted by $\mathbf{t}\in\mathbb{R}^C$. The last $C$ axes are called the \emph{exclusive communities}, with \emph{exclusive affiliation features (EAF)}, denoted by $\mathbf{s} \in \mathbb{R}^C$.\footnote{More generally, one can define a different number of inclusive and exclusive communities.} We define the concatenated affiliation feature by $\mathbf{f} = (\mathbf{t},\mathbf{s})\in\mathbb{R}^{2C}$. Given a sequence of affiliation features  $\mathbf{F}=\{\mathbf{f}_n\}_{n=1}^N\in\mathbb{R}^{N\times 2C}$, IeClam defines the probability of a single edge by
\begin{equation}
\begin{split}
    P(n\sim m| \mathbf{f}_n,\mathbf{f}_m) & = 1-\exp\big({-\mathbf{t}_n^{\top } \mathbf{t}_m +  \mathbf{s}_n^{\top } \mathbf{s}_m}\big) \\  
   &  = 1-\exp\big({-\mathbf{f}_n^{\top } \mathbf{L} \mathbf{f}_m}\big),
    \label{eq:LorFeature}
    \end{split}
\end{equation}
where $\mathbf{L} = {\rm diag}(1,...1,-1,...-1)$. The bilinear form $(\mathbf{u},\mathbf{v})\mapsto \mathbf{u}^\top  \mathbf{L} \mathbf{v}$ is called the \emph{Lorentz inner product}, and has its roots in special relativity. See Appendix \ref{The Lorentz Inner Product} for more details on the Lorentz inner product. 

Note that $\mathbf{L}$ is not positive-definite, so it does not actually define an inner product. Moreover, $\mathbf{f}_n^\top  \mathbf{L} \mathbf{f}_m$ can be negative even if $\mathbf{f}_n,\mathbf{f}_m\in \mathbb{R}_+^C$. To guarantee that (\ref{eq:LorFeature}) defines a proper probability between $0$ and $1$, we limit the affiliation space as follows.  
%
\begin{definition}
  A \emph{cone of non-negativity} is a  subset $\mathcal{C}$ of $\mathbb{R}^{2C}$ such that for every $\mathbf{f},\mathbf{g}\in \mathcal{C}$ we have 
$\mathbf{f}^\top  \mathbf{L} \mathbf{g} \geq 0$.  
\end{definition}
If we limit the affiliation space to be a cone of non-negativity, then IeClam gives well defined probabilities in $[0,1]$.
In our experiments, we restrict ourselves to the following simple construction of a cone of non-negativity, noting that it is not the only possible construction. 
%
%For each $c\in[C]$, we take the direct product of axis $c$ with axis $C+c$, and denote it by $\mathcal{C}^c=\sim\mathbb{R}^2$. Given $\mathbf{f}\in\mathbf{C}^{2C}$, we call $(t^c,s^c)$ the entry of $\mathbf{f}$ in $\mathcal{C}^c$. 
\begin{definition}
 The \emph{pairwise cone} $\mathcal{T}$ is defined to be the set of affiliation features $\mathbf{f}=(\mathbf{t},\mathbf{s})\in\mathbb{R}^{2C}$ such that for every $c\in[C]$ we have $-t^c \leq s^c\leq t^c$.    
\end{definition}
It is easy to see that $\mathcal{T}$ is a cone of non-negativity. Indeed, for any $\mathbf{f}_1=(\mathbf{t}_1,\mathbf{s}_1), \mathbf{f}_2=(\mathbf{t}_2,\mathbf{s}_2) \in \mathcal{T}$,  we have  
 \[\forall c\in[C]: \quad  t_1^c t_2^c - s_1^c s_2^c \geq 0,\] so
\[\mathbf{f}_1^\top  \mathbf{L} \mathbf{f}_2 = \sum_{c=1}^C t_1^c t_2^c - s_1^c s_2^c \geq 0.\] 
%
%
As opposed to the BigClam decoder, the IeClam model can approximate a bipartite graph with a small number of communities. Namely, a bipartite graph with two (disjoint) sides $\mathcal{U},\mathcal{V}\subset[N]$ and constant probability for edges between $\mathcal{U}$ and $\mathcal{V}$ can be represented using $C=1$. Here,  all of the nodes in $\mathcal{U}$ are encoded to $(a,a)$,  and all of the nodes of $\mathcal{V}$ are encoded to $(a,-a)$, for some $a\geq 0$. This gives zero probability for edges within each part, and probability $1- e^{-2a^2}$ for edges between the parts. %1- e^{-2a^2}%\ron{now write the construction that approximates a full bipartite graph (where the probability of having a edge between any two nodes from opposite sides is a constant).}

\begin{remark}
    The above construction gives an interpretation for each pair of axes $(t^c,s^c)$ in $\mathcal{T}$ as a \emph{generalized community} which can model anything between a clique and a bipartite component.
\end{remark}

%for every $u \in U$ and  $w \in W$, a complete bipartite graph can be encoded using $\mathbf{f}_u = (a,a)$  and $\mathbf{f}_w = (a,-a)$, for some scalar $a\in \mathbb{R}/{0}$. In this construction, $P(u\sim w) = 1- e^{-2a^2}$ which goes to $1$ when $a \rightarrow \infty$, and $P(\neg(u_1\sim u_2)) = P(\neg(w_1\sim w_2)) = 0$ for all $u_1,u_2 \in U$ and $w_1, w_2 \in W$ . 

Given AFs $\mathbf{F}$ in a cone, IeClam is seen as a decoder, by decoding $\mathbf{F}$ into the random graph $G(\mathbf{F})$ with node set $[N]$ and independent Bernoulli distributed edges with probabilities $\mathbf{P}=\{P(n\sim m|\mathbf{f}_n, \mathbf{f}_m)\}_{n,m=1}^N$.

For affiliation features $\mathbf{F}\in \mathcal{T}^N$, the probability of the graph $([N],E)$ given $\mathbf{F}$ of IeClam is
\begingroup
\footnotesize
\begin{equation}
\begin{aligned}
   &  P(E|\mathbf{F}) \\
   & = \prod_{n \in [N]}\Bigg(
        \sqrt{\prod_{m \in \mathcal{N}(n)}(1-e^{-\mathbf{f}_n^\top  \mathbf{L} \mathbf{f}_m})\prod_{m \notin \mathcal{N}(n)} e^{-\mathbf{f}_n^\top  \mathbf{L} \mathbf{f}_m}}\Bigg).
\end{aligned}
\label{eqn:IeClam_prob}
\end{equation}
\endgroup
Like BigClam, IeClam is optimized by maximizing the log likelihood with gradient descent
\begin{equation}
\begin{aligned}
   &  2 l(\mathbf{F}) =
    \\ &\sumgraphu \Big(\sumneighbors{}\log(1 - e^{-{\mathbf{f}_n}^\top  \mathbf{L} \mathbf{f}_m})  -   \sumstrangers{}  \mathbf{f}_n ^\top  \mathbf{L} \mathbf{f}_m\Big)
    \end{aligned}
\label{eqn:IeClam_loss}
\end{equation}
This loss can be efficiently implemented on sparse graph by the formulation
\begingroup
\begin{equation}
\footnotesize
\begin{aligned}
\label{eqn:loss_reparam_ie}
   &  2l(\mathbf{F})=\\
   & \sumgraphu \Big(\sumneighbors{\log(e^{\flf} - 1)}  - \mathbf{f}_n^\top  \mathbf{L}{\sumgraph{\mathbf{f}_m}} + \mathbf{f}_n^\top  \mathbf{L}\mathbf{f}_n\Big).
\end{aligned}
\end{equation}
\endgroup
The gradient of the loss for node $n$ is
\begin{align}
\label{eqn:BCupdate_ie}
    \nabla_{\mathbf{f}_n}l = 
   \mathbf{L}\bigg( \sumneighbors{\mathbf{f}_m \big(1-e^{-\flf}
    \big)^{-1}} - \sumgraph{\mathbf{f}_m} + \mathbf{f}_n \bigg).
\end{align}
Notice that all of the calculations are the same as BigClam, up to replacing the dot product by the Lorenz inner product.
%We can see that any reconstruction achieved using BigClam with $C$ communities can also be achieved using IeClam with $2C$ IE communities. In practice, when graphs have strong inclusive connections, IeClam can preform worse than BigClam due to overfitting.

% \textcolor{cyan}{add an explanation for allowing the features in the time cone so that the product is non negative. nodes on opposide sides of the cone have a high probability of connecting, as well as features with high t and low s}
% We refer to $\mathbf{t}$ communities as "inclusive", since if the affiliation of two nodes to a $\mathbf{t}$ type community is of the same sign it increases the probability of connection. Similarly, if the affiliations of two nodes to an $\mathbf{s}$ community is of an opposite sign it increases the probability of connection.

% The inclusive-exclusice (IE) modification replaces the Hilbert code space with a version of the Minkowski space where there are the same amount of "time" coordinates as there are "space" coordinates \textcolor{cyan}{citation? i think from books you need to give the page}. 

% For the probability model to hold, we need the argument of the exponent to be positive. This was easy when there were only inclusive communities and indeed, BigClam was restricted to positive features To avoid divergence of the model. 

% In order to achieve positivity in the exponent argument we restrict the domain to pairs of "time-cones".
% \begin{definition}
%     For a node $n\in[N]$, a timecone is defined to be the pair $\mathbf{\tau}_n[i] = (t_n[i], s_n[i]) = (f_n[i], f_n[i+k])$  
% \end{definition}
% For all time cones we restrict $-t_n[i] \leq s_n[i] \leq t_n[i]$.
% The definition of timecones assures that the argument of the exponent is always non positive. Despite being a subset of ${\mathbf{f}_n}^\top  B \mathbf{f} \geq 0$, restriction of features to the timecones eases the computation time and breaks the symmetry of the initial IE definition, creating generalized communities which can be interpreted as each having an inclusive and exclusive part.

% In some cases, BigClam can still outperform IeClam as graphs can have a strong community structure that can be overfitted with IeClam.

\subsection{Community Affiliation Models With Prior}

%\textcolor{red}{Talk about the limit of BigClam and IeClam as generative models- no way to approximate different graphon this way.}
BigClam and IeClam are not  generative graph models. Indeed, these methods only fit a conditional probability of the graph, conditioned on the AF values, but the methods do not learn the probability of the AFs over the affiliation space. Hence, the total probability of $E\wedge\mathbf{F}$ is not defined. %If one was trying to generate a graph using only the features learned using this model, they could only either use the learned node features, or interpolate them in some way that is not rigorous. In the other direction, there is no option to obtain the probability of a test graph given a training graph, only by running the method again for the test graph.
%
%Therefore, neither of the models can be considered a generative graph model, as they do not provide a distribution over the AS.

%\textcolor{red}{Explain our solution - define a prior distribution in the affiliation space.}
To extend IeClam (and similarly BigClam) into probabilistic generative models, we define a prior probability distribution over the affiliation cone space $\mathcal{C}\subset \mathbb{R}^{2C}$, with probability density function $p:\mathbb{R}^{2C}\rightarrow [0,\infty)$ supported on $\mathcal{C}$. Using Bayes law, we now obtain the joint probability $P(E \wedge \mathbf{F})$ of the edges and community affiliation features  via the probability density function
\begin{align*}
    p(E ,\mathbf{F}) = P(E|\mathbf{F}) p(\mathbf{F}).
\end{align*}
We assume that the prior probabilities of all nodes are independent, namely,  
\[p(\mathbf{F}) = \underset{n\in [N]}{\prod} p(\mathbf{f}_n).\]
 Hence, the probability densities that a dyad $(n,m)$ is an edge or non-edge are
\[p(n\sim m, \mathbf{f_n}, \mathbf{f_m}) = p(\mathbf{f}_n)p(\mathbf{f}_m)(1 - e^{-{\mathbf{f}_n}^\top  \mathbf{L} {\mathbf{f}_m}}),\]
\[p(\neg(n\sim m), \mathbf{f_n}, \mathbf{f_m}) = p(\mathbf{f}_n)p(\mathbf{f}_m)e^{-{\mathbf{f}_n}^\top  \mathbf{L} {\mathbf{f}_m}}.\]
As before, we assume that the probabilities of different edges are independent, which gives
\begin{equation}
\scriptsize
\begin{aligned}
\label{eqn:likelihood_prior}
    &p(E, \mathbf{F}) =  \\ &\prod_{n \in [N]} p(\mathbf{f}_n)
        \sqrt{\prod_{m \in \mathcal{N}(n)}P(n\sim m|\mathbf{F}_m)\prod_{m \notin \mathcal{N}(n)}{P(\neg(n\sim m)|\mathbf{f}_n, \mathbf{f}_m)}
        }. 
 \end{aligned}
 \end{equation}
Now, the log likelihood loss is
\begingroup
\begin{equation}
\footnotesize
\begin{aligned}
\label{eqn:loss_reparam_pie}
   &  l(\mathbf{F})= \sumgraphu \Bigg(\log(p(\mathbf{f}_n)) + \\
   & \quad  \frac{1}{2} \Big(\sumneighbors{\log(e^{\flf} - 1)}  -  \mathbf{f}_n^\top  \mathbf{L}{\sumgraph{\mathbf{f}_m}} + \mathbf{f}_n^\top  \mathbf{L}\mathbf{f}_n \Big) \Bigg).
\end{aligned}
\end{equation}
\endgroup
We call this extension of IeClam PieClam (Prior Inclusive Exclusive Cluster Affiliation Model). We similarly extend BigClam to PClam (Prior Cluster Affiliation Model) by replacing $\mathbf{L}$ in (\ref{eqn:loss_reparam_pie}) with the identity matrix. 

Observe that the PieClam loss is similar to the IeClam loss only, with the addition of the prior, acting as a per node regularization term. The prior attracts all nodes to areas of higher probability during the optimization dynamics.

In order to sample from the above generative models, we first sample features $\{\mathbf{f}_n\}_{n\in [N]}$ according to $p$, and then connect them using either the BigClam or IeClam conditional probability.
%
To model the prior in practice, we use \emph{realNVP},  which is a \emph{normalizing flow} neural network model \cite{dinh2016density}. For more details on normalizing flows, see Appendix \ref{extended_sec:normflows}.



\paragraph{PieClam for graphs with node features.}
\label{sec:pieclam for graphs with node features}
%\textcolor{red}{Say that the prior maps to the direct product space of the affiliation space and the feature sapce.}

So far, we have used only the topology of the graph, not considering node features. 
We extend PieClam (and PClam) to graph-signals $([N], E, \mathbf{X})$ as follows. Concatenate the feature space of $\mathbf{X}$ with the affiliation space, and learn the prior on this combined space. This only affects the prior $p$. The conditional edge probabilities are  defined only in terms of the affiliation features, as before.

\section{Universality of PieClam and IeClam}

In this section, we define the universality of graph autoencoders, and prove that IeClam and PieClam are universal, while BigClam and PClam are not. The motivation behind the universality definition is that we would like to uniformly choose the dimension of the code space, such that every graph can be approximated by decoding some points in this fixed code space. Namely, we would like one \emph{universal} decoder that works for all graph, as opposed to choosing the dimension of the code space depending on the graph.

\subsection{General Graph Autoencoders}

Next, we define a general decoder that defines edge probabilities  by operating on pairs of points in a code space. 

\begin{definition}
   A \emph{pairwaise decoder} over the code spaces $\mathbb{R}^M$ is a mapping $D_M:\mathbb{R}^{2M}\rightarrow[0,1]$. Given $N$ points in the code space $\mathbf{z}=\{z_n\in\mathbb{R}^M\}_{n=1}^N$, the decoded graph $G_N(\mathbf{z})$ is the weighted graph with adjacency matrix  
\[\mathbf{D}_M(\mathbf{z})=\big(D_M(z_n,z_k)\big)_{n,k=1}^N.\]
\end{definition}

Clam models are special cases of pairwise decoders.

%The pairwise decoders of of the Clam models involve a construction with an exponential, which cannot attain probability 1.

\subsection{Log Cut Distance}

Our definition of universality has the following form: for every error tolerance $\epsilon>0$, there is a choice of the dimension $M(\epsilon)$ of the code space such that every graph can be approximated up to error $\epsilon$ by decoding some points in this space. To formalize the ``up to error $\epsilon$'' statement, we present in this section a new graph similarity measure which we call the log cut distance.
Our construction is based on a well-known graph similarity measure called the cut norm.

\begin{definition}
  The \emph{cut norm} of a matrix $\mathbf{X}\in\mathbb{R}^{N\times N}$ is defined to be
    \begin{equation}
    \label{eq:GraphCutNorm0}  \|\mathbf{X}\|_{\square}:=\frac{1}{N^2}\sup_{\mathcal{U},\mathcal{V}\subset[N]}\Big|\sum_{i\in \mathcal{U}}\sum_{j\in \mathcal{V}}x_{i,j}\Big|.
\end{equation}  
\end{definition}
The \emph{cut metric} $\|\mathbf{A}-\mathbf{B} \|_{\square}$ between two adjacency matrices $\mathbf{A}$ and $\mathbf{B}$ is interpreted as the difference between the edge densities of $\mathbf{A}$ and $\mathbf{B}$ on the block $\mathcal{U}\times\mathcal{V}$ on which their edge densities are the most different.

The following graph similarity measure modifies the cut norm, making it appropriate for graphs with random edges over a fixed node set.
\begin{definition}
\label{def:KL1}
Given two random graphs over the nodes set $[N]$, with independently Bernoulli distributed edges, with probabilities $\mathbf{P}=\{p_{n,m}\}_{n,m\in[N]}$ and $\mathbf{Q}=\{q_{n,m}\}_{n,m\in[N]}$ respectively, their \emph{log cut distance} is defined to be 
\begin{equation}
\footnotesize
    \begin{split}
       & D_{\square}(\mathbf{P} || \mathbf{Q}):=  \inf_{0< e,d\leq 1} \Bigg( e+d+ \\
       &  \quad \quad \  \frac{1}{N^2}\sup_{\mathcal{U},\mathcal{V}\subset[N]}\Big|\log\Big(\prod_{n\in \mathcal{U}}\prod_{m\in \mathcal{V}}
\frac{1-(1-e)p_{n,m}}{1-(1-d)q_{n,m}}\Big)\Big|\Bigg).
    \end{split}
    \label{eq:KL1}
\end{equation}
\end{definition}
The second term in (\ref{eq:KL1}) is the cut distance $\| \tilde{\mathbf{P}}-\tilde{\mathbf{Q}}\|_{\square}$ between the matrix $\tilde{\mathbf{P}}$ with entries 
\[\tilde{p}_{n,m}=-\log(1-(1-e)p_{n,m})\]
and the matrix $\tilde{\mathbf{Q}}$
 with entries 
 \[\tilde{q}_{n,m}=-\log(1-(1-d)q_{n,m}).\] Namely, the cut distance between the log likelihoods of non-edges. The parameters $e$ and $d$ make the $[0,1]$-valued probabilities valid inputs to the log. The goal of $e,d$ is to regularize the probability of the edges, where higher regularization is penalized via the additive term $e+d$ in (\ref{eq:KL1}).

%\ron{Change to cut relative entropy? Or something else? since we don't have the probability times log probability... Change it here an in the abstract to log cut distance?}
% The above similarity measure is related  to the \emph{Kullback-Leibler (KL) divergence}. Namely, 
For each choice of a cut $\mathcal{U},\mathcal{V}\subset[N]$,  the term
 \begin{equation}
 \label{eq:KL2}
     \frac{1}{N^2}\log\Big(\prod_{n\in \mathcal{U}}\prod_{m\in \mathcal{V}}
\frac{1-(1-e)p_{n,m}}{1-(1-d)q_{n,m}}\Big)
 \end{equation}
is somewhat similar in structure to an un-normalized KL divergence, or distance of log likelihoods, between the non-edge probabilies of the graphs $\mathbf{P}$ and $\mathbf{Q}$ over the dyads between  $\mathcal{U}$ and $\mathcal{V}$. Here, ``un-normalized'' means that the dyads are drawn  uniformly with probabilities $1/N^2$, but the sum of probabilities is $|\mathcal{U}|\cdot|\mathcal{V}|/N^2$ and not 1. The un-normalized uniform distribution discourages the supremum inside the definition of $D_{\square}$ from choosing small blocks for maximizing (\ref{eq:KL2}). Note that normalized uniform distributions would lead $D_{\square}$ to choose  small blocks, which do not reflect meaningful empirical estimates of the edge statistics (the edge densities of small blocks would not be interpretable as expected number of edges). To conclude, $D_{\square}(\mathbf{P} || \mathbf{Q})$ is interpreted as the maximal divergence between $\mathbf{P}$ and $\mathbf{Q}$ over all blocks, up to the best regularizers $e,d$.

In our analysis we compute the log cut distance between the random decoded graph $\mathbf{P}$ and the deterministic target graph $\mathbf{A}$. While $\mathbf{A}$ has edge probabilities in $\{0,1\}$, $\mathbf{P}$ has edge probabilities in $[0,1)$ for Clam models. Therefore, we only require regularization  for $\mathbf{A}$. We hence consider the following modified version of Definition \ref{def:KL1}.
\begin{definition}
\label{def:KL2}
Given an unweighted graph with adjacency matrix $\mathbf{A}\in\{0,1\}^{N\times N}$ and a random graph over the nodes set $[N]$, with independently Bernoulli distributed edges with probabilities $\mathbf{P}=\{0\leq p_{n,m}<1\}_{n,m\in[N]}$, the \emph{log cut distance} between $\mathbf{P}$ and $\mathbf{A}$ is defined to be
    \[\footnotesize
    \begin{split}
    & D_{\square}(\mathbf{P} || \mathbf{A}):= \\
    & \inf_{0< d\leq1} \Bigg( d+ \frac{1}{N^2}\sup_{\mathcal{U},\mathcal{V}\subset[N]}\Big|\log\Big(\prod_{n\in \mathcal{U}}\prod_{m\in \mathcal{V}}
\frac{1-p_{n,m}}{1-(1-d)a_{n,m}}\Big)\Big|\Bigg).
\end{split}\]
\end{definition}

Lastly, in case both $\mathbf{P}$ and $\mathbf{Q}$ are $[0,1)$-valued, a simple version of the log cut distance is (\ref{eq:KL1}) with the choice $e=d=0$, namely,
 \[D^0_{\square}(\mathbf{P} || \mathbf{Q}):=\]
 \begin{equation}
    \label{eq:KL3}
         \frac{1}{N^2}\sup_{\mathcal{U},\mathcal{V}\subset[N]}\Big|\log\Big(\prod_{n\in \mathcal{U}}\prod_{m\in \mathcal{V}}
\frac{1-p_{n,m}}{1-q_{n,m}}\Big)\Big|.
    \end{equation}



\subsection{Universal Graph Autoencoders}

We are now ready to define the universality of general pairwise autoencoders. Motivated by the fact that a Clam autoencoder is actually a family of autoencoders, parameterized by the number of communities, we also define general pairwise decoders as families.

\begin{definition}
    A family of code spaces $\mathbb{R}^M$ and corresponding pairwaise decoders $D_M:\mathbb{R}^{2M}\rightarrow[0,1]$, parametrized by $M\in\mathbb{N}$, is called  \emph{universal} if for every $\epsilon>0$ there is $M\in\mathbb{N}$ (which depends only on $\epsilon$) such that for every $N\in\mathbb{N}$ and every graph  with adjacency matrix $\mathbf{A}$ and $N$ nodes there are $N$ points in the code space $\{z_n\in\mathbb{R}^M\}_{n=1}^N$ such that 
     \[D_{\square}(\mathbf{D}_M(\mathbf{z}) || \mathbf{A}) <\epsilon.\] 
\end{definition}

\subsection{BigClam and PClam are Not Universal}
\label{BigClam is Not Universal}

We now show that BigClam (and hence also PClam) is not a universal autoencoder since it cannot approximate bipartite graphs.
Consider the bipartite graph $\mathbf{B}$ with $N$ nodes at each part, and probability $1-e^{-a^2}$ for an edge between the two parts, and $0$ within each part. 
Since in this case all probabilities are less than 1, we can use (\ref{eq:KL3}) as the definition of the log cut distance. The analysis for Definition \ref{def:KL2} extends naturally.

Let $\mathbf{P}$ be a decoded BigClam graph. 
Our goal is to show that there is no way to make $D^0_{\square}(\mathbf{P} || \mathbf{Q})$ small by choosing the dimension $C$ uniformly with respect to $N$. In fact, we will show BigClam cannot approximate a bipartite graph at all.\footnote{In Appendix \ref{BigClam With No Self Loops Approximating Bipartites} we show that one can approximate a bipartite graph of $2N$ nodes using $C=N^2$ classes in BigClam if the model ignores self-loops.} 

\begin{claim}
\label{claim:Clam_bi}
    Under the above construction, 
    \[D^0_{\square}(\mathbf{P} || \mathbf{B}) \geq \frac{a^2}{16}.\]
    As a result, BigClam is not a universal autoencoder.
\end{claim}
The proof is given in Appendix \ref{Proof That BigClam Is Not Universal}.
We note that one can similarly show that BigClam is not universal also with respect to the log cut distance of Definition \ref{def:KL2}.


\subsection{Universality of IeClam and PieClam}
We are now ready to show that IeClam (and hence also PieClam) is a universal autoencoder. The proofs of the following two theorems are in Appendix \ref{Proof of the Universality of PieClam} and \ref{Proof of the Universality of IeClam in a Cone of Non-negativity}.

We give two versions for the universality result. The first is without a cone restriction, and requires a relatively small number of communities for the given error tolerance. The corresponding decoder produces edge weights that can be negative. The second theorem is restricted to the pairwise cone of non-negativity, and has pessimistic  asymptotics for the required number of communities given an error tolerance. This decoder is guaranteed to produce proper edge probabilities in $[0,1)$.



\begin{theorem}
\label{thm:IEUniversality}
    For every epsilon $\epsilon>0$, every $N\in\mathbb{N}$, and every adjacency matrix $\mathbf{A}\in[0,1]^{N\times N}$, there are $N$ affiliation features $\mathbf{F}\in \mathbb{R}^{2K}$ of dimension  $K=-9\log(\epsilon/2)^2/\epsilon^2$ such that the corresponding IeClam model $\mathbf{P}=\{P(n\sim m|\mathbf{f}_n, \mathbf{f}_m)\}_{n,m=1}^N$ satisfies 
     \[D_{\square}(\mathbf{P} || \mathbf{A}) <\epsilon.\] 
    Here, the log cut distance is from Definition \ref{def:KL2}.
    As a result, 
    IeClam and PieClam are universal autoencoders with code space $\mathbb{R}^{2K}$.
\end{theorem}


\begin{theorem}
\label{thm:IEUniversality2}
    For every epsilon $\epsilon>0$, every $N\in\mathbb{N}$, and every adjacency matrix $\mathbf{A}\in[0,1]^{N\times N}$, there are $N$ affiliation features $\mathbf{F}$ in the cone of pairwise non-negativity  $\mathcal{T}\subset \mathbb{R}^{2C}$ of dimension  $C=2^{4\lceil -\log(\epsilon/2)^2/\epsilon^2\rceil}$  such that the corresponding IeClam model $\mathbf{P}=\{P(n\sim m|\mathbf{f}_n, \mathbf{f}_m)\}_{n,m=1}^N$ satisfies 
     \[D_{\square}(\mathbf{P} || \mathbf{A}) <\epsilon.\] 
    Here, the log cut distance is from Definition \ref{def:KL2}.
    As a result, 
    IeClam and PieClam are universal autoencoders with code space $\mathcal{T}$.
\end{theorem}


\section{Experiments}
\label{Experiments}

%\textcolor{cyan}{should also write about the fact that the model works well without node features. (actually, on the real anomalies, the features almost don't matter)}

%We conduct experiments that showcase and support our theoretical constructions.

\subsection{Reconstructing Synthetic Priors}
\label{sec:reconstruction_prior}

We consider a ground-truth synthetic prior $p:\mathbb{R}^C\rightarrow[0,\infty)$ in PClam. %, modeled via a realNVP normalizing flow \cite{dinh2016density}. 
We sample $N=500$ points from the prior, decode the corresponding PClam graph, and sample a simple graph from the random Bernoulli edges. Then, given these sampled graph, we fit to it a PClam and  model. In Figure~\ref{fig:prior_reconstruction} we compare the ground-truth prior to the reconstructed prior. We observe that even though our method is unsupervised, it still manages to capture the prior qualitatively well. More details are given in Appendix \ref{Extended Details on Experiments}.

\begin{figure}[!t]
    \begin{minipage}{0.3\columnwidth}
        \includegraphics[width=\textwidth]{CameraReady/LaTeX/images/r.png}
        
    \end{minipage}\hfill
    \begin{minipage}{0.3\columnwidth}
        \includegraphics[width=\textwidth]{CameraReady/LaTeX/images/c.png}
        
    \end{minipage}\hfill
    \begin{minipage}{0.3\columnwidth}
        \includegraphics[width=\textwidth]{CameraReady/LaTeX/images/ppp.png}
        
    \end{minipage}
    \caption{Left to right: Synthetic prior in an affiliation space of two inclusive communities.  Reconstructed prior by PClam with normalizing flow.  Reconstructed affiliation features by PClam.}
    \label{fig:prior_reconstruction}
\end{figure}

%\lipsum[1-5] % Dummy text


\subsection{Reconstructing Synthetic SBMs}
\label{sec:reconstruction_sbm}

In Figures \ref{fig:sbm_PieClam} and \ref{fig:sbm_pclam} we consider a synthetic SBM, and sample simple graph with $N=210$ nodes from it.  We then fit A PClam and PieClam model to it. %, and estimate the log cut distance\footnote{We estimate the cut distance underlying the log cut distance using the method from \cite{Alon_cut_estimation2004}.} between the Clam model and the SBM. 
The SBM is not dominant diagonal, so it cannot be well approximated by the PClam model (which finds a nested community structure), while the PieClam model approximates it well qualitatively. 
 %Figure~** illustrates that Clam and PClam models are unable to approximate SBMs which are not diagonally dominant, while IeClam and PieClam can. 
  Additional details are given in Appendix \ref{Extended Details on Experiments}.

\begin{figure}[!t]
    \begin{minipage}{0.29\columnwidth}
        \includegraphics[width=\textwidth]{CameraReady/LaTeX/images/e.png}
        
    \end{minipage}\hfill
    \begin{minipage}{0.29\columnwidth}
        \includegraphics[width=\textwidth]{CameraReady/LaTeX/images/d.png}
        
    \end{minipage}\hfill
    \begin{minipage}{0.2\columnwidth}
        \includegraphics[width=\textwidth]{CameraReady/LaTeX/images/j.png}
        
    \end{minipage}
    \begin{minipage}{0.2\columnwidth}
        \includegraphics[width=\textwidth]{CameraReady/LaTeX/images/k.png}
        
    \end{minipage}
    \caption{Left to right:  Adjacency matrix sampled from SBM with three classes and 9 blocks.  Adjacency matrix of the fitted PieClam graph, with two inclusive and two exclusive communities.  Affiliation features of the PieClam matrix in $\mathcal{T}$ projected to $(t^1,s^1)$ and $(t^2,s^2)$.}
    \label{fig:sbm_PieClam}
\end{figure}

\begin{figure}[!t]
    \begin{minipage}{0.3\columnwidth}
        \includegraphics[width=\textwidth]{CameraReady/LaTeX/images/i.png}
        
    \end{minipage}\hfill
    \begin{minipage}{0.3\columnwidth}
        \includegraphics[width=\textwidth]{CameraReady/LaTeX/images/h.png}
        
    \end{minipage}\hfill
    \begin{minipage}{0.3\columnwidth}
        \includegraphics[width=\textwidth]{CameraReady/LaTeX/images/f.png}
        
    \end{minipage}
    \caption{Left to right: Adjacency matrix sampled from SBM. Fitted  PClam adjacency matrix based on two inclusive communities.  Affiliations features of the PCLam matrix.}
    \label{fig:sbm_pclam}
\end{figure}

\subsection{Anomaly Detection}
\label{sec:anomaly_detection}

In unsupervised node anomaly detection, one is given a graph with node features, where some of the nodes are unknown anomalies. The goal is to detect these anomalous nodes, without supervising on any example of normal or anomalous nodes, by only using the structure of the graph and the node features.  
%
We use Clam models for node anomaly detection by fitting the Clam model to the graph and flagging nodes as anomalous if they satisfy the following different criteria.
\begin{itemize}
    \item 
    \textbf{(S)} \textbf{Star probability:} Given any Clam model, a node $n$ is called anomalous if $\prod_{m\in\mathcal{N}(n)}P(n\sim m| \mathbf{F})<\delta$. 
    \item 
    \textbf{(P)} \textbf{Prior probability:} Given any Clam model with prior, a node $n$ is called anomalous if $p(\mathbf{f}_n)<\delta$. 
    \item 
     \textbf{(PS)} \textbf{Prior star probability:} Given any Clam model with prior, a node $n$ is called anomalous if $p(\mathbf{f}_n)\prod_{m\in\mathcal{N}(n)}P(n\sim m| \mathbf{F})<\delta$. 
\end{itemize}
 We reduce the dimension of the node features of the input graphs to $100$ using truncated SVD, unless the dimension of the features is smaller than $100$ in which case we only normalize them to have zero mean and standard deviation one.   We use an affiliation space embedding dimension of 30 for $\mathbf{F}$ for PieClam and IeClam, and $24$ for BigClam.
Every Clam method starts with a random embedding $\mathbf{F}$.  Clam models with prior are trained with the following steps. 
 ($\mathbf{F}$-$t$): given a fixed prior $p$, optimize only $\mathbf{F}$ for $t$ steps. ($p$-$t$): given a fixed embedding $\mathbf{F}$, optimize only $p$ for $t$ steps. For regularization, in each iteration of of $p$ we add Gaussian noise with amplitude 0.01 to the affiliation features. We train PieClam with scheme $\mathbf{F}$-500 $\rightarrow$ $p$-1300 $\rightarrow$ $\mathbf{F}$-500 $p$-1300 with learning rate of $2e^{-6}$ on $\mathbf{F}$  and $1e^{-6}$ on $p$.  
For the models which only optimize $\mathbf{F}$ (IeClam and BigClam)  we use the following configurations: We train IeClam on 2500 iterations with learning rate 1e-6 using 30 communities. we train BIGClam on 2200 iterations with learning rate 1e-6.
More details on hyper-parameters are given in Appendix \ref{Extended Details on Experiments}.


In Table~\ref{Table1} we compare the performance of Clam methods to DOMINANT \cite{ding2019dominant}, AnomalyDAE \cite{fan2020anomalydae}, OCGNN \cite{wang2021ocgnn}, AEGIS \cite{ding2021aegis}, GAAN \cite{chen2020gaan} and TAM \cite{qiao2024tam} on the datasets Reddit, Elliptic, and Photo. The hyper-parameters of the competing methods are taken as the recommended values from the respective papers. The results are taken from Table 1 in \cite{qiao2024ggad}.
We observe that our methods are first and second place on all datasets. Moreover, S- IeClam, PS-PieClam and S BigClam each beats the competing methods in two out of the three datasets.


%elliptic
%accuracies_anomaly['vanilla_star'][-1]= 0.43547641111559626
%        accuracies_anomaly['prior'][-1]= 0.6092925420504949
%        accuracies_anomaly['prior_star'][-1]= 0.550598487402

%reddit
%accuracies_anomaly['vanilla_star'][-1]= %0.6104237623089773
 %       accuracies_anomaly['prior'][-1]= 0.5671389289452801
 %       accuracies_anomaly['prior_star'][-1]= 0.6124879959487293

%photo
%accuracies_anomaly['vanilla_star'][-1]= 0.5564428423968186
%        accuracies_anomaly['prior'][-1]= 0.42471794085192105
%        accuracies_anomaly['prior_star'][-1]= 0.5065087026473599

        

\begin{table}[t]
%\footnotesize
\centering
\begin{tabular}{|c|c|c|c|}
    \hline
  Method  & Reddit& Elliptic& Photo\\
    \hline

 (S)- IeClam& \textbf{0.639}& 0.440&\textbf{0.592}\\
 (S) - PieClam& 0.610& 0.435&0.556\\
 (P) - PieClam& 0.567& \textbf{0.610}&0.425\\
 (PS) - PieClam& *0.612& \underline{0.551}&0.507\\
 (S) - BigClam & \underline{0.637}& 0.434&\underline{0.581}\\
 \hline
 DOMINANT  & 0.511& 0.296&0.514\\
 AnomalyDAE &0.509& *0.496&0.507\\
 OCGNN  &0.525 & 0.258&0.531\\
 AEGIS &0.535 & 0.455&0.552\\
 GAAN &0.522 &0.259 &0.430\\
 TAM &0.606 & 0.404 & *0.568\\
 \hline
\end{tabular}
\caption{Comparison of Clam anomaly detectors with competing methods. % to DOMINANT \cite{ding2019dominant}, AnomalyDAE \cite{fan2020anomalydae}, OCGNN \cite{wang2021ocgnn}, AEGIS \cite{ding2021aegis}, GAAN \cite{chen2020gaan} and TAM \cite{qiao2024tam}. 
 First place in \textbf{boldface}, second with \underline{underline}, third with *star. We observe that our methods are first and second place on all datasets. Moreover, S- IeClam, PS-PieClam and S BigClam each beats the competing methods in two out of the three datasets. The accuracy metric is areas under curve (AUC).}
\label{Table1}
\end{table}

\begingroup
\color{blue}
\subsection{Link Prediction}
In supervised link prediction, one is given a graph for which some of the connections are unknown. The goal is to train the model using the known connections and predict the connectivity of the so called omitted dyads. In the following, we predict the probability of an edge using the conditional probability given in \ref{eq:LorFeature}. Since the prior is not used for prediction, there is no added benefit in concatenating the node features to the affiliation features as we did in \ref{sec:anomaly_detection} and so the link prediction reported in this paper \textbf{only uses the graph structure}. After assigning probabilities to all of the omitted dyads, we calculate the AUC score of the classification.
We use the experimental setting presented in \cite{zhou2022disenlink}. For every dataset, we omit $10\%$ of the edges as a test set and find the hyperparameters by talking a validation set of another $5\%$ for every hyperparameter configuration. In addition to omitting edges, we also omit 5 non-edges for every omitted edge. Finally, we run the model $10$ times on the test set with the fitted hyperparameters. We run our experiments on the datasets Squirrel and Texas \cite{pei2020texasquirrel}, Photo \cite{shchur2018photo} and Facebook100's Johns Hopkins \cite{lim2021hopkins} and compare PieClam to the baselines AA \cite{adamic2003friends}, VGAE \cite{kipf2016variational}, GAT \cite{GAT}, LINKX \cite{lim2021linkx} and Disenlink \cite{zhou2022disenlink}. The results are presented in Table \ref{table: link table}, and the baseline results are taken from Table 2 in \cite{zhou2022disenlink}. An in depth explanation on the experimental setup and a table with confidence intervals is given in the Appendix.
\begin{table}[t]
    \centering
    \begin{tabular}{|c|c|c|c|c|}
        \hline
        Method  & Squirrel& Photo & Texas & JH55 \\
        \hline
        PieClam & \textbf{98.8} & \textbf{98.4} &   \textbf{86.4}           & 95.6*     \\
        \hline
        BigClam & \underline{98.5}& 97.4*  & 78.8*& 94.9 \\
        VGAE    &   98.2   &94.9  & 68.6 & 92.8 \\
        GAT     &   98.0   &97.3  & 68.5 &   94.3\\
        LINKX   &   98.1   &97.0  & 75.8 &   93.4 \\
        AA      &   97.1 &97.4  & 53.1 &  \underline{96.1}  \\
        DisenLink & 98.3* & \underline{97.9} & \underline{80.7} &   \textbf{97.5} \\
        \hline
    \end{tabular}
    \caption{Comparison of Clam link predictors with competing models. First place in \textbf{boldface}, second with \underline{underline}, third with *star}
    \label{table: link table}
\end{table}

\endgroup


%We note that all of the competing models require the use of node attributes to recognize anomalies. One advantage of our models is that we often do not require the graph to have node attribute to achieve competitive results (in \ron{name the graphs in which we do not need features to be good}). 

\section{Conclusion}

We introduced PieClam, a new probabilistic graph generative model. PieClam models graphs via embedding the nodes into an inclusive and exclusive communities space, learning a prior distribution in this space, and decoding pairs of points in this space to edge probabilities, such that points are more likely to be connected the more inclusive communities and the less exclusive communities they share. We showed that PieClam is a universal autoencoder, able to approximate any graph, where the budget of parameters (the number of communities) can be predefined, irrespective of any property of specific graph, not even the number of nodes. Our experiments show that PieClam achieves competitive results when used in graph anomaly detection.

One limitation of PieClam is that, for attributed graphs, it only models the node features through the prior in the community affiliation space, but not via the conditional probabilities of the edges (given the community affiliations). Future work will deal with extending PieClam to also include the node (or edge) features in the edge conditional probabilities. 
Another limitation of our analysis is that the log cut distance is mainly appropriate for dense graphs. Future work will extend this metric to sparse graphs. This can be done, e.g., similarly to the sparse constructions in \cite{Levie_intersecting_reg_24}.




\bibliography{aaai25}

\[\]

\textbf{\Large Reproducibility Checklist}

This paper:

\begin{itemize}
    \item 
 Includes a conceptual outline and/or pseudocode description of AI methods introduced (\textbf{yes}/partial/no/NA)
    \item
Clearly delineates statements that are opinions, hypothesis, and speculation from objective facts and results (\textbf{yes}/no)
    \item
Provides well marked pedagogical references for less-familiare readers to gain background necessary to replicate the paper (\textbf{yes}/no)
\end{itemize}

Does this paper make theoretical contributions? (\textbf{yes}/no)

If yes, please complete the list below.
\begin{itemize}    
    \item
All assumptions and restrictions are stated clearly and formally. (\textbf{yes}/partial/no)
    \item
All novel claims are stated formally (e.g., in theorem statements). (\textbf{yes}/partial/no)
    \item
Proofs of all novel claims are included. (\textbf{yes}/partial/no)
    \item
Proof sketches or intuitions are given for complex and/or novel results. (\textbf{yes}/partial/no)
    \item
Appropriate citations to theoretical tools used are given. (\textbf{yes}/partial/no)
    \item
All theoretical claims are demonstrated empirically to hold. (\textbf{yes}/partial/no/NA)
    \item
All experimental code used to eliminate or disprove claims is included. (\textbf{yes}/no/NA)
\end{itemize}
Does this paper rely on one or more datasets? (\textbf{yes}/no)

If yes, please complete the list below.
\begin{itemize}
    \item
A motivation is given for why the experiments are conducted on the selected datasets (\textbf{yes}/partial/no/NA)
    \item
All novel datasets introduced in this paper are included in a data appendix. (\textbf{yes}/partial/no/NA)
    \item
All novel datasets introduced in this paper will be made publicly available upon publication of the paper with a license that allows free usage for research purposes. (\textbf{yes}/partial/no/NA)
    \item
All datasets drawn from the existing literature (potentially including authors’ own previously published work) are accompanied by appropriate citations. (\textbf{yes}/no/NA)
    \item
All datasets drawn from the existing literature (potentially including authors’ own previously published work) are publicly available. (\textbf{yes}/partial/no/NA)
    \item
All datasets that are not publicly available are described in detail, with explanation why publicly available alternatives are not scientifically satisficing. (yes/partial/no/\textbf{NA})
\end{itemize}
Does this paper include computational experiments? (\textbf{yes}/no)

If yes, please complete the list below.
\begin{itemize}
\item
Any code required for pre-processing data is included in the appendix. (\textbf{yes}/partial/no).
\item
All source code required for conducting and analyzing the experiments is included in a code appendix. (\textbf{yes}/partial/no)
    \item
All source code required for conducting and analyzing the experiments will be made publicly available upon publication of the paper with a license that allows free usage for research purposes. (\textbf{yes}/partial/no)
\item
All source code implementing new methods have comments detailing the implementation, with references to the paper where each step comes from (\textbf{yes}/partial/no)
\item
If an algorithm depends on randomness, then the method used for setting seeds is described in a way sufficient to allow replication of results. (\textbf{yes}/partial/no/NA)
\item
This paper specifies the computing infrastructure used for running experiments (hardware and software), including GPU/CPU models; amount of memory; operating system; names and versions of relevant software libraries and frameworks. (\textbf{yes}/partial/no)
\item
This paper formally describes evaluation metrics used and explains the motivation for choosing these metrics. (\textbf{yes}/partial/no)
\item
This paper states the number of algorithm runs used to compute each reported result. (\textbf{yes}/no)
\item
Analysis of experiments goes beyond single-dimensional summaries of performance (e.g., average; median) to include measures of variation, confidence, or other distributional information. (\textbf{yes}/no)
\item
The significance of any improvement or decrease in performance is judged using appropriate statistical tests (e.g., Wilcoxon signed-rank). (\textbf{yes}/partial/no)
\item
This paper lists all final (hyper-)parameters used for each model/algorithm in the paper’s experiments. (\textbf{yes}/partial/no/NA)
\item
This paper states the number and range of values tried per (hyper-) parameter during development of the paper, along with the criterion used for selecting the final parameter setting. (\textbf{yes}/partial/no/NA)
\end{itemize}


\clearpage

\appendix

\begin{center}
{\huge \textbf{Supplementary Material}}


\end{center}

\section{Proofs}

\subsection{Proof That BigClam Is Not Universal}
\label{Proof That BigClam Is Not Universal}
In this subsection we Prove Claim \ref{claim:Clam_bi}. 
Consider the bipartite graph $\mathbf{B}$ with $N$ nodes at each part, and probability $1-e^{-a^2}$ for an edge between the two parts, and $0$ within each part. Consider (\ref{eq:KL3}) as the definition of the log cut distance. Let $\mathbf{P}$ be a decoded BigClam graph from the affiliation features $\mathbf{F}$.

Denote by $\tilde{\mathbf{P}}$ the matrix with entries 
\[\tilde{p}_{n,m}=-\log(1-p_{n,m})=\mathbf{f}_n^\top\mathbf{f}_m,\] 
and by $\tilde{\mathbf{B}}$ the matrix  with entries $\tilde{b}_{n,m}=-\log(1-b_{n,m})$.
We show that there is no way to make $\| \tilde{\mathbf{P}}-\tilde{\mathbf{Q}}\|_{\square}$ small.

\begin{claim}
    Under the above construction, 
    \[D^0_{\square}(\mathbf{P} || \mathbf{B}) \geq \frac{a^2}{16}.\]
    As a result, BigClam is not a universal autoencoder.
\end{claim}


\begin{proof}
    
Note that $\tilde{b}_{n,m}=a^2$ if $n,m$ are in opposite parts, and $\tilde{b}_{n,m}=0$ if $n,m$ are on the same side.
We index the nodes such that $[N]$ is the first side of the graph, and $[N]+N$ is the second side. For $n\in[N]$ we denote $\mathbf{q}_n=\mathbf{f}_n$ and  $\mathbf{y}_n=\mathbf{f}_{n+N}$.
Next, we use the identity 
\[D^0_{\square}(\mathbf{P} || \mathbf{Q}) = \| \tilde{P}-\tilde{B}\|_{\square },\]
and bound the right-hand-side from below.

First, by the definition of cut norm, for every $\mathcal{U},\mathcal{V}\subset[2N]$,
\[\| \tilde{P}-\tilde{B}\|_{\square} \geq \Big|\frac{1}{4N^2}\sum_{n\in \mathcal{U}}\sum_{m\in \mathcal{V}} (\tilde{p}_{n,m} - \tilde{b}_{n,m})\Big|.\]
Hence, for $\mathcal{U}_1=\mathcal{V}_1=[N]$, $\mathcal{U}_2=\mathcal{V}_2=[N]+1$, $\mathcal{U}_3=[N],\mathcal{V}_3=[N]+N$, and $\mathcal{U}_4=[N]+N,\mathcal{V}_4=[N]$, we have
\[\| \tilde{P}-\tilde{B}\|_{\square} \geq \]
\[\frac{1}{16N^2}\sum_{j=1}^4 \Big|\sum_{n\in \mathcal{U}_j}\sum_{n\in \mathcal{V}_j} (\tilde{p}_{n,m} - \tilde{b}_{n,m})\Big| \]
\[= \frac{1}{16N^2}\sum_{n=1}^N\sum_{m=1}^{N} \mathbf{q}_n^\top \mathbf{q}_m  + \frac{1}{16N^2}\sum_{n=1}^N\sum_{m=1}^{N} \mathbf{y}_n^\top \mathbf{y}_m \]
\begin{equation}
    \label{eq:ClamNpUni0}
    +\frac{1}{8N^2}\sum_{n=1}^N\sum_{m=1}^{N} ( a^2 -\mathbf{q}_n^\top \mathbf{y}_m).
\end{equation}
Denote 
\[\mathbf{q}=\frac{1}{N}\sum_{n=1}^N \mathbf{q}_n, \quad \mathbf{y}=\frac{1}{N}\sum_{n=1}^N \mathbf{y}_n.\]
With these notations, (\ref{eq:ClamNpUni0}) can be written as
\[\| \tilde{P}-\tilde{B}\|_{\square} \geq \]
\[\frac{1}{16}\mathbf{q}^\top \mathbf{q}  + \frac{1}{16}\mathbf{y}^\top \mathbf{y} +\frac{1}{8}( a^2 -\mathbf{q}^\top \mathbf{y})\]
\[=\frac{1}{16}\Big(a^2 + (\mathbf{q}^\top-\mathbf{y}^\top)(\mathbf{q}-\mathbf{y})\Big) \geq \frac{a^2}{16}.\]
\end{proof}

We note that one can similarly show that BigClam is not universal also with respect to the log cut distance of Definition \ref{def:KL2}.


\begingroup
\color{blue}


\subsection{Proof That a Monotonous Inner Product Decoder Is Not Universal in Log Cut}
In this subsection we Prove Claim \ref{claim:Clam_bi}. 
Consider the bipartite graph $\mathbf{B}$ with $N$ nodes at each part, and probability $1$ for an edge between the two parts, and $0$ within each part. Consider Definition \ref{def:KL2} as the definition of the log cut distance.
Let $g$ such that $g: \mathbb{R} \rightarrow [0,1]$ and $g$ is ascending. 
% Also, let $M$, and $\Pi: \mathbb{R}^M \times \mathbb{R}^M \rightarrow [0,1]$ such that for any $\mathbf{a},\mathbf{b} \in \mathbb{R}^M$: $\Pi(\mathbf{a}, \mathbf{b}) = g(\mathbf{a}^\top \mathbf{b})$.
Let $\{f_n\}_{n=1}^{2N}$ be some affiliation vector set for $\mathbf{B}$ and let $\mathbf{P}$ be the matrix such that $p_{nm} = Pr(n\sim m \mid \mathbf{f}_n, \mathbf{f}_m) = g(\mathbf{f}_n^\top \mathbf{f}_m)$ be a random graph reconstruction of $\mathbf{B}$. Note that BigClam is a special case of this construction.

Denote by $\tilde{\mathbf{P}}$ the matrix with entries 
\[\tilde{p}_{n,m}=-\log(1-p_{n,m})=-\log(1-g(\mathbf{f}_n^\top \mathbf{f}_m)),\]
Let $d\in [0,1]$ and denote by $\tilde{\mathbf{B}}^d$ the matrix  with entries $\tilde{b}^d_{n,m}=-\log(1-(1-d)b_{n,m})$. \ron{What is $\mathbf{B}$? If it's the same as in my construction then why does your bound not depend on $a$?}
\textcolor{cyan}{In here I'm using \textbf{definition \ref{def:KL2}} for the log cut and so the bipartite does not depend on $a$. it's just a regular bipartite: $b_{mn} = 0$ if they are in the same part and 1 if they are in different parts}
\begin{claim}
    Under the above construction, 
    \[D_{\square}(\mathbf{P} || \mathbf{B}) \geq 4\frac{1}{16}.\]
    As a result, any autoencoder that uses a monotonous ascending function of the inner product between affiliation features for decoding is not a universal autoencoder.
\end{claim}


\begin{proof}
    
Note that $\tilde{b}^d_{n,m}=-log(d)$ if $n,m$ are in opposite sides, and $\tilde{b}^d_{n,m}=0$ if $n,m$ are on the same side.
We index the nodes such that $[N]$ is the first side of the graph $\mathcal{C}_1$, and $[N]+N$ is the second side $\mathcal{C}_2$. For $n\in[N]$ we denote $\mathbf{q}_n=\mathbf{f}_n$ and  $\mathbf{y}_n=\mathbf{f}_{n+N}$.
Next, we use the definition 
\[D_{\square}(\mathbf{P} || \mathbf{B}) = \underset{d\in[0,1]}{\inf}(d + \| \tilde{P}-\tilde{B}^d\|_{\square }),\]
and bound the right-hand-side for our arbitrarily selected $d$ from below:

First, by the definition of cut norm, for every $\mathcal{U},\mathcal{V}\subset[2N]$,
\[\| \tilde{P}-\tilde{B}^d\|_{\square} \geq \Big|\frac{1}{4N^2}\sum_{n\in \mathcal{U}}\sum_{m\in \mathcal{V}} (\tilde{p}_{n,m} - \tilde{b}^d_{n,m})\Big|.\]
Hence, for $\mathcal{U}_1=\mathcal{V}_1=\mathcal{C}_1$, $\mathcal{U}_2=\mathcal{V}_2=\mathcal{C}_2$, $\mathcal{U}_3=\mathcal{C}_1,\mathcal{V}_3=\mathcal{C}_2$, and $\mathcal{U}_4=\mathcal{C}_2,\mathcal{V}_4=\mathcal{C}_1$, we have

\[\| \tilde{P}-\tilde{B}^d\|_{\square} \geq \]
\[\frac{1}{16N^2}\sum_{j=1}^4 \Big|\sum_{n\in \mathcal{U}_j}\sum_{m\in \mathcal{V}_j} (\tilde{p}_{n,m} - \tilde{b}^d_{n,m})\Big| \]
\[= \frac{1}{16N^2}\bigg(\sum_{n=1}^N\sum_{m=1}^{N} -\log(1-g(\mathbf{q}_n^\top \mathbf{q}_m))  +  \]
\begin{equation}
    +\sum_{n=1}^N\sum_{m=1}^{N} -\log(1-g(\mathbf{y}_n^\top \mathbf{y}_m))
\end{equation}
\begin{equation}
    \label{eq:ClamNotUnibound}
    +2\Big| \sum_{n=1}^N\sum_{m=1}^{N} ( -log(d) +\log(1-g(\mathbf{q}_n^\top \mathbf{y}_m))\Big|\bigg).
\end{equation}
Denote 
\[\xi:[0,1] \rightarrow \mathbb{R}, \quad \xi(x) = -\log(1-g(x))\]
\[\langle \mathcal{C}_i, \mathcal{C}_i\rangle\ = \frac{1}{N^2}\sum_{n=\mathcal{C}_i}\sum_{m\in \mathcal{C}_j} \xi(\mathbf{f}_n^\top \mathbf{f}_m))\]
Notice that $\xi \geq 0$ and is ascending, and so, by matching each of the elements in the sums for $\langle \mathcal{C}_i, \mathcal{C}_j\rangle$ we can easily see that \textcolor{cyan}{((((DAN: easily see is because WLG $\mathbf{a}^\top\mathbf{b} \leq \mathbf{a}^\top \mathbf{a}$ so $\xi(\mathbf{a}^\top\mathbf{b}) \leq \xi(\mathbf{a}^\top \mathbf{a}) \leq \xi(\mathbf{a}^\top \mathbf{a}) + \xi(\mathbf{b}^\top \mathbf{b})$)))))}
\begin{equation}
    \label{eq:meanIneq}
    \langle \mathcal{C}_1, \mathcal{C}_2 \rangle \leq \langle \mathcal{C}_1, \mathcal{C}_1 \rangle + \langle \mathcal{C}_2, \mathcal{C}_2 \rangle  
\end{equation}
\[\langle \mathcal{C}_1, \mathcal{C}_2 \rangle = \langle \mathcal{C}_2, \mathcal{C}_1 \rangle \]
With these notations, (\ref{eq:ClamNotUnibound}) can be written as
\[\| \tilde{P}-\tilde{B}^d\|_{\square} \geq \]
\[\frac{1}{16}\bigg(\langle \mathcal{C}_1 \mathcal{C}_1 \rangle  + \langle \mathcal{C}_2 \mathcal{C}_2 \rangle +2\Big|( -\log(d) -\langle \mathcal{C}_1 \mathcal{C}_2 \rangle\Big|\bigg)\]
\textbf{Case 1}: If $-\log(d) \geq \langle \mathcal{C}_1 \mathcal{C}_2 \rangle$:
\[\| \tilde{P}-\tilde{B}^d\|_{\square} \geq \]
\[\frac{1}{16}\bigg(\langle \mathcal{C}_1 \mathcal{C}_1 \rangle  + \langle \mathcal{C}_2 \mathcal{C}_2 \rangle -2\log(d) -2\langle \mathcal{C}_1 \mathcal{C}_2 \rangle\bigg) \geq\]
\[\frac{1}{16}\bigg(-2\log(d) -\langle \mathcal{C}_1 \mathcal{C}_2 \rangle\bigg) \geq\]
\[\frac{1}{16}\bigg(-\log(d) + \langle \mathcal{C}_1 \mathcal{C}_2 \rangle - \langle \mathcal{C}_1 \mathcal{C}_2 \rangle\bigg) =\]
\[-\frac{1}{16}\log(d)\]

\textbf{Case 2}: If $-\log(d) \leq \langle \mathcal{C}_1 \mathcal{C}_2 \rangle$:
\[\| \tilde{P}-\tilde{B}^d\|_{\square} \geq \]
\[\frac{1}{16}\bigg(\langle \mathcal{C}_1 \mathcal{C}_1 \rangle  + \langle \mathcal{C}_2 \mathcal{C}_2 \rangle + 2\langle \mathcal{C}_1 \mathcal{C}_2 \rangle +2\log(d) \bigg) \geq\]
\[\frac{1}{16}\bigg(\langle \mathcal{C}_1 \mathcal{C}_1 \rangle + \langle \mathcal{C}_2 \mathcal{C}_2 \rangle\bigg) \geq\]
\[-\frac{1}{16}\log(d)\]
In both cases, by using Equation \ref{eq:meanIneq} we see that
\[d + \| \tilde{P}-\tilde{B}^d\|_{\square} \geq \]
\[d-\frac{1}{16}\log(d)\]
For $d\in[0,1]$, this function has one minimum at $d=1/16$ and so we conclude that for any $d\in[0,1]$
\[d + \| \tilde{P}-\tilde{B}^d\|_{\square} \geq 4\frac{1}{16}\]
This inequality holds for any choice of $d$ and so we conclude
\[D^0_{\square}(\mathbf{P} || \mathbf{B}) \geq 4\frac{1}{16}.\]
\end{proof}


\textbf{Note}: If $g$ is descending, use a graph with two fully connected communities of the same size as the two problems are equivalent.
\endgroup

\subsection{BigClam With No Self Loops Approximating Bipartite Graphs}
\label{BigClam With No Self Loops Approximating Bipartites}

Consider the above bipartite graph  $\mathbf{B}$ with $N$ nodes at each part, and probability $1-e^{-a^2}$ for an edge between the two parts, and $0$ within each part. 
If we redefine the BigClam decoder to have no self-loops, namely, $\mathbf{P}$ has entries 
$p_{n,m}=P(n\sim m|\mathbf{f}_n, \mathbf{f}_m)$ for $n\neq m$, and $p_{n,m}=0$ for $n=m$, then one can obtain a bipartite $\mathbf{P}$ with $C=N^2$ communities as follows. 

In the following analysis, an addition or multiplication of a set by a scalar is defined to be the addition or multiplication of every element in the set by this scalar.
Encode each node $n\in[N]$ in part 1 to $\mathbf{f}_n$ with
$f_n^c=a$ for $c\in [N]+n(N-1)$ and $f_n^c=0$ otherwise. Encode every node $n\in [N]+N$ from side 2 to $\mathbf{f}_n$ with
$f_n^c=a$ for $c\in N([N]-1)+n$ and $f_n^c=0$ otherwise. It is easy to see that the corresponding $\mathbf{P}$ is bipartite with edge probability between the parts being $1-e^{-a^2}$.






\subsection{Proof of the Universality of PieClam}
\label{Proof of the Universality of PieClam}

The proof is based on a version of the weak regularity lemma for intersecting communities. While the standard weak regularity lemma \cite{weakReg,Szemeredi_analyst} partitions the graph into disjoint communities, it is well known that allowing the communities to overlap allows using much less communities, which improves the asymptotics of the approximation. The regularity lemma was used in the context of graph machine learning in \cite{Levie_reg_23,Levie_intersecting_reg_24}. To formalize the relevant version of the weak regularity theorem for our analysis, we first need to cite a definition from \cite{Levie_intersecting_reg_24}.

\begin{definition}
    A (hard) intersecting community graph (ICG) with $N$ nodes and  $K$ communities is a matrix $\mathbf{C}\in\mathbb{R}^{N\times N}$ of the following form. There exist $K$ column vectors $\mathbf{Q}=\big(\mathbf{q}_k\in\{0,1\}^N\big))_{k=1}^K\in \{0,1\}^{N\times K}$ and $k$ coefficients $\mathbf{r}=(r_k\in\mathbb{R})_{k=1}^K\in\mathbb{R}^K$ such that
    \[\mathbf{C}= \mathbf{Q}{\rm diag}(\mathbf{r})\mathbf{Q}^{\top }.\]
\end{definition}

%\ron{Prove in the supplamentry material the regularity lemma for non-negative $A$ approximated by non-negatieve $C$. For this, we can replace the 0 values of $A$ by small $\delta>0$. Then, cancel a DC part (full community) to $C$... Or just go over the proof...}

%\ron{Follow the standard proof, and demand positivity of the approximant each time. For that follow the proof of the lemma in Hilbert space in the special case of functions. Seems like the Hilbert regularity lemma should work with this restriction. Indeed, I am looking for a step where there is not a big improvement.}

%\ron{However, the thing I am allowed to take inner product against in this lemma would also be something that together with prevous approximation would be positive. This is fine! Since I only want to take inner product against tensor indicators with positive coefficient!!! So it works!}

The following is a special case of the weak regularity lemma from \cite{Levie_intersecting_reg_24}, up to the small modification to the adjacency matrix, allowing it to have values in $[0,R]$ instead of $[0,1]$. %, it is well known as a step in the proof of the weak regularity lemma, see, .e.g  [Frieze, Fast. , Lovasz. Analyst]. For completness, we prove it in Appendix ***.

\begin{theorem}
\label{thm:reg_sprase}
   % Let $1\leq p\leq2$. % and $D\in\sN$.
     Let $\mathbf{A}\in[0,R]^{N\times N}$ be an adjacency matrix of a graph with $N$ nodes. Let $\epsilon>0$. Denote $K=\frac{9R^2}{4\epsilon^2}$. Then, there exists a hard ICG $\mathbf{C}$ with $K$ communities such that 
    \begin{equation}
        \label{eq:reg_sparse}
    \|\mathbf{A}-\mathbf{C}\|_{\square} \leq \epsilon.
    \end{equation}
\end{theorem}

\begin{proof}[Proof of Theorem \ref{thm:IEUniversality}]
 Let $\epsilon>0$. Let $\mathbf{A}\in[0,1]^{N\times N}$ be an adjacency matrix. Let $0<d\leq 1$.

Consider the matrix  $\tilde{\mathbf{A}}$ with entries
\[\tilde{a}_{n,m}=-\log(1-(1-d)a_{n,m}).\]
In the following construction, we build IeClam affiliation features $\mathbf{F}$ and we want
\[1-\exp(-\mathbf{f}_n^{\top}\mathbf{L}\mathbf{f}_m)\approx (1-d)a_{n,m}.\]
%So the regularity lemma should be against the matrix $\mathbf{X}$ with entries
%\[x_{u,v}=-\log(1-(1-d)a_{u,v}).\]
Note that $-\log(1-(1-d)a_{n,m})$ is increasing in $(1-d)a_{n,m}$. For $(1-d)a_{n,m}=0$ the value of this function is 0, and for $(1-d)a_{n,m}=1-d$ it is $R=-\log(d)$. Choose $d=\epsilon/2$. For this specific choice of $d$, if we replace $d$ by $\epsilon/2$ in the definition of $D_{\square}$ and omit the infimum, we get an upper bound of $D_{\square}(\mathbf{P},\mathbf{A})$.
 

%We now use the overlapping weak regularity lemma.


Using the overlaping weak regularity lemma, we approximate $\tilde{\mathbf{A}}$ by an ICG with 
\[K=\frac{-9\log(\epsilon/2)^2}{\epsilon^2}\]
communities,
\[\mathbf{C}= \mathbf{Q}{\rm diag}(\mathbf{r})\mathbf{Q}^{\top },\]
such that
\[\|\tilde{\mathbf{A}}-\mathbf{C}\|_{\square} \leq \epsilon/2.\]

Let $\mathbf{r}_+={\rm ReLU}(\mathbf{r})$ and $\mathbf{r}_-={\rm ReLU}(-\mathbf{r})$. Denote
\[\mathbf{C}_+=\mathbf{Q}{\rm diag}(\sqrt{\mathbf{r}_+})\]
and
\[\mathbf{C}_-=\mathbf{Q}{\rm diag}(\sqrt{\mathbf{r}_-}).\]
Denote the rows of $\mathbf{C}_+$ by $\mathbf{t}_n$ and the rows of $\mathbf{C}_+$ \dan{(it should be $C_-$)} by $\mathbf{s}_n$, for $n=1,\ldots,N$. For each $n\in[N]$ we concatenate $(\mathbf{t}_n,\mathbf{s}_n)$ to define the affiliation feature $\mathbf{f}_n$, with the first $K$ coordinates being the inclusive communities, and the last $K$ coordinates being the exclusive communities. Denote the corresponding IeClam matrix by $\mathbf{P}$.

It is easy to see that $\mathbf{f}_n^\top\mathbf{L}\mathbf{f}_m$ is the $(n,m)$ entry of $\mathbf{Q}{\rm diag}(\mathbf{r})\mathbf{Q}^{\top }$. 
%Hence, since $\mathbf{C}$ is non-negative, the set $\{\mathbf{f}_n\}_{n\in[N]}$ belongs to some cone $\mathbf{C}$.
This also proves that
\[\|-\log(1-(1-\epsilon)\mathbf{A})+\log(1-\mathbf{P})\|_{\square}\]
\[=\|\tilde{\mathbf{A}}-\mathbf{C}\|_{\square} \leq \epsilon/2.\] 
%The claim now follows by adding $\epsilon$ to the above formula, and then replacing $\epsilon$ by $\epsilon/2$.
%
We can now summarize
\[D_{\square}(\mathbf{P}||\mathbf{A}) \leq \]
\[\epsilon/2+\frac{1}{N^2}\sup_{\mathcal{U},\mathcal{V}\subset[N]}\Big|\log\Big(\prod_{n\in \mathcal{U}}\prod_{m\in \mathcal{V}}
\frac{1-p_{n,m}}{1-(1-\epsilon/2)a_{n,m}}\Big)\Big|\]
\[ =
\epsilon/2+\frac{1}{N^2}\sup_{\mathcal{U},\mathcal{V}\subset[N]}\Big|\sum_{n\in \mathcal{U}}\sum_{m\in \mathcal{V}}\Big(-\log\big(1-(1-\epsilon)a_{n,m}\big) \]
\[\quad \quad \quad \quad\quad \quad \quad \quad + \log\big(1-p_{n,m})\big)\Big)\Big|\]
\[=\epsilon/2 + \| \tilde{\mathbf{P}}-\mathbf{C} \|_{\square} = \epsilon. \]
   
\end{proof}


\subsection{Proof of the Universality of IeClam in the Pairwise Cone of Non-negativity}
\label{Proof of the Universality of IeClam in a Cone of Non-negativity}

For this result, we use the standard weak regularity lemma for non-intersecting classes.  It is based on the weak regularity lemma from \cite{weakReg,Szemeredi_analyst}, see also Lemma 9.3 and Corollary 9.13 from \cite{cut-homo3}.

\begin{definition}
    A \emph{block matrix} $\mathbf{B}$ with $K$ classes is a symmetric matrix $\mathbf{B}\in[0,\infty)^{N\times N}$ for which there exists a partition of $[N]$ into $K$ disjoint sets, called \emph{classes}, $\mathcal{C}_1,\ldots,\mathcal{C}_K$ (with $\cup\mathcal{C}_j=[N]$), such that for every pair of classes $i,j\in[K]$, there is a constant $c_{i,j}\geq 0$ such that  $b_{n,m}=c_{i,j}$ for any two nodes $n\in\mathcal{C}_i$ and $m\in\mathcal{C}_j$.
\end{definition}

%Since we need our approximating block matrix to be non-negative, we
%Denote the \emph{Frobenius} norm of a matrix $\mathbf{X}\in\mathbb{R}^{N\times N}$ by
%\[\|\mathbf{X}\|_{\rm F} = \sqrt{\sum_{n,m=1}^N x_{n,m}^2}.\]
\begin{theorem}
\label{thm:reg_sprase2}
   % Let $1\leq p\leq2$. % and $D\in\sN$.
     Let $\mathbf{A}\in[0,R]^{N\times N}$ be an adjacency matrix of a graph with $N$ nodes. Let $\epsilon>0$. Denote $K=2^{2\lceil R^2/\epsilon^2\rceil}$.         Then, there exists a block matrix  $\mathbf{B}$ with $K$ (disjoint) classes such that 
    \begin{equation}
        \label{eq:reg_sparse2}
    \|\mathbf{A}-\mathbf{B}\|_{\square} \leq \epsilon.
    \end{equation}
    %and 
    %\[\|\mathbf{A}-\mathbf{B}\|_{\rm F} = \min_{\mathbf{B}'} \|\mathbf{A}-\mathbf{B}'\|_{\rm F},\]
    %where the minimum is over the block matrices $\mathbf{B}'$ with $K$ (disjoint) classes. 
\end{theorem}
%Note that Theorem \cite{Levie_intersecting_reg_24} is for intersecting classes. Our Theorem \ref{thm:reg_sprase2} follows \cite{Levie_intersecting_reg_24}  by two simple modifications.
%First, by scaling $[0,R]$-valued adjacency matrices to be $[0,1]$ valued. Second, by taking deriving $K=2^M$ non-intersecting classes  $\mathcal{C}=\{\mathcal{C}_j\}_{j=1}^K$ from the $M=9R^2/4\epsilon^2$ intersecting classes $\mathcal{Q}=\{\mathcal{Q}_j\}_{j=1}^M$ via
% \[\tilde{\mathcal{C}}=\big\{C\subset [0,1] \ |\ \exists ~n\in[N]: ~C=\cap\{\mathcal{Q}_j\in\mathcal{Q}|n\in \mathcal{Q}_j\}\big\},\]
% and adding to $\tilde{\mathcal{C}}$ as many copies of the empty set to obtain $K=2^M$ disjoint sets $\mathcal{C}.$

%The last part of Theorem \ref{thm:reg_sprase2} is useful for showing that we may take $\mathbf{B}$ to be non-negative. 
%Indeed, since $\mathbf{A}$ is non-negative, replacing any negative block of $\mathbf{B}$ by a block of zeros can only improve the Frobenius approximation. 

We stress that Theorem \ref{thm:reg_sprase2} guarantees non-negative block values of $\mathbf{B}$, while in Theorem \ref{thm:reg_sprase}, in general, the matrix $\mathbf{C}$ may have negative entries.

%Hence, as a corollary of Theorem \ref{thm:reg_sprase2}, under the same assumptions, there exists a non-negative $\mathbf{B}$ that satisfies $\|\mathbf{A}-\mathbf{B}\|_{\square} \leq \epsilon$. 







\begin{proof}[Proof of Theorem \ref{thm:IEUniversality2}]

We start similarly to the proof of Theorem \ref{thm:IEUniversality}. Let $\epsilon>0$. Let $\mathbf{A}\in[0,1]^{N\times N}$ be an adjacency matrix. 
Consider the matrix  $\tilde{\mathbf{A}}\in[0,-\log(\epsilon/2)]$ with entries
\[\tilde{a}_{n,m}=-\log(1-(1-\epsilon/2)a_{n,m}).\]
In the following, we build IeClam affiliation features $\mathbf{F}$  such that
\[1-\exp(-\mathbf{f}_n^{\top}\mathbf{L}\mathbf{f}_m)\approx (1-\epsilon/2)a_{n,m}.\]

By the weak regularity lemma (Theorem \ref{thm:reg_sprase2}), we approximate there is a non-negative block matrix $\mathbf{B}$ with $K=2^{2\lceil -\log(\epsilon/2)^2/\epsilon^2\rceil}$ classes $\mathcal{C}_1,\ldots,\mathcal{C}_K$, such that
\[\|\tilde{\mathbf{A}}-\mathbf{B}\|_{\square} \leq \epsilon.\]

We now take the affiliation space to have $C=K^2$ inclusive communities  and $C=K^2$ exclusive communities.

For each $n\in[N]$, let $k_n$ be the class such that $c\in\mathcal{C}_{n_k}$.  For each pair of classes $i,j\in[K]$, let $c_{i,j}\geq 0$ denote the edge weight between $\mathcal{C}_i$ and $\mathcal{C}_j$.
%Build a feature for every block. $\mathbf{f}_n$ is:

%I have enough dimensions to make each $n$ that belong to the same class to interact as I wish with all of the other classes.

%Let $c_{i,j}$ be the block values. let $k_n$ be the class of node $n$. Then we build the features by:

The feature of each $n\in[N]$ at the inclusive channel $c=(K-1)k_n+k_n$ is $t_n^c=\sqrt{c_{k_n,k_n}}$. It is $t_n^c=\sqrt{c_{k_n,j}/4}$ at inclusive channels $c=(K-1)k_n+j$ and $c=(K-1)j+k_n$, and $s_n^c=-\sqrt{c_{k_n,j}/4}$ at exclusive channels $c=(K-1)k_n+j$ and $c=(K-1)j+k_n$. In all other channels $t^c_n$ and $s^c_n$ are zero. Note that $\mathbf{f}_n$ belongs to the cone of pairwise non-negativity $\mathcal{T}\subset \mathbb{R}^{2K^2}$.

It is now direct to see that  $\mathbf{f}_n^\top\mathbf{L}\mathbf{f_m}=c_{k_n,k_m}$.  As a result, as in the proof of Theorem \ref{thm:IEUniversality}, we get
\[D_{\square}(\mathbf{P}||\mathbf{A}) \leq \]
\[\epsilon/2+\frac{1}{N^2}\sup_{\mathcal{U},\mathcal{V}\subset[N]}\Big|\log\Big(\prod_{n\in \mathcal{U}}\prod_{m\in \mathcal{V}}
\frac{1-p_{n,m}}{1-(1-\epsilon/2)a_{n,m}}\Big)\Big|\]
\[ =
\epsilon/2+\frac{1}{N^2}\sup_{\mathcal{U},\mathcal{V}\subset[N]}\Big|\sum_{n\in \mathcal{U}}\sum_{m\in \mathcal{V}}\Big(-\log\big(1-(1-\epsilon)a_{n,m}\big) \]
\[\quad \quad \quad \quad\quad \quad \quad \quad + \log\big(1-p_{n,m})\big)\Big)\Big|\]
\[=\epsilon/2 + \| \tilde{\mathbf{P}}-\mathbf{B} \|_{\square} = \epsilon. \]
%$Kk_n+j$, and $Kj+k_n$ is $c_{i,j}\geq0$, and exclusive channels $Kk_n+j$ and $Kj+k_n$ if $c_{i,j}<0$.
%At inclusive $Kk_n+k_n$ it is $\sqrt{|c_{k_n,j}|}$ if $c_{k_n,j}\geq 0$. At exclusive $Kk_n+k_n$ it is $\sqrt{|c_{k_n,j}|}$ if $c_{k_n,j}< 0$.  the value at all other channels is zero.


%$[2K^2]$. the feature of $n$ is $\sqrt{|c_{k_n,j}|/2}$ at inclusive channels $Kk_n+j$ and $Kj+k_n$ is $c_{i,j}\geq0$, and exclusive channels $Kk_n+j$ and $Kj+k_n$ if $c_{i,j}<0$.
%At inclusive $Kk_n+k_n$ it is $\sqrt{|c_{k_n,j}|}$ if $c_{k_n,j}\geq 0$. At exclusive $Kk_n+k_n$ it is $\sqrt{|c_{k_n,j}|}$ if $c_{k_n,j}< 0$.  the value at all other channels is zero.
  
\end{proof}
\ron{formalize what is universality in log likelihood the following is a lemma which proves universality}
\begingroup
\color{blue}

\subsection{Proof that Convergence in Log Cut Implies Convergence in Log Likelihood}
For a binary matrix $A\in\{0,1\}^{N\times N}$ We denote 
    \begin{equation*}    
    \sum_{m\sim n} = \sum_{\substack{n\in [N] \\ m \in [N] \\ A_{nm} = 1}}
    \end{equation*}
    and
    \begin{equation}
    \label{eq:summation_notation}
    \sum_{m\not\sim n} = \sum_{\substack{n\in [N] \\ m \in [N] \\ A_{nm} = 0}}.
    \end{equation}
    
\begin{definition}
\label{def:log_likelihood}
    Let $A\in \{0,1\}^{N\times N}$ and $P \in [0,1)^{N\times N}.$ 
    The \textbf{log likelihood loss} between $A$ and $P$ is defined to be:
    \[l(P,A) = \sum_{m\sim n } \log(p_{nm}) + \sum_{m\not\sim n}\log(1 - p_{nm}).\]
    \ron{For this definition, you need to assume $p_{n,m}>0$ for every $n\sim m$. You can also assume simply that $p_{n,m}>0$ for every $n,m$, as our methods satisfy this.}  
\end{definition}
\begin{lemma}
    Let $\frac{1}{N^2}>\epsilon > 0$, $A\in \{0,1\}^{N\times N}$ and $P \in [0,1)^{N\times N}$ such that $D_\Box(P\|A) < \epsilon$, then 
    \[|{\rm ll}(P,A)| < 5N^2E\epsilon\] 
    where $E$ is the number of edges in $A$. 
\end{lemma}

\ron{This does not seem to give universality. Given $\epsilon$, it is not true that you can choose a dimension of the code space such that the log likelihood has distance less than $\epsilon$ from zero, independently of $N$. For any $\epsilon$, if $N>\frac{1}{\sqrt{\epsilon}}$ then your theorem does not give anything...}

\begin{proof}
We start by proving a number of claims.
    \begin{claim}
    There exists $0<d\leq 1$ such that \[\epsilon > d + \|\log(1-p_{nm}) - \log(1 - (1-d)a_{nm})\|_\Box\]
    \label{claim:d_exists}
    \end{claim}
    \begin{proof}
        If we assume the opposite towards contradiction we get that for every $d$
        \[\epsilon \leq d + \|\log(1-p_{nm}) - \log(1 - (1-d)a_{nm})\|_\Box\]
        and therefore $ \epsilon \leq D_\Box(P||A) $ in contradiction. 
    \end{proof}
    \begin{claim}
    \label{claim:elementwise_cut}
        Let $S$ and $T$ such that 
        \[\epsilon > \|S-T\|_\Box\]
        then 
        \[\forall (n,m) \in [N]^2: \epsilon > \frac{1}{N^2}|S_{nm} - T_{nm}|\]
    \end{claim}
    \begin{proof}
        For the sets $\{m\}$ and $\{n\}$ we get that
        \[\epsilon > \|S-T\|_\Box = \frac{1}{N^2}\underset{U,V}{\sup} |\sum\limits_{k \in U} \sum\limits_{l \in V} (S_{kl} - T_{kl})| \geq \]

        \[\frac{1}{N^2}|\sum\limits_{k \in \{n\}}\sum\limits_{l \in \{m\}} (S_{kl} - T_{kl})| =  \frac{1}{N^2}|S_{nm} - T_{nm}|.\] 
    \end{proof}
    By Claim \ref{claim:d_exists} we deduce that there exists $d$ such that 
    \begin{align}
    \label{eq:epsilon_d_epsilon_sup}
        \varepsilon &> d\notag \\
        \varepsilon &> \|\log(1-p_{nm}) - \log(1-(1-d)a_{nm})\|_\Box
    \end{align}
    
    From Equation \ref{eq:epsilon_d_epsilon_sup} and the supremum definition
    \begingroup
    \footnotesize
    \begin{align}
    \label{eq:ll_epsilon}
        \epsilon &> \frac{1}{N^2}\Big|\sum_{\substack{n\sim m}}(\log(1-p_{nm}) - \log(d)) + \sum_{\substack{n\not\sim m}}\log(1-p_{nm}) \Big| =\notag\\ 
        = & \frac{1}{N^2}\Big|\sum_{\substack{n\sim m}}\big(-\log(d) + \log(1-p_{nm})\big) - \notag\\ 
        &\quad\quad\quad\quad\quad\quad - \sum_{\substack{n\sim m}}\log(p_{nm}) - |ll(P,A)|\Big|
    \end{align}
    \endgroup
    % And by Claim \ref{claim:elementwise_cut} it follows that for all $(n,m) \in [N]^2$
    % \begin{equation}
    %     \label{eq:log1_pnm}
    % \epsilon > \frac{1}{N^2}\Big|\log(1-p_{nm}) - \log(1-(1-d)a_{nm})\Big|.
    % \end{equation}
    To handle the $(-\underset{{n\sim m}}{\sum}\log(p_{nm}))$ term in Equation \ref{eq:ll_epsilon} we prove the following claim.
    \begin{claim}
    \label{claim:log_pnm}
    $|\log(p_{nm})| < 3\epsilon$
    \end{claim} 
    \begin{proof}
        We use Claim \ref{claim:elementwise_cut} on the pairs $(n,m) \in [N]^2$ such that $a_{nm}= 1$, so for every such pair it holds that
        % , in which case Equation \ref{eq:log1_pnm} reads
        \begin{equation}
        \label{eq:neighbors_small}    
        \epsilon > \frac{1}{N^2}\Big|\log(1-p_{nm}) - \log(d)\Big|
        \end{equation}
        If $\log(d) < \log(1-p_{nm})$, we get that
        \[\frac{1-p_{nm}}{d} < e^{N^2\epsilon}\]
        which implies (recalling that $\epsilon < \frac{1}{N^2}$)
        \begin{align*}
          -\log(&p_{nm}) < -\log(1-de^{N^2\epsilon}) < -\log(1 - \epsilon e^{N^2\epsilon})< \notag \\
          & < \epsilon e^{N^2\epsilon} - \mathcal{O}\Big( ( \epsilon e^{N^2\epsilon^2}) \Big) < 3\epsilon
        \end{align*}

    If $\log(d) > \log(1-p_{nm})$ then
    \begin{align}
    &-\log(1-p_{nm}) > -\log(\epsilon)\notag \\
    & \iff p_{nm} > 1-\epsilon \notag\\
    & \iff -\log(p_{nm}) < \epsilon - \mathcal{O}(\epsilon^2) < \epsilon       \end{align}
    \end{proof}

    We now use Equation \ref{eq:neighbors_small} and Claim \ref{claim:log_pnm} on Equation \ref{eq:ll_epsilon} to bound $|ll(P,A)|$. By the triangle inequality
    \ron{$4N^2\epsilon$ in the left hand side?}
    \begin{align}
        &N^2\epsilon >   \bigg|\Big|\sum_{\substack{n\sim m}}\big(-\log(d) + \log(1-p_{nm})\big) - \sum_{\substack{n\sim m}}\log(p_{nm})\Big| \notag\\ 
        &\quad\quad\quad\quad\quad\quad  - |ll(P,A)|\bigg|.
    \end{align}
    Finally
    \[|ll(P,A)| < (N^2 + N^2E + 3E)\epsilon< 5N^2E\epsilon.\]
    

\end{proof}

\subsection{Proof that a Sigmoid Inner Product Decoder Is Not Universal in Log Likelihood}
Consider the bipartite graph $\mathbf{B}\in \{0,1\}^{2N\times 2N} = (b_{nm})_{n,m \in [2N]}$ with $N$ nodes in each part and denote the parts by $\mathcal{C}_i$ where $i\in \{1,2\}$. Arrange the nodes such that $\mathcal{C}_1 = [N]$ and $\mathcal{C}_2 = [N]+ N$. 



Consider Definition (\ref{def:log_likelihood}) as the definition of the log likelihood. Let $\sigma: \mathbb{R} \rightarrow [0,1]$ be an ascending function.
Also, let $K$, and $\Pi: \mathbb{R}^K \times \mathbb{R}^K \rightarrow [0,1]$ such that $\Pi(\mathbf{f}_n, \mathbf{f}_m) = \sigma(\mathbf{f}_n^\top \mathbf{f}_m)$.
Let $\mathbf{P}$ such that $p_{nm} = Pr(n\sim m \mid \mathbf{f}_n, \mathbf{f}_m) = \Pi(\mathbf{f}_n, \mathbf{f}_m)$ be a random graph reconstruction from the affiliation features $\mathbf{F}$. Note that BigClam is a special case of this construction.

\begin{claim}
    Under the above construction, 
    \[|ll(P,B)| \geq -\log(\frac{1}{2})\]
    As a result, any autoencoder that uses a monotonous ascending function of the inner product between affiliation features for decoding is not universal.
\end{claim}
 
   \begin{proof}
    
For $n\in[N]$ we denote $\mathbf{q}_n=\mathbf{f}_n$ and  $\mathbf{y}_n=\mathbf{f}_{n+N}$.
Let $a > 0$ such that $|ll(P, B)| \leq a$.
Therefore
    \begingroup
    \footnotesize
    \begin{align}
    a >&  -2\sum\limits_{n=1}^N\sum\limits_{m=1}^N\log(\sigma(\mathbf{q}_n^\top\mathbf{y}_m)) - \sum\limits_{n=1}^N\sum\limits_{m=1}^N\log(1 - \sigma(\mathbf{y}_n^\top \mathbf{y}_m)) \notag\\
    &- \sum\limits_{n=1}^N\sum\limits_{m=1}^N\log(1 - \sigma(\mathbf{q}_n^\top \mathbf{q}_m))
    \end{align}
    \endgroup
Since all terms are of the same sign it holds that each of them is smaller than $a$.
\begin{equation}
\label{eq:small_1st}
-\log(\sigma(\mathbf{q}_n^\top \mathbf{y}_m))< a
\end{equation}

and (WLG for $\mathbf{q}$)
\begin{equation}
\label{eq:small_2nd}
-\log(1-\sigma(\mathbf{q}_n^\top \mathbf{q}_m))< a
\end{equation}
So for every $n,m \in [N]$
\[\sigma^{-1}(1-e^{-a})>\max(\|\mathbf{q}_n\|^2, \|\mathbf{y}_m\|^2)>\mathbf{q}_n^\top\mathbf{y}_m>\sigma^{-1}(e^{-a})\]
$\sigma^{-1}: [0,1] \rightarrow \mathbb{R}$ is also ascending so we must have $1-e^{-a} > e^{-a} \iff a> \log(\frac{1}{2})$ which means that for $a$ to be larger than $|ll(P,B)|$, $a$ has to be larger than $\log(\frac{1}{2})$ and so
\[|ll(P,B)|\geq\log(\frac{1}{2})\]

\end{proof}

\endgroup

\section{Extended Related Work}
\label{Extended Related Work}





\subsection{Message Passing Algorithms and Networks}


 %Message Passing Neural Networks (MPNNs) process graphs by iteratively updating the features of each node using data from its adjacent nodes. Information is exchanged between nodes along the graph's edges, a process referred to as \emph{message passing}. Each node combines the incoming messages from its neighbors using an \emph{aggregation scheme}, with common methods being summing, averaging, or taking the coordinate-wise maximum of the messages. Most graph neural networks applied in practice are specific instances of MPNNs \cite{MPNN,GCN,GAT}.

%\ron{Just keep it this way I will write it}

The message passing algorithm is a general architecture for processing graph-signals. An MPNN operates on the graph data by aggregating the features in the neighborhood of each node, allowing for information to travel along the edges. The first example of this scheme (Message Passing) was originally suggested by Pearl et al \cite{pearl1982belief},  and was combined with a neural network in \cite{duvenaud2015convolutional}, and later generalized in \cite{gilmer2017neural}. Most graph neural networks applied in practice are specific instances of MPNNs \cite{MPNN,GCN,GAT}. %In this architecture, the features of each node act as a hidden state that is updated by the following equations.

In MPNNs, information is exchanged between nodes along the graph's edges. Each node combines the incoming messages from its neighbors using an \emph{aggregation scheme}, with common methods being summing, averaging, or taking the coordinate-wise maximum of the messages.
Let $T \in \mathbb{N}$ represent the number of layers, and define two sequences of positive integers $(c_t)_{t=0}^T$ and $(d_t)_{t=0}^T$ representing the feature dimensions in the hidden layers $\{\mathcal{G}_t\}_{t=0}^T = \{\mathbb{R}^{c_t}\}_{t=0}^T$ and  $\{\mathcal{S}_t\}_{t=0}^T = \{\mathbb{R}^{d_t}\}_{t=0}^T$. Define the message functions as  $M^t : \mathcal{S}_t \times \mathcal{S}_t \rightarrow \mathcal{G}_t$  and unpdate functions as $U^t: \mathcal{G}_t \times \mathcal{S}_t \rightarrow \mathcal{S}_{t+1}$.
The features $\mathbf{f}_n^{t+1} \in \mathcal{S}_{t+1}$ at layer $t+1$ of the nodes $n\in[N]$ are computed from the features $\mathbf{f}_m^t  \in \mathcal{S}_t$ by 
\begin{align*}
    \mathbf{m}^{t}_n &= \underset{k\in \mathcal{N}(n)}{\sum}{M^t(\mathbf{f}^t_n,\mathbf{f}^t_k)}\\
    \mathbf{f}^{t+1}_n &= U^t(\mathbf{m}^t_n, \mathbf{f}^t_n).
\end{align*}
Here, $M^t$, $U^t$ and $m^t$ are called the message function, update function and mail at time $t$ respectively.  The summation over the messages can also be replaced for any node by any function $Agg_n^t: \prod_{|\mathcal{N}(n)|} \mathcal{G}_t \rightarrow \mathcal{G}_t$
which is permutation invariant. 
%At each iteration, data flows into each node from it's immediate proximity. 

At step $T$ a readout space can be defined $\mathcal{R}_T = \mathbb{R}^{b_T}$, with a permutation invariant readout function $R^T$, e.g., summing, averaging, or taking the max of all of the nodes of the graph. %so an intermediate calculation can be performed on the entire graph with some order invariant function:
%\begin{align}
%    \hat{y}^t = R^t(\{\mathbf{f}^{t}_n|n\in [N]\}).
%\end{align}
This produces a vector representation of the whole graph. %vector out of the entire graph.
%If any of $M^t$, $U^t$ or $R^t$ is a neural network, the algorithm is referred to as a message passing neural network or MPNN. The message passing algorithm can also be applied to unattributed graphs by setting each node with a local structural feature (e.g. the node's degree).



 

%\paragraph{Graph Representation Learning}
%Graph Representation Learning, also known as Graph Embedding, is a subfield of machine learning that focuses on learning vector representations for entities in a graph, such as nodes, edges, or subgraphs. \cite{hamilton2017representation} \ron{Add: [Graph Representation Learning: A Survey
%FENXIAO CHEN, YUNCHENG WANG, BIN WANG AND C.-C. JAY KU]}. 

%\ron{Write some things that you can do with this embedding. Why do we want to ember the graph? Extend the discussion I wrote in the introduction.}

%\ron{You have to talk about GNNs in this context. This is the most popular type of graph representation learning architecgture.}

% I already wrote about graph representation learning in the intro, so no need to write again here.



 
%In this work we focus on graph autoencoders based on node embedding. \ron{(,} and on unsupervised models. \ron{(? Extend or delete)}

\subsection{Deep Generative Models}
\label{extended_sec:generative_models}
Generative models in machine learning assume that training data is generated by some underlying probability distribution. One goal in this context is to approximate this distribution, or build a model that approximates a random sampler of data points from this distribution. Hence, generative models can be used to generate synthetic data, mimicking training data by sampling from the distribution 
\cite{harshvardhan2020comprehensive,dinh2016density,mohamed2016learning}, 
or to infer the probability of unseen data by substituting it into the probability function. 
The latter can be useful in tasks like anomaly detection \cite{pang2021anomalynotgraph} in which the model can asses whether a sample is probable under the learned model. 
Two examples of generative models are Generative Adversarial Networks (GANs) \cite{goodfellow2014gan,radford2015dcgan} and Variational Autoencoders (VAEs) \cite{kingma2013vae}, which are used both for inference and generation. 

VAE models consist of an encoder and a decoder. The encoder maps data from a high-dimensional \emph{data space} to a lower-dimensional, simpler, \emph{code space}. The decoder reverses this process, transforming data from the code space back into the data space. The code space serves as a bottleneck, capturing the essential features of the training data, which is often high-dimensional (e.g., images, social network graphs) and thus more complex than the code space. If the model trains by encoding followed by decoding, then minimizing the difference between the input and its decoded version, it is called an \emph{autoencoder}.

In a VAE, training involves encoding each data point to a known distribution (typically Gaussian), sampling from this distribution, decoding it, and then minimizing the difference between the distribution of the original data and the decoded data. For a survey on VAEs, see \cite{tschannen2018vaesurvay}. 
%In general, an autoencoder refe \ron{??}

%AGAN is another type of generative model. It is defined by pitting the encoder and decoder against each other in a ``min-max game'' loss. \ron{Why is GAN an autoencoder? You train one model to classify normal or fake inputs, and another to generate fake inputs. So where is the autoencoder here?} In this setup, the decoder is called the generator, and the encoder is the discriminator \ron{Is this how GAN's work?} The generator creates samples from the code space, while the discriminator assigns a probability of realness to images it receives. During training, the discriminator learns to distinguish between real and generated images, while the generator learns to fool the discriminator. For more information, see \cite{creswell2018gan_review}.

\subsection{Graph Generative Models}
%When dealing with graph type data, generative models provide an interpretable framework for many tasks.  %Graph representation learning, which is a field that generalizes graph generative models, by transforming the graph data into a lower dimensional vector space, while preserving the graph's structural and attribute information \cite{hamilton2017representation}, \cite{yang2013bigclam} \ron{You are very repetitive... We already discussed generative models.}. 
%Building on the foundation of graph representation learning, 
Graph generative models learn a probability distributions of graphs. % on the code space, making it possible for them to create new graphs based on a code space representation, and to estimate the probability of different graph structures under the training data \ron{What code space!? Not all generative models are autoencodels.}. 
 Such models allow for various tasks %that representation models cannot achieve 
  where the goal is not only to analyze existing graphs but also to predict or simulate new graph data.

Some classical generative models are pre-defined probabilistic models, e.g., the Erdős–Rényi model, Preferential attachment, Watts–Strogatz model, and more. See \cite{newman2003classical_graph_generative}  for a review. Other graph generative models are learned from data, e.g., \cite{lee2019review, kipf2016variational,you2018graphrnn,bojchevski2018netgan}. %\ron{LIST OF CITATIONS}\textcolor{cyan}{cited}.  %SBMs   are learned by defining a  
%\ron{We already have a section about SBMs! Whay are you writing everything twice?}
%probabilistic model given a graph which achieves maximum probability on the training graph, usually by defining a log likelihood function and maximizing it with various variational methods,\ron{more citations.} such as Expectation Minimization, \cite{holland1983sbm}, ,maximum likelihood estimation (MLE) \cite{karrer2011dcsbm}, Evidence Lower Bound  \cite{airoldi2008mmsbm} and other variational methods.

%\ron{say how they are learned}. 
%\ron{Explain: Autoencoder} \textcolor{cyan}{I explained it before explaining VAE}, \ron{any other way? Survey.} \textcolor{cyan}{not sure what you mean here, there are two reviews.}

%\ron{Why twice?! You already discussed applications above. So integrate these two parts to one paragraph that will be here.} \textcolor{cyan}{done}
Applications of generative graph models include social network analysis \cite{wang2018graphgan,grover2019graphite,harshvardhan2020comprehensive}.  anomaly detection \ref{extended_sec:anomaly_detection}, graph synthesis \cite{you2018graphrnn, guo2022systematic}, data augmentation \cite{ding2019data_augmentation}, and protein interaction modeling (e.g. for drug manufacturing) \cite{de2018molgan,ingraham2019generative}, to name a few.  For a review see \cite{ma2021comprehensive}.   

%Probabilistic generative models for graphs have a rich history including famous models such as stochastic block models (SBMs) and Bayesian network models \cite{newman2018networks} \ron{write the reference directly after its respective text. Otherwise we don't know which reference refer to what.}. Worth mentioning in this case are SBMs [\cite{lee2019review}] and the BigClam [\cite{Overlapping2012}] on which we will elaborate in a following section as they have a deep relation to PieClam.

\paragraph{Deep Graph Autoencoders.}
\label{sec:deep graph autoencoders}
In graph deep learning-based autoencoders,  one estimates the data distribution by learning to embed the nodes to a code space, in which the data distribution is defined to be some standard distribution,  e.g., Gaussian, in such a way that the encoded nodes can be recunstructed back to the graph with small error. Graph VAEs \cite{kipf2016variational,grover2019graphite,JMLR:v21:19-671,mehta2019stochastic} embed the data into the code space by minimizing the evidence lower bound loss comprised of the decoding loss and the KL divergence between the encoded distribution and a Gaussian prior. %, modulating both the recreation and distribution in the latent space. \ron{Recreation is a word used for hobbies and such}

\paragraph{GAN-based Graph generative models.}

In  \cite{wang2018graphgan}, a GAN method for graphs generates a neighborhood for each of the nodes and the discriminator gives a probability score for each edge. This method also formulate a graph version of softmax, and offer a random walk based generating strategy. Another GAN model that is used for anomaly detection is GAAN \cite{chen2020gaan}. In GAAN, the ground truth and generated node attributes are encoded into a latent space from which the adjacency between any two nodes is decoded using a sigmoid of the inner product between their latent features. 

\paragraph{Normalizing Flows-based Graph Models.}
Normalizing flows models \cite{dinh2016density}  also have adaptations for graph data. For example, the work of \cite{liu2019graph,madhawa2019graphnvp} offers a version of coupling blocks that use message passing neural networks.  See more details on normalizing flows in Section \ref{extended_sec:normflows}.

%\paragraph{Graph Generative Models as Message Passing.}

%Many generative graph models employ an MPNN scheme for encoding node features. and indeed, all of the models described above can be formulized as MPNNs if they are not formulized that way already \textcolor{cyan}{should i show on all of them? which ones to show?} \ron{This is a bold statement. Are you sure? All of them??}\\danny{pretty sure. from what i saw it was only GAAN that wasn't directly, but it multiplied by the adjacency which i think should make it an mpnn as only neighbors affect each feature. it is pretty bold, if we can leave it be, let's.}

\subsection{Stochastic block models} 
\label{Stochastic block models} 

A stochastic block model (SBM) is a generative model for random graphs. 
A basic stochastic block model is defined by specifying a number of \emph{classes} $K$, the probability $p_k$ of a random node being in block $k$, for $k\in[K]$, where $\sum_kp_k=1$, and an array of values $\mathbf{C}=\{c_{k,l}\}_{k,l=1}^J\in [0,1]^{J\times K}$ indicating edge probabilities between classes. Each node  of a randomly generated graph of $N$ nodes is independently chosen to belong to one of the classes at random, with probabilities  $\{p_k\}_{k\in[K]}$.  Then, the edges of the graph are chosen independently at random according to the following rule. For each $n\in[N]$, denote by $k_n\in[K]$ the class of $n$. Each dyad $(n,m)\in[N]^2$ is chosen to be an edge in probability $c_{k_n,k_m}$. Namely, the entries $a_{n,m}$ of the adjacency matrix of the random graph are independent random Bernoulli variables. See \cite{lee2019review} for a review on SBMs.

For each node $n$, denote by $\mathbf{f}_n\in\{0,1\}^K$ the vector such that $f_n^c=1$ if and only if $k_n=c$. Hence, in a basic SBM, the presence of an edge between nodes $n$ and $m$ follows a Bernoulli distribution with parameter $P(n\sim m|\mathbf{f}_m,\mathbf{f}_m,\mathbf{C}) = {\mathbf{f}_n}^\top  \mathbf{C} \mathbf{f}_m$, where the adjacency matrix is $\mathbf{P} = \mathbf{F}  \mathbf{C} \mathbf{F}^\top$, where $\mathbf{F}\in\mathbb{R}^{N\times K}$ is the matrix where each row $n$ has the feature $\mathbf{f}_n$ \cite{nowicki2001estimation}. 
This model can be extended to intersecting classes, where now $\mathbf{f}_n$ can have more than one nonzero entry \cite{morup2011infinite,miller2009overlappingsbm1,palla2012overlappingsbm2}.

In a \emph{Bernouli-Poisson SBM} \cite{yang2013BigClam,zhou2015bernoullipoisson, shchur2019overlapping}, the probability for a non-edge is modeled by a Poisson distribution. The idea is that the more classes $n$ and $m$ share, the higher the probability that there is an edge between them.  Hence,
\begin{align}
\label{eqn:ber_poi}    
P(n\sim m|\mathbf{f}_m,\mathbf{f}_m,\mathbf{C}) =1 - e^{-{\mathbf{f}_n}^\top  \mathbf{C} \mathbf{f}_m}, 
\end{align}
with the expected number of edges being ${\mathbf{f}_n}^\top  \mathbf{C} \mathbf{f}_m$.

In both the Bernouli and Bernouli-Poisson models, the probabilisic model of the entire graph is given by a product of the probabilities of all of the events
\smallskip
\begin{align*}
\label{eqn:prob_prod}
    & P(E|\mathbf{F},\mathbf{C}) = \\ & \sqrt{\underset{n \in [N]} {\prod}\Big(\underset{m\in \mathcal{N}(n)}{\prod}P(n \sim m)\underset{m \notin \mathcal{N}(n)}{\prod} P(\neg(n \sim m))\Big)}.
\end{align*}
Here, the square root is taken since we assume that the graph is undirected so the product goes over all of the edges twice \cite{aicher2015wsbm,morup2011infinite}. 

When fitting an SBM to a graph, both the class affiliations of nodes and the block structure $\mathbf{C}$ are learned %\ron{[*citations*]}using variational techniques 
\cite{snijders1997sbm_estimation1,nowicki2001sbm_estimation2,latouche2012variational_sbm_estimation}.


\subsection{Community Affiliation Models}

Community detection is a fundamental task in network analysis, aiming to identify groups of nodes that are more densely connected internally than with the rest of the network. This process is useful for understanding the structure and function of complex networks, such as social, biological, and information networks \cite{fortunato2016community}.

Although there exist models in which each node belongs to only one community as in traditional SBM models \cite{holland1983stochastic}, and some relatively new deep learning models such as \cite{cavallari2017learning},
it was shown that real-world networks often exhibit overlapping communities \cite{yang2014structure,yang2013BigClam}, where nodes belong to multiple groups and the probability of connectivity increases the more communities two nodes share. This indeed makes intuitive sense when looking at, e.g., social networks, where the more common interests and social circles people share the more they are likely to connect.  

An example of a community affiliation model that can generate new graphs is AGM \cite{Overlapping2012}, which classifies all of the nodes in a graph into several communities, where each community $k$ has a probability $p_k$ for two member nodes to connect. If we denote by $C_{n,m}$ the set of communities two nodes $n,m \in [N]$  share, then the probability of an edge between $n$ and $m$ is
\[
p(n\sim m) = 1 - \prod_{k \in C_{nm}} (1 - p_k).
\]
To sample a new graph from a trained model, nodes are sampled and assigned communities based on the relative sizes of communities in the training graph, and edges are connected based on their mutual community memberships.  
 
Community affiliations can also be continuous, where nodes have varying degrees of membership in multiple communities.
Community Affiliation models with continuous communities include Mixed Membership SBM (MMSBM)  \cite{airoldi2008mixed} which is an SBM which allows nodes to have mixed memberships in multiple communities, and BigClam \cite{yang2013BigClam}, which is a Bernouli Poisson model model that scales efficiently for sparse graphs. 


It is also worth mentioning the NOCD model presented in \cite{shchur2019overlapping} which uses the Bernouli Poisson loss, but embeds the nodes of the graph into the code space using a learned GNN. Another related paper is \cite{sun2019vgraph}, which models the features and communities separately, learning latent features from which it estimates the community affiliation. 

While community affiliations are typically non-negative, there are models where affiliations can be negative. An example is the Signed Stochastic Block Model (SSBM) \cite{jiang2015stochastic}.

%A well cited model for fitting continuous community affiliations to a graph is the Cluster Affiliation Model for Big Networks, or BigClam\cite{Overlapping2012} which also uses the Bernoulli-Poisson model. 

% \begin{definition}{Community Affiliation Model}

%     A community affiliation model (CAM) $M(V, C(\cdot), W(\cdot, \cdot))$ gives a \textbf{community vector} $C_u$ to any element $u\in V$, and a weight $W_{uv}$ to any $u,v \in V$ according to the \textbf{reconstruction function} $W_{uv} = W_{uv}(C_u, C_v)$ , which is defined to be continuous in the case when $Range(C)$ is continuous.

%     The range of the community affiliation function $C$ can signify subsets of  $V$ in the case that $Range(C) = \{0,1\}^d$ for some $d\in \mathbb{N}$, or a continuous affiliation scoring in the case where $Range(C)$ is continuous. 

% The elements of $C(\cdot)$ are called communities, and the community affiliation function is denoted by $F$ in the continuous case.

%     A CAM will be called non-overlapping in the discrete case when $\forall u,v \in V:  C_u \cap C_v \in \{C_u, \{\}\}$  and will be called overlapping if there are non empty intersections. 
% \end{definition}

% A community affiliation model is useful in describing a graph structure using an embedding of the nodes into some metric space (usually $(\mathbb{R}^d, l_2)$) as it replaces the graph's adjacency $E$ with the weight function $W$. 

% In the scope of this work we will refer to continuous ranged overlapping communities which we will denote as $F$. Also, we will treat $W$ as a conditional probability, namely $W(F_u, F_v) = P(E_{uv}|F_u, F_v)$.




\subsection{Anomaly Detection}
\label{extended_sec:anomaly_detection}
%nodes the behave in an inconsistent way with the rest of the graph.
Anomaly detection in graphs aims to identify nodes or subgraphs that deviate significantly from the typical patterns within the graph. This is useful in various applications such as network security, fraud detection, and social network analysis. In the unsupervised formulation of the problem, there is no labeled data in the training process, We highlight five works in this direction. GAAN \cite{chen2020gaan} and AEGIS \cite{ding2021aegis} use a generative adversarial approach, training a discriminator to distinguish real and fake nodes. Dominant \cite{ding2019dominant} and AnomalyDAE \cite{fan2020anomalydae} identify anomalies via reconstruction errors of a graph autoencoder. 

Another work is \cite{qiao2024ggad}, which considers a semi-supervised setting. Here, there is a relatively small number of available labeled normal nodes during training (nodes that are known to not be anomalies), and the goal is to predict the anomalies in the unknown nodes.

%All of the aforementioned models require the use of node attributes to recognize anomalies. One advantage of our models is that we do not require the graph to have node attribute to achieve competitive results. 
\begingroup
\color{blue}
\subsection{Link Prediction}
\label{sec:link prediction theory}
Link prediction in graphs aims to predict missing or future connections between nodes based on the existing structure and node features. This task is critical in various domains such as social networks, biological networks, and recommendation systems. In the supervised setting, a subset of dyads is withheld as a test set, with their labeles (edges or non-edges) omitted. The goal is to predict the connectivity of these omitted dyads using the remaining data, which may also include node features.

We highlight five notable works in this area. The Adamic-Adar index \cite{adamic2003friends} henceforth referred to as AA, is a second-order heuristic that assumes shared neighbors with lower degrees contribute more to the likelihood of a link, while high-degree neighbors are less significant. VGAE \cite{kipf2016variational} is a variational graph autoencoder, discussed in Section \ref{sec:deep graph autoencoders}. LINKX \cite{lim2021linkx} decouples the feature and structure representations in a graph and uses a simple framework to process this information separately. GAT \cite{GAT} applies an attention mechanism in the message passing process, allowing the model to weight neighboring nodes based on their importance. Another recent work, DisenLink \cite{zhou2022disenlink}, learns disentangled representations of nodes by capturing multiple latent factors, allowing it to model various aspects of node connectivity for link prediction. All models mentioned utilize node features accept for AA.
\endgroup

\subsection{Normalizing Flows}
\label{extended_sec:normflows}

Normalizing flows are deep learning algorithms that estimate  probability distributions, and allow an efficient sampling from this distribution. They do so by constructing an invertible coordinate transformation between the unknown target probability space and a standard probability space with a well known distribution (e.g. Gaussian). This transformation is modeled as a deep neural network, composed of a series of basic transformations called \emph{flows}. 

%\ron{More details. You are using notations that you did not explain and define. Start by giving notations to the two spaces, and to the learned transformation between them.}
\paragraph{Notations.}
Let $\mathcal{F}$ represent the space of the target data, with an unknown probability density function $p_\mathcal{F}:\mathcal{F} \rightarrow [0,\infty)$. We denote the elements of $\mathcal{F}$ by $\mathbf{f}$.

Let $\mathcal{Z}$ represent a latent space with a known probability density funcion $p_\mathcal{Z}:\mathcal{Z} \rightarrow [0,\infty)$. We denote the elements of $\mathcal{Z}$ by $\mathbf{z}$. The density function $p_\mathcal{Z}$ is often chosen to be a standard isotropic Gaussian with zero mean and identity covariance, denoted by $\mathcal{N}(0,\mathbf{I})$. 
%

\paragraph{Goal.} The goal in normalizing flows is to learn an invertible transformation $T_\theta: \mathcal{F} \rightarrow  \mathcal{Z}$ (parameterized by $\theta$) which maps the target space $\mathcal{F}$ to the latent space $\mathcal{Z}$, and preserves probabilities. Namely, $T_{\theta}$ should satisfy: for every measurable subset $F\subset \mathcal{F}$ we have
\[\int_F T_{\theta}(\mathbf{f}) p_{\mathcal{F}}(\mathbf{f})d\mathbf{f} = \int_{T_{\theta}(F)}  p_{\mathcal{Z}}(\mathbf{z})d\mathbf{z}.\]
Here $T_{\theta}(F)= \{z\in\mathcal{Z}\ |\ \exists \mathbf{f}\in F\ {\text{such that}}\ T_{\theta}(\mathbf{f})=\mathbf{z}\}$.
Such a transformation can be seen as a change of variable.

\paragraph{Density Transformation.}
Given an invertible transformation $T_\theta$, the target density $p_\mathcal{F}(\mathbf{f})$ can be expressed in terms of the latent density $p_\mathcal{Z}(\mathbf{z})$. By upholding the constraint that either probability density function has integral 1 over their respective spaces, one can deduce the change of variable formula
\[p_\mathcal{F}(\mathbf{f}) = p_\mathcal{Z}(T_\theta(\mathbf{f})) \bigg| \det \bigg( \frac{\partial T_\theta (\mathbf{f})}{\partial  \mathbf{f}}  \bigg) \bigg|. \]
Here, $\frac{\partial T_\theta (\mathbf{f})}{\partial  \mathbf{f}}$ denotes the Jacobian matrix of the transformation $T_\theta$ with respect to $\mathbf{F}$ and $\det(\cdot)$ denotes its determinant.

\paragraph{Sequential Composition of Flows.}

In practice, the transformation $T_\theta$ is modeled as a composition of $L\in \mathbb{N}$ simpler invertible transformations $\{T_{\theta^i}\}_{i=1}^L$ between the consecutive spaces $\{\mathcal{Z}^i\}$ where $\mathcal{Z}^L = \mathcal{F} $ and $\mathcal{Z}^1 = \mathcal{Z}$ in the previous notations. The invertible mappings $T_{\theta^i}: \mathcal{Z}^i \rightarrow \mathcal{Z}^{i-1}$ are called flows, and we have 
\[T_\theta = T_{\theta^L} \circ T_{\theta^{L-1}} \circ \dots \circ T_{\theta^1}\]
where $\circ$ denotes function composition. The overall density function for $\mathbf{f}$ then becomes
\[p_\mathcal{F}(\mathbf{f}) = p_\mathcal{Z}(T_\theta(\mathbf{f}))\prod_{i=1}^{L} \bigg| \det \bigg( \frac{\partial T_\theta^i (\mathbf{f}_i)}{\partial  \mathbf{f}_i}  \bigg) \bigg|. \]
Here, $\mathbf{f}_i$ is the intermediate representation after applying the first $i-1$ flows.  
%
The algorithm is optimized by maximum log likelihood, namely, by maximizing
\begin{align*}
    & \log\big(p_\mathcal{F}(\mathbf{f})\big) \\
    & = \log\big(p_\mathcal{Z}(T_\theta(\mathbf{f}))\big) + \sum_{i=1}^{L} \log\Bigg(\bigg| \det \bigg( \frac{\partial T_\theta^i (\mathbf{f}_i)}{\partial  \mathbf{f}_i}  \bigg) \bigg|\Bigg), 
\end{align*}
using gradient descent.

Since $T_\theta$ and every $T_\theta^i$ are invertible, information can also flow in the opposite direction in order to sample data. First, a point $\mathbf{z}$ is sampled $\mathbf{z} \sim p_\mathcal{Z}$ and the inverse transformation $T_\theta^{-1}$ is applied to map the sample $\mathbf{z}$ into the data space $\mathcal{F}$. Since $T$ is composed of a series of flows, the generated point in the data space is 
\[\mathbf{f} = T_\theta^{-1}(\mathbf{z}) = T_{\theta^1}^{-1} \circ T_{\theta^2}^{-1} \circ \dots \circ T_{\theta^L}^{-1}(\mathbf{z}).\]
%The result $\mathbf{f}$ is a data point in the original data space, which has been generated to follow the distribution modeled by the normalizing flow.

\paragraph{RealNVP.}
One popular method of implementing normalizing flow is the \emph{Real Valued Non-Volume Preserving (RealNVP)}. In this model, the flows $T_{\theta^i}$ are modeled as mappings called \emph{coupling blocks}. In a coupling block, the input  vector $\mathbf{f}^i$ (or $\mathbf{z}^i$, in the generative direction) is split into two vectors: $\mathbf{f}^i_A$ and $\mathbf{f}^i_B$, each the same dimension. Namely, $\mathbf{f}^i=(\mathbf{f}^i_A,\mathbf{f}^i_B)$.
The flow $T_{\theta^i}$ is then defined to be
\begin{align}
    \mathbf{f}_A^{i-1}& = \mathbf{f}_A^i \\
    \mathbf{f}_B^{i-1} &= \mathbf{f}_B^i \odot \mathbf{s}_{\theta^i}(\mathbf{f}_A^i) + \mathbf{t}_{\theta^i}(\mathbf{f}_A^i) 
\end{align}
Where $(\mathbf{s}_{\theta^i}(\cdot), \mathbf{t}_{\theta^i}(\cdot))$ are Multi Layer Perceptrons (MLPs) parameterized by $\theta^i$, and $\odot$ is elementwise product. 
In addition, the elements of $\mathbf{f}^i$ are permuted at every $i$ before applying the coupling block, using predefined permutations, that are parts of the hyperparameters of the model.
It is easy to see that the Jacobian $\bigg( \frac{\partial T_\theta^i (\mathbf{f}_i)}{\partial  \mathbf{f}_i}  \bigg)$ is a diagonal matrix with $1$ for the first $\dim(\mathbf{f}_A^i)$ elements, and $s_{\theta^i}(\mathbf{f}_A^i)$ for the last $\dim(\mathbf{f}_B^i)$ elements. This  makes the log of the determinant of the Jacobian
\begin{equation}
\log\Bigg(\bigg| \det \bigg( \frac{\partial T_\theta^i (\mathbf{f}_i)}{\partial  \mathbf{f}_i}  \bigg) \bigg|\Bigg) = \sum_{j=0}^{\dim(\mathbf{f}_B^i)}\log \Big( \mathbf{s}_{\theta^i}(\mathbf{f}_A^i)_j \Big). 
 \label{eqn_ext:log_det_jac}   
\end{equation}
Hence, this determinant is easily computed in practice.

For generation, the inverse transformation can be calculated easily as
\begin{align*}
\mathbf{z}_A^{i+1} &= \mathbf{z}_A^{i} \\
\mathbf{z}_B^{i+1} &= \frac{\mathbf{z}_B^{i} - \mathbf{t}(\mathbf{z}_A^i)}{\mathbf{s}(\mathbf{z}_A^i)}
\end{align*}
where the division is elementwise.
Here, the log determinant of the Jacobian can be calculated similarly to (\ref{eqn_ext:log_det_jac}),  applied to $1/\mathbf{s}(\mathbf{z}_A^i)$.

\subsection{The Lorentz Inner Product}
\label{The Lorentz Inner Product}

The Lorentz inner product is a bilinear form used in the context of special relativity to describe the spacetime structure. In a four-dimensional spacetime, the Lorentz inner product between two vectors \(\mathbf{v}\) and \(\mathbf{w}\) is given by:

\[
\langle \mathbf{v}, \mathbf{w} \rangle = v_0w_0 - v_1w_1 - v_2w_2 - v_3w_3
\]
Here, \(v_0\) and \(w_0\) represent the "time" components, while \(v_1, v_2, v_3\) and \(w_1, w_2, w_3\) represent the "spatial" components. The negative sign in front of the time component is what distinguishes the Lorentz inner product from the standard Euclidean inner product, making it suitable for modeling the geometry of spacetime where time and space are treated differently. Note that due to the subtraction, the Lorenz inner product is in fact not an inner product as it is not positive definite. In the context of special relativity, the points for which the Lorenz product remain positive define the so called "light cone" structure of spacetime, separating events into those that are causally connected and those that are not. 

The Lorentz inner product is a specific example of a broader class of inner products known as \textbf{pseudo-Euclidean inner products}. In a pseudo-Euclidean space, the inner product can have a mixture of positive and negative signs, leading to different geometric properties. These spaces generalize the concept of Euclidean space by allowing for non-positive definite metrics.
%For example, in a space with signature \((p, q)\), where \(p\) dimensions have positive sign and \(q\) dimensions have negative sign, the inner product takes the form
%\[
%\langle \mathbf{v}, \mathbf{w} \rangle = \sum_{i=1}^{p} v_i w_i - \sum_{j=1}^{q} v_j w_j.
%\]
%This form extends the concept of the Lorentz inner product to higher-dimensional pseudo-Euclidean spaces, where the structure of the space can vary depending on the relative number of positive and negative dimensions. These products, as the Lorenz "inner product" do not satisfy the positivity condition of an inner product.

\section{PClam and PieClam as Graphons}


A graphon is a model which can be seen as a graph generative model that extends SBMs.
A graphon \cite{graphon1,cut-homo3} can be seen as a weighted graph with a ``continuous'' node set $[0,1]$. % or more accurately, the nodes are parameterized by an atom-less standard probability space called the \emph{graphon domain}. Since all such graphon domains are equivalent to $[0,1]$ with the standard Lebesgue measure (up to a measure preserving bijection), we take $[0,1]$ as the node set. %Hence, a graphon is defined as a measurable function $W:[0,1]^2\rightarrow\mathbb{R}$. 
%




\begin{definition}
    The space of graphons $\mathcal{W}$ is defined to be the set of all measurable function  $W:[0,1]^2\rightarrow [0,1]$ which are symmetric, namely $W(x,y)=W(y,x)$. 
\end{definition}
%
%
The edge weight $W(x,y)$ of a graphon $W\in \mathcal{W}$ can be seen as the probability of having an edge between node $x$ and node $y$.   Given a graphon $W$, a random graph is generated by sampling  independent uniform nodes $\{X_n\}$ from the graphon domain $[0,1]$, and connecting each pair $X_n,X_m$ in probability $W(X_n,X_m)$ to obtain the edges of the graph.

We next show that Clam models with prior are special cases of graphon models. 
Note that the prior $p$ defines a standard atomless probability spaces over the community affiliation space $\mathbb{R}^K$. Since all standard atomeless probability spaces are equivalent, there is a probability preserving a.e. bijection  $\xi_p:[0,1]\rightarrow\mathbb{R}^K$ that maps the prior probability to the uniform probability over $[0,1]$. Now, the PieClam model for generating a graph of $N$ nodes can be written as follows.
\begin{itemize}
    \item 
    Sample $N$ points $\{X_n\}_{n=1}^N\subset[0,1]$ uniformly. Observe that $\{\xi_p(X_n)\}_{n=1}^N\subset \mathbb{R}^K$ are indepndent samples via the probability density $p$.
    \item 
    Connect the points according to the BigClam of IeClam models $P(n\sim m| \xi_p(X_n),\xi_p(X_m))$ (\ref{eq:ClamEdge}) or (\ref{eq:LorFeature}).
\end{itemize}
This shows that PClam and PieClam coincide with the generative graphon model $W(x,y)=P(n\sim m| \xi_p(X),\xi_p(Y))$ where $P$ is defined either by  (\ref{eq:ClamEdge}) or by (\ref{eq:LorFeature}).


\section{Extended Details on Experiments}
\label{Extended Details on Experiments}

All of the experiments were run on Nvidia GeForce RTX 4090 and Nvidia L40 GPUs.
Our code can be found in: 

\url{https://anonymous.4open.science/r/PieClam-49C4} 


\begin{table*}[t]
%\footnotesize
\centering
\begin{tabular}{|c|c|c|c|}
    \hline
  Method  & Reddit& Elliptic& Photo\\
    \hline

 (S)- IeClam& \textbf{0.639}& 0.440&\textbf{0.592}\\
 (S) - PieClam& 0.6320 ± 0.0054&  
0.4354 ± 0.0005&*0.5727 ± 0.0151\\
 (P) - PieClam& 0.5089 ± 0.0219& \textbf{0.6201 ± 0.0176}&0.4252 ± 0.0074\\
 (PS) - PieClam& *0.6329 ± 0.0052& \underline{0.5294 ± 0.0089}&0.5289 ± 0.0127\\
 (S) - BigClam & \underline{0.637}& 0.434&\underline{0.581}\\
 \hline
 DOMINANT  & 0.511& 0.296&0.514\\
 AnomalyDAE &0.509& *0.496&0.507\\
 OCGNN  &0.525 & 0.258&0.531\\
 AEGIS &0.535 & 0.455&0.552\\
 GAAN &0.522 &0.259 &0.430\\
 TAM &0.606 & 0.404 & 0.568\\
 \hline
\end{tabular}
\caption{Comparison of Clam anomaly detectors with competing methods. % to DOMINANT \cite{ding2019dominant}, AnomalyDAE \cite{fan2020anomalydae}, OCGNN \cite{wang2021ocgnn}, AEGIS \cite{ding2021aegis}, GAAN \cite{chen2020gaan} and TAM \cite{qiao2024tam}. 
 First place in \textbf{boldface}, second with \underline{underline}, third with *star. We observe that our methods are first place on all datasets. Moreover, S- IeClam, PS-PieClam and S BigClam each beats the competing methods in two out of the three datasets . The accuracy metric is areas under curve (AUC).}
\label{Table2}
\end{table*}



\begin{table*}[t]
%\footnotesize
\centering
\begin{tabular}{|c|c|c|c|}
    \hline
  Method  & Reddit & Elliptic & Photo \\
    \hline
  (S) - PieClam & 0.6320 ± 0.0054 & 0.4354 ± 0.0005 & \textbf{0.5727 ± 0.0151} \\
  (P) - PieClam & 0.5089 ± 0.0219 & \textbf{0.6201 ± 0.0176} & 0.4252 ± 0.0074 \\
  (PS) - PieClam & *0.6329 ± 0.0052 & \underline{0.5294 ± 0.0089} & 0.5289 ± 0.0127 \\
    \hline
  (S) - PieClam (No Attr) & \underline{0.6346 ± 0.0015} & 0.4349 ± 0.0009 & *0.5696 ± 0.0065 \\
  (P) - PieClam (No Attr) & 0.4530 ± 0.0142 & *0.4826 ± 0.0134 & 0.4977 ± 0.0747 \\
  (PS) - PieClam (No Attr) & \textbf{0.6351 ± 0.0015} & 0.4349 ± 0.0009 & \underline{0.5697 ± 0.0065} \\
    \hline
\end{tabular}
\caption{Comparison of attributed and unattributed PieClam anomaly detection.}
\label{Table22}
\end{table*}



\subsection{Additional Architecture Details in PieClam and PClam}

%\ron{\large Say exactly in which experiments each thing is used.}
%\textcolor{cyan}{fvdv}

%Say the important stuff in practice, e.g., the added affiliation noise when training to regularize the prior, densification, choice of the number of communities, etc.
\paragraph{Added Affiliation Noise.}
When training the prior, overfitting may cause the probability to spike around certain areas in the affiliation space, e.g., around the affiliation features of the nodes. To avoid this issue, we add gaussian noise to the affiliation vector as a regularization at each step of training the normalizing flow prior model. The normalizing flow model then transforms the same point to a slightly different location in the code space. The noise optimization therefore provides a resolution to the prior, smoothing the distribution in the affiliation space.  Noise addition is the primary regularization method we have used when training the prior in all of our experiments.

% \ron{in the *** experiments [list the experiments where you used noise, and write the amount fo noise for each experiment]}. 

%\ron{What is the amount of noise you add? what variance in which experiment?}


\paragraph{Densification.}
\label{sec: densification definition}
For very sparse datasets, Clam models may not behave well. In such cases, the community structure is sometimes unstable, where the same node can find itself in different communities based on slightly different initial conditions. In order to strengthen the connections within communities, we apply a two-hop densification scheme on very sparse graphs. Namely, we connect two disjoint nodes $n,m$ with an edge if there is a third node $k$ for which $n \sim k$ and $m\sim k$. We find that this scheme improves anomaly detection on Elliptic, Photo and Reddit datasets. 
Densification may strengthen the community structure in some datasets, but can also destroy it in others. In the experiments on synthetic datasets we did not use densification, as these graphs are relatively dense.
\textcolor{blue}{after a closer examination, the elliptic dataset is not improved by densification. this is probably because of the heterophilic properties of the dataset. Instead of going into detail on heterophilic properties, i suggest we look huristically on the datasets and say something like "reddit and photo are social networks so we expect them to have an inclusive structure whereas elliptic for which the nodes represent licit and illicit bitcoin transactions and edges represent flow of money, we would expect illicit transactions to cover their tracks and sift the money into licit transactions therefore having more of an exclusive structure". or we can leave the results as are.... 
}
%as it has made those graphs fully connected. Another example of graphs that will lose their structure because of densification is a bipartite graph. A bipartite graph consists of two clusters for which the probability of an edge between two nodes in different clusters is higher than the probability of connection within each of the clusters. Therefore, if two-hop edges are added to a bipartite graph it might create edges within the cluster between that are both connected to a node in the other cluster. \textcolor{cyan}{should we say we haven't found a specific condition? something about communities and density of edges between them.... maybe make a hypothesis?} 

%LINK TABLE!!!!
\begin{table*}[t]
    \centering
    \begin{tabular}{|c|c|c|c|c|}
        \hline
        Method  & Squirrel        & Photo           & Texas        & John's Hopkins 55 \\
        \hline
        PieClam & \textbf{98.8 ±0.0} & \textbf{98.4 ±0.0}     &   \textbf{86.4 ±0.4}           & 95.6 ±0.0*     \\
        \hline
        BigClam & \underline{98.5 ± 0.0}& 97.4 ± 0.0*  & 78.8 ± 0.5*& 94.9 ± 0.0 \\
        VGAE    &   98.2 ±0.1   &94.9 ±0.8  & 68.6 ±4.2 & 92.8 ±0.2 \\
        GAT     &   98.0 ±0.0   &97.3 ±0.3  & 68.5 ±5.4 &   94.3 ±0.5\\
        LINKX   &   98.1 ±0.3   &97.0 ±0.2  & 75.8 ±4.7 &   93.4 ±0.3 \\
        AA      &   97.1 ±0.4 &97.4 ±0.4  & 53.1 ±6.2 &  \underline{96.1 ±0.5}  \\
        DisenLink & 98.3 ±0.1* & \underline{97.9 ±0.1} & \underline{80.7 ±4.0} &   \textbf{97.5 ±0.1} \\
        \hline
    \end{tabular}
    \caption{Comparison of Clam link predictors with competing models. First place in \textbf{boldface}, second with \underline{underline}, third with *star}
    \label{table: link table}
\end{table*}
% LINK TABLE!!!!!


\subsection{Anomaly detection}


\paragraph{Metric.} In order to measure the accuracy of the anomaly detection classification we use the area under the Receiver Operating Characteristic (ROC) curve \cite{peterson1954roc}. The curve is a plot of the true positive rate (TPR) against the false positive rate (FPR) for the range of the threshold values between the value that classifies all samples as true and the value that classifies all as false. The area under the curve signifies the general tendency of the curve toward the point (0,1) for which FPR$= 0$ and TPR$= 1$.


\paragraph{Additional experiment.} We add additional experiments on anomaly detection using PieClam, where we averaged the accuracy over 10 rounds and included error bars.

In this setup we compose three iteration steps of the form $\mathbf{F}$-500 $\rightarrow$ $p$-1300 with initial learning rates of $3e-6$ for the features and $2e-6$ for the prior, and initial noise amplitude of $0.05$. The learning rates and noise amplitude were halved after every $\mathbf{F} \rightarrow p$ iteration.  The results are in Table~\ref{Table2}.

Furthermore, in this section we also run our anomaly detection methods without using the node features, namely, only using the graph structure. We find that we comparable results on Photo and Reddit with or without node features, but that the detection on Elliptic decreases significantly. Still, anomaly detection on Elliptic without node features is competitive with state of the art models that do use node features.  The results are in Table~\ref{Table22}.

%In the with-or-without attributes experiment, we show the results of the study described above \ron{Which one? You discussed many things} with $95\%$ confidence intervals . We also analyze the significance of the effect of adding attributes using a Student t-test. \ron{Have you done it? If not, delete.}



%\paragraph{Scheduler.}
%In the additional anomaly detection experiments we halved both the feature and prior learning rates after every $\mathbf{F} \rightarrow p$ cycle, a method known as scheduling \cite{lecun1998scheduler}. Scheduling was not used in the synthetic experiments.  


We proceed to compare the optimization that includes the node features to the optimization that didn't. The comparison is presented in 






\paragraph{T-Test Comparison Summary.}
We conducted t-tests to determine if node features/attributes ("With Features") significantly improve performance compared to not using attributes ("Without Features"), across the Reddit, Elliptic, and Photo datasets for each method (S, P, PS). Even though S method is not affected directly by the affiliation vectors, the latter can affect the convergence of the prior in the affiliation space and have an indirect effect.
\begin{itemize}
    \item \textbf{Reddit Dataset}: Attributes did not significantly improve performance for the "Star" (S) and "Prior Star" (PS) methods. However, "Prior" (P) did show a significant improvement with attributes (p-value=0.0002).
    \item \textbf{Elliptic Dataset}: "Prior" (P) and "Prior Star" (PS) both showed significant improvements with attributes (p-value$<$0.0001). However, the "Star" (S) method did not significantly benefit from attributes.
    \item \textbf{Photo Dataset}: The "Star" (S) method did not show significant difference (p-value=0.6864). The "Prior Star" (PS) method, however, did show a significant benefit from attributes (p-value$<$0.0001), and "Prior" (P) had a borderline significant result (p-value=0.0529).
\end{itemize}

In conclusion, attributes did not make a difference for the "Star" (S) method across all datasets (as would be expected), and for the "Prior Star" (PS) method in the Reddit dataset. However, the "Prior" (P) and "Prior Star" (PS) methods mostly showed improvement in the Elliptic dataset. %\textcolor{cyan}{i think that means that the prior wasn't really helping in the anomaly detection task besides the attributes use}



%\textcolor{cyan}{i'm leaving this table here in case you need it for something.... i dont know how to make the tables fit the whole page sorry it looks bad}

%\textcolor{cyan}{ morning talk. where exactly used every architectural detail, move the tables upload the figures.}



%\subsection{Reconstruction Of Prior}
%In this section we extend Section \ref{sec:reconstruction_prior}. As mentioned,
%we consider ground-truth synthetic priors $p:\mathbb{R}^C\rightarrow[0,\infty)$ on the affiliation space, and sample graphs from the corresponidng Clam models with prior. 
%We sample $N=500$ points from every prior, decode the corresponding PClam or PieClam graph in the appropriate space, and sample a simple graph from the random Bernoulli edges. Then, given these sampled graph, we fit to it a PClam or PieClam model accordingly.
%The priors. The results are shown in Figure~\ron{**}.
\begingroup
\color{blue}
\subsection{Link Prediction}
\label{sec: link prediction experiment}
\paragraph{Probabilistic Model with Omitted Dyads.} 
When doing link prediction (as explained in length in Section \ref{sec:link prediction theory}) we remove a percentage of all dyads from computation, meaning that the omitted dyads are neither edges nor non-edges. Having non-affiliated dyads will manifest in the probabilistic model.

For the revised probabilistic model we define the following sets containing unordered pairs of edges. Remembering that $E$ is the set of edges, we define $\bar{E}$ as the set of non-edges. We also denote the sets of omitted edges and non-edges by $\check{E}$ and $\check{\bar{E}}$. Finally, we denote the set of all of the omitted dyads by $\check{\mathcal{E}} = \check{E}\dot\bigcup \check{\bar E}$. The Clam loss neglecting the omitted dyads $\hat{l}$ will now be
\begin{equation}
\footnotesize    
\begin{split}
&2\hat{l}(\mathbf{F}, E, \check{\mathcal{E})} = \\
&\underset{(n,m)\in E\setminus\check{E}}{\sum} \log(1-e^{-\mathbf{f}_n^\top \mathbf{L} \mathbf{f}_m}) - \underset{(n,m)\in \bar{E}\setminus\check{\bar{E}}}{\sum} \mathbf{f}_n^\top \mathbf{L} \mathbf{f}_m
\end{split}
\end{equation}

Which is equivalent to:

\begin{equation}
\footnotesize
\begin{split}
& \underset{(n,m)\in E\setminus\check{E}}{\sum} \log(e^{\mathbf{f}_n^\top \mathbf{L} \mathbf{f}_m}-1) - \underset{(n,m)\in \mathcal{E}}{\sum} \mathbf{f}_n^\top \mathbf{L} \mathbf{f}_m + \sum_{n \in [N]} \mathbf{f}_n^\top \mathbf{L} \mathbf{f}_m\\ 
&+ \underset{(n,m) \in \check{\mathcal{E}}}{\sum}\mathbf{f}_n^\top \mathbf{L} \mathbf{f}_m.
\end{split}
\label{def: link loss}
\end{equation}
The first three terms in \ref{def: link loss} are the original Clam loss as described in equations \ref{eqn:comp_loss} and \ref{eqn:loss_reparam_ie}, on the retained edge set $E\setminus \check{E}$. The second term is a sum of the feature product for all of the omitted dyads. The number of operations is now $O(\max \big\{|E|, |\check{E}|, |\check{\bar{E}}| \big\})$, and the Clam optimization is modeled as a heterogeneous MPNN. For the priored models, the prior is added to the loss in the same way as described in Equation \ref{eqn:loss_reparam_pie}.



\paragraph{Experimental Setup}
\label{sec: link experimental setup}
In the experiments, we evaluate the model’s ability to classify edges and non-edges by constructing a ROC curve as described in Section \ref{extended_sec:anomaly_detection}.

Note that, since the prior is not used in classification, there is no added benefit from utilizing node features when training the prior, as proposed in \ref{sec:pieclam for graphs with node features}. Therefore, the link prediction algorithm relies solely on the graph’s topological structure and the prior functions only as regularization.

We use the scheme from \cite{zhou2022disenlink}, where $10\%$ of the edges are omitted as a test set. For each omitted edge, five times that number of non-edges are also omitted. Additionally, for each run of the optimization another $5\%$ of the remaining training edges are picked at random and set aside as a validation set, which is used to optimize hyperparameters for each dataset separately. Each experiment is repeated 10 times, and we calculate the average and standard deviation to report the results along with confidence bounds.

\paragraph{Datasets and Baselines.}
We conduct these experiments on four datasets: Facebook100’s Johns Hopkins 55 \cite{lim2021hopkins}, Amazon Photo \cite{shchur2018photo}, WebKB’s Texas and Wikipedia’s Squirrel \cite{pei2020texasquirrel}. We compare ourselves to the baseline methods AA \cite{adamic2003friends}, VGAE \cite{kipf2016variational}, GAT \cite{GAT}, LINKX \cite{lim2021linkx}, DisenLink \cite{zhou2022disenlink} and BigClam.


\paragraph{Hyperparameter Configuration.}
As was mentioned in \ref{sec: link experimental setup}, hyperparameters are selected separately for each dataset and model by randomly selecting a validation set outside of the test set. Since all four Clam models are a special case of PieClam, we present only the best performing model and treat it as a hyperparameter. We describe the hyperparameter configuration for each of the models in Table \ref{table: link hypers table}. The columns of the table signify the Clam model, AS dimension, number of feature iterations, feature step size, number of prior weight iterations, the prior weight step size, the added noise amplitude, $\mathbf{s}$ community regularization and the first function in the alternation. The following hyperparameters are constant across all datasets: feature step size (5e-6), l$_1$ regularization ($1$), Noise amplitude ($0.1$) the number of alternations ($7$), scheduler step size ($3$) and the scheduler step size ($0.5$).

\begin{table*}[t]
    \centering
    \begin{tabular}{|c|c|c|c|c|c|c|c|c|}
        \hline
        &Model  & Dim. & Feat. Iter.  & Prior Iter. & Prior Step & S Reg. & First Optimization\\
        \hline
        Squirrel &  PieClam   &     90 &   700   & 2000 & 5e-6     & 0.1 & Feature   \\
        Photo    &  PClam     & 150    &  700    & 1300  & 3e-6  & - & Feature \\
        Texas    &   PieClam  &  22    &  2000   &   1500 & 5e-6 &  0.0& Prior       \\
        JH55     &   PClam    &   200 & 750     &  1800   & 3e-6   & - &Feature\\
        
        \hline
    \end{tabular}
    \caption{Link prediction hyperparameters for each of the datasets}
    \label{table: hypers table}
\end{table*}




\paragraph{Results}
The results on the datasets and the comparison to the other baseline models is presented in Table \ref{table: link table}. Baseline results are taken from Table $2$ in \cite{zhou2022disenlink}.

% link prediction table
% \begin{table*}[t]
%     \centering
%     \begin{tabular}{|c|c|c|c|c|}
%         \hline
%         Method  & Squirrel        & Photo           & Texas        & John's Hopkins 55 \\
%         \hline
%         PieClam & \textbf{98.8± 0.0} & \textbf{98.4±0.0}     &   \textbf{83.2±0.6}           & 95.6±0.0*     \\
%         \hline
%         BigClam & \underline{98.5 ± 0.0}& 97.4 ± 0.0*  & 78.8 ± 0.5*& 94.9 ± 0.0 \\
%         VGAE    &   98.2 ±0.1   &94.9 ±0.8  & 68.6 ±4.2 & 92.8 ±0.2 \\
%         GAT     &   98.0 ±0.0   &97.3 ±0.3  & 68.5 ±5.4 &   94.3 ±0.5\\
%         LINKX   &   98.1 ±0.3   &97.0 ±0.2  & 75.8 ±4.7 &   93.4 ±0.3 \\
%         AA      &   97.1 ±0.4 &97.4 ±0.4  & 53.1 ±6.2 &  \underline{96.1 ±0.5}  \\
%         DisenLink & 98.3 ±0.1* & \underline{97.9 ±0.1} & \underline{80.7 ±4.0} &   \textbf{97.5 ±0.1} \\
%         \hline
%     \end{tabular}
%     \caption{Comparison of Clam link predictors with competing models. First place in \textbf{boldface}, second with \underline{underline}, third with *star}
%     \label{table: link table}
% \end{table*}
\endgroup


\begingroup
\color{blue}
\subsection{Reconstructing Synthetic Priors}
\label{sec:reconstruction_prior}
We consider two ground-truth synthetic priors $p_{big}:\mathbb{R}^2_+\rightarrow[0,\infty)$ in PClam, and $p_{ie}:\mathcal{T}\rightarrow[0,\infty)$ in PieClam. %, modeled via a realNVP normalizing flow \cite{dinh2016density}. 
We sample $N=500$ points from each of the priors, decode the corresponding PClam and PieClam graphs, and sample simple graphs from the random Bernoulli edges as shown in figure \ref{fig:ground truth 2d reconstruction}. Then, given these sampled graphs, we reset their affiliation features to random vectors and fit a BigClam and a PClam to the graph sampled from $p_{big}$, and an IeClam and a PieClam to the graph that was sampled from $p_{ie}$.


\begin{figure}[t]        
\begin{minipage}{0.6\columnwidth}
        \includegraphics[width=\textwidth]{CameraReady/LaTeX/images/synthetic_2d/circ_ie_gt.png}
        \end{minipage}

        \begin{minipage}{0.6\columnwidth}
        \includegraphics[width=\textwidth]{CameraReady/LaTeX/images/synthetic_2d/moons_big_gt.png}    
        \end{minipage}  
        \caption{Right to left: ground truth synthetic prior and sampled graph. Top: graph sampled in $\mathcal{T}$ using PieClam. Bottom: graph sampled in $\mathbb{R}^2_+$ using Pclam.}
        \label{fig:ground truth 2d reconstruction}
\end{figure}

The optimization results are shown in Figures \ref{fig: unpriored models reconstruction} and \ref{fig: priored reconstructions}. Both the priored and unpriored models effectively learned the feature reconstructions in the feature spaces $\mathcal{T}$ and $\mathbb{R}_+^2$. Notably, the features in the priored models are more localized, supporting the claim that the prior acts as a regularizer.

The priors themselves are well-approximated, as illustrated in Figure \ref{fig: priored reconstructions}, with the quality of approximation improving as more nodes are sampled from the priors.


        
    % \item Priors have been reconstructed to a similar degree.%\pause
    % \item Shape also reconstructed with BigClam and IeClam.%\pause
    % \item Prior regularizes results.
    % \end{itemize}
    
\begin{figure}[t]
    \centering
    % Top row with small images
    \begin{minipage}{0.17\textwidth}
        \centering
        \includegraphics[width=\textwidth]{CameraReady/LaTeX/images/synthetic_2d/circ_ie_reconstruct.png}
    \end{minipage}
    \hspace{0.02\textwidth}
    \begin{minipage}{0.2\textwidth}
        \centering
        \includegraphics[width=\textwidth]{CameraReady/LaTeX/images/synthetic_2d/moons_big_reconstruct.png}
    \end{minipage}
    \caption{Left to Right: IeClam and BigClam Reconstructions of features sampled from synthetic priors}
    \label{fig: unpriored models reconstruction}
\end{figure}
    % Bottom row with large images


\begin{figure}[t]
\centering
    \begin{minipage}{0.3\textwidth}
        \centering
        \includegraphics[width=\textwidth]{CameraReady/LaTeX/images/synthetic_2d/piegam_circ_reconstruct.png}
    \end{minipage}
    
    \vspace{1em} % Add some vertical space between images
    \begin{minipage}{0.3\textwidth}
        \centering
        \includegraphics[width=\textwidth]{CameraReady/LaTeX/images/synthetic_2d/moons_pclam_reconstructed.png}
    \end{minipage}
    \caption{Left to right: reconstructed nodes, reconstructed prior. Top to botton: PieClam in $\mathcal{T}$, PClam in $\mathbb{R_+^2}$}
    \label{fig: priored reconstructions}
\end{figure}

\endgroup



\subsection{SBM Reconstruction And Distance Convergence}
We extend the experiment in Section \ref{sec:reconstruction_sbm} by considering another  synthetic SBM, which is off-diagonal dominant. As in Section \ref{sec:reconstruction_sbm}, we sample a simple graph with $N=210$ nodes from the SBM. The non zero probability blocks have a probability of $0.5$. We estimate the log cut distance between the Clam models (BigClam and IeClam) and the SBM during training. We consider for both methods affiliation spaces of dimensions $2,4$ and $6$. In addition, we calculate the cut distance and l2 errors between the Clam models and the SBM. The results are shown in Figures~\ref{fig:sbm_pclam1}--\ref{fig:sbm_pclam6}.



%BIGCLAM
\afterpage{
\begin{figure*}[!t]
    \begin{minipage}{0.5\columnwidth}
\includegraphics[width=\textwidth]{CameraReady/LaTeX/images/sbm_reconstruction/halfdiag.png}       \end{minipage}\hfill  
\begin{minipage}{0.5\columnwidth}
\includegraphics[width=\textwidth]{CameraReady/LaTeX/images/sbm_reconstruction/halfdiag_bigclam/2_com/adj_cropped.png}       
\end{minipage}\hfill  
\begin{minipage}{1.1\columnwidth}
\includegraphics[width=\textwidth]{CameraReady/LaTeX/images/sbm_reconstruction/halfdiag_bigclam/2_com/distances.png}
       \end{minipage}
    \caption{Left to right: Target SBM. Fitted BigClam graph with two communities. Error as a function of optimization iteration where error is, left to right,  log cut distance, cut distance, l2 distance. After convergence, the log cut distance between the SBM and BigClam is 0.0776.}
    \label{fig:sbm_pclam1}
\end{figure*}

\begin{figure*}[!t]
    \begin{minipage}{0.5\columnwidth}
\includegraphics[width=\textwidth]{CameraReady/LaTeX/images/sbm_reconstruction/halfdiag.png}       \end{minipage}\hfill  
\begin{minipage}{0.5\columnwidth}
\includegraphics[width=\textwidth]{CameraReady/LaTeX/images/sbm_reconstruction/halfdiag_bigclam/4_com/adj_cropped.png}       \end{minipage}\hfill  
\begin{minipage}{1.1\columnwidth}
\includegraphics[width=\textwidth]{CameraReady/LaTeX/images/sbm_reconstruction/halfdiag_bigclam/4_com/distances.png}
       \end{minipage}
    \caption{Left to right: Target SBM. Fitted BigClam graph with four communities. Error as a function of optimization iteration where error is, left to right,  log cut distance, cut distance, l2 distance. After convergence, the log cut distance between the SBM and BigClam is 0.0775.}
    \label{fig:sbm_pclam2}
\end{figure*}

\begin{figure*}[!t]
    \begin{minipage}{0.5\columnwidth}
\includegraphics[width=\textwidth]{CameraReady/LaTeX/images/sbm_reconstruction/halfdiag.png}       \end{minipage}\hfill  
\begin{minipage}{0.5\columnwidth}
\includegraphics[width=\textwidth]{CameraReady/LaTeX/images/sbm_reconstruction/halfdiag_bigclam/6_com/adj_cropped.png}       \end{minipage}\hfill  
\begin{minipage}{1.1\columnwidth}
\includegraphics[width=\textwidth]{CameraReady/LaTeX/images/sbm_reconstruction/halfdiag_bigclam/6_com/distances.png}
       \end{minipage}
    \caption{Left to right: Target SBM. Fitted BigClam graph with six communities. Error as a function of optimization iteration where error is, left to right,  log cut distance, cut distance, l2 distance. After convergence, the log cut distance between the SBM and BigClam is 0.0776.}
    \label{fig:sbm_pclam3}
\end{figure*}
% ====== bigclam =======

% IECLAM 

\begin{figure*}[!t]
    \begin{minipage}{0.5\columnwidth}
\includegraphics[width=\textwidth]{CameraReady/LaTeX/images/sbm_reconstruction/halfdiag.png}       \end{minipage}\hfill  
\begin{minipage}{0.5\columnwidth}
\includegraphics[width=\textwidth]{CameraReady/LaTeX/images/sbm_reconstruction/halfdiag_ieclam/2_com/adj_cropped.png}       \end{minipage}\hfill  
\begin{minipage}{1.1\columnwidth}
\includegraphics[width=\textwidth]{CameraReady/LaTeX/images/sbm_reconstruction/halfdiag_ieclam/2_com/distances_2com.png}
       \end{minipage}
    \caption{Left to right: Target SBM. Fitted IeClam graph with two communities. Error as a function of optimization iteration where error is, left to right,  log cut distance, cut distance, l2 distance. After convergence, the log cut distance between the SBM and IeClam is 0.0662.}
    \label{fig:sbm_pclam4}
\end{figure*}

\begin{figure*}[!t]
    \begin{minipage}{0.5\columnwidth}
\includegraphics[width=\textwidth]{CameraReady/LaTeX/images/sbm_reconstruction/halfdiag.png}       \end{minipage}\hfill  
\begin{minipage}{0.5\columnwidth}
\includegraphics[width=\textwidth]{CameraReady/LaTeX/images/sbm_reconstruction/halfdiag_ieclam/4_com/adj_cropped.png}       \end{minipage}\hfill  
\begin{minipage}{1.1\columnwidth}
\includegraphics[width=\textwidth]{CameraReady/LaTeX/images/sbm_reconstruction/halfdiag_ieclam/4_com/dists.png}
       \end{minipage}
    \caption{Left to right: Target SBM. Fitted IeClam graph with four communities (two inclusive and two exclusive). Error as a function of optimization iteration where error is, left to right,  log cut distance, cut distance, l2 distance. After convergence, the log cut distance between the SBM and IeClam is 0.0312.}
    \label{fig:sbm_pclam5}
\end{figure*}

\begin{figure*}[!t]
    \begin{minipage}{0.5\columnwidth}
\includegraphics[width=\textwidth]{CameraReady/LaTeX/images/sbm_reconstruction/halfdiag.png}       \end{minipage}\hfill  
\begin{minipage}{0.5\columnwidth}
\includegraphics[width=\textwidth]{CameraReady/LaTeX/images/sbm_reconstruction/halfdiag_ieclam/6_com/adj_cropped.png}       \end{minipage}\hfill  
\begin{minipage}{1.1\columnwidth}
\includegraphics[width=\textwidth]{CameraReady/LaTeX/images/sbm_reconstruction/halfdiag_ieclam/6_com/distances.png}
       \end{minipage}
    \caption{Left to right: Target SBM. Fitted IeClam graph with six communities  (three inclusive and three exclusive). Error as a function of optimization iteration where error is, left to right,  log cut distance, cut distance, l2 distance. After convergence, the log cut distance between the SBM and IeClam is 0.0309.}
    \label{fig:sbm_pclam6}
\end{figure*}
}


\begingroup
\color{blue}
\subsection{Ablation Study}
\label{sec:ablation}
% \paragraph{Learning Rates and Number of Iterations.}
% A correct ratio between the learning rates and number of iterations of the prior and affiliation features is critical to the approximation's success. On one hand, if the features do not change enough in the $f$ iteration, 

% A learning rate on the prior that will be too large could collapse all of the features to one point. The learning rate of the features is affected both by the . The prior learning rate doesn't affect this velocity. instead, a learning rate of the prior that is too big will not give the nodes enough time to separate and will exert a pull on them that will make them eventually collapse to one point.

\paragraph{l$_1$ and $\mathbf{s}$ Regularization.}
The l$_1$ regularization is adapted from the BigClam original architecture and shows improvement when optimizing the models. This regularization type adds a constant update in all of the directions, which means that for every node $n$ and it's affiliation feature $\mathbf{f}_n$, components that are not updated are moved towards 0 regardless of their magnitude. This constant push towards 0 forces sparsity in the affiliation matrix and also acts as a regularization. The $\mathbf{s}$ regularization is an l$_1$ regularization of the $\mathbf{s}$ features and has shown effective for the Squirrel dataset.


\paragraph{Noise Addition.} Due to the high expressivity of the normalizing flow model, the neural network tends to learn a narrow distribution around the affiliation features. This overfitting tendency cannot be alleviated by simply reducing the network parameters without significantly compromising the expressivity of the prior. Noise addition offers a way to control the prior’s resolution while preserving its ability to learn complex distributions. An example comparing optimization results with and without noise addition is shown in Figure \ref{fig: noise comparison}.


\begin{figure}[t]
\centering
    \begin{minipage}{0.4\textwidth}
        \centering
        \includegraphics[width=\textwidth]{CameraReady/LaTeX/images/synthetic_2d/moons_pclam_reconstructed.png}
    \end{minipage}
    
    \vspace{1em} % Add some vertical space between images
    \begin{minipage}{0.4\textwidth}
        \centering
        \includegraphics[width=\textwidth]{CameraReady/LaTeX/images/ablation/no_noise.png}
    \end{minipage}
    \caption{Left to right: reconstructed nodes, reconstructed prior. Top to botton: With and without noise addition.}
    \label{fig: noise comparison}
\end{figure}


\paragraph{Densification.}
\label{sec:ablation densification}
The densification method described in Section \ref{sec: densification definition} can yield varying results depending on the internal structure of the graph. To illustrate the impact of densification, we conducted 10 anomaly detection experiments, following the setup from Section \ref{sec:anomaly_detection}, but without applying densification. The results are shown in Table \ref{TableAblation}.
\begin{table}[H]
\centering
\small
\begin{tabular}{|c|c|c|c|}
    \hline
  Method  & Reddit & Elliptic & Photo \\
    \hline
  (S) - PieClam & \textbf{0.6320 ± 0.0054} & \textbf{0.4354 ± 0.0005} & \textbf{0.5727 ± 0.0151} \\
  (S) - PieClam (No Dens.) & 0.5480 ± 0.0022 & 0.4009 ± 0.0021& 0.5293 ± 0.0046\\
\hline
  (P) - PieClam & \underline{0.5089 ± 0.0219} & 0.6201 ± 0.0176 & \underline{0.4252 ± 0.0074} \\
  (P) - PieClam (No Dens.) & \underline{0.5088 ± 0.0383}& \textbf{0.6378 ± 0.0245}& \underline{0.4325 ± 0.0089}\\
    \hline
  
  (PS) - PieClam & \textbf{0.6329 ± 0.0052} & 0.5294 ± 0.0089 & \textbf{0.5289 ± 0.0127} \\
  (PS) - PieClam (No Dens,) & 0.5466 ± 0.0047& \textbf{0.6181 ± 0.0233}& 0.4538 ± 0.0030\\
    \hline
\end{tabular}
\caption{Comparison of PieClam on densified and non densified datasets. \textbf{Boldface} denotes a significantly better result (difference of over 2 standard deviations), \underline{underline} denotes an insignificant difference.}
\label{TableAblation}
\end{table}

As seen in Table \ref{TableAblation}, densification generally improved results across the datasets, except for the Elliptic dataset using the P and PS methods, where the undensified datasets performed significantly better. For the Reddit and Photo datasets with the P method, the difference was not significant.

% An open question for future work is identifying the conditions under which densification should be used. It is hypothesized that this decision may depend on the graph's homophily as defined in \cite{zhu2020graph_homophily}. In highly homophilic graphs, densification may preserve and strengthen community structures, whereas in bipartite or heterophilic structures, it may blur distinctions by overly connecting nodes within the same group. This is evident in Table \ref{TableAblation}, where densification improved anomaly detection on the homophilic Reddit and Photo graphs but worsened performance on the heterophilic Elliptic dataset.

\endgroup

%\subsection{Ablation Study}
%\textcolor{cyan}{should i refer only to the anomaly detection or also to the shape experiments?}
%\paragraph{Large Learning Rate for Feature Optimization}
%If the learning rate for the feature optimization is too large...

%\paragraph{Relationship Between Prior and Feature Optimization Parameters}
%Write what happens when you change the archtiecture.

%Without feature noise?

%Different nbumber of iterations?

%Optimizing the prior together with the Clam features?


\end{document}

\clearpage

\[\]

\[\]


\[\]

\[\]


\[\]


\section{Notes and Junk}

\[\]


\subsection{Proof of he non-negative overlapping regularity lemma}

In this section, we denote $\ip{w}{v}= w^\top v$.

\begin{lemma} 
\label{fact:szlemma}
Let $\mathcal{K}_1, \mathcal{K}_2,\ldots$ be arbitrary nonempty subsets of $\mathbb{R}^J$ for some integer $J$. Then, for every $\epsilon>0$ and $\mathbf{v}\in \mathbb{R}^J$ with $v_j\geq 0$ for every $j\in[J]$, there is  $m\leq \lceil 1/\epsilon^2 \rceil$  and  $\mathbf{v}^{(i)}\in \mathcal{K}_i$ and $\gamma_i \in \mathbb{R}$,  $i\in[m]$, such that
\[\sum_{i=1}^{m}\gamma_i v^{(i)}_j \geq 0\]
for every $j\in[J]$, and such that every $\mathbf{w}\in \mathcal{K}_{m+1}$ which satisfies $w_j\geq 0$ for every $j\in[J]$ 
also satisfies
\begin{align}
\label{l:hilbert}
    \bigg|\ip{\mathbf{w}}{\mathbf{v}-(\sum_{i=1}^{m}\gamma_i \mathbf{v}^{(i)})}\bigg|\leq \epsilon\ \norm{w}\norm{v}.
\end{align}
\end{lemma}


\begin{proof}
Let $\epsilon>0$. For every $k$, let
\[\eta_k = \inf_{\boldsymbol{\gamma}, \mathbf{v}}\norm{v-\sum_{i=1}^k\gamma_iv_i}\]
where we denote $\boldsymbol{\Gamma}= \{\gamma_1,\ldots,\gamma_k\} \in \mathbb{R}^k$ and 
$\mathbf{V}=\{v_1,\ldots,v_k\}\in\mathcal{K}_1\times \ldots \times \mathcal{K}_k$, and where the infimum is over $\boldsymbol{\Gamma}$, $\mathbf{V}$ such that
\[\sum_{i=1}^{m}\gamma_i v^{(i)}_j \geq 0\]
for every $j\in[J]$.
Note that $\norm{v}\geq \eta_1\geq\eta_2\geq\ldots\geq 0$. Therefore, there is an $m\leq \lceil 1/\epsilon^2\rceil$ such that $\eta_m \leq \eta_{m+1} + \epsilon^2\norm{v}^2$. Otherwise, for every $m\leq \lceil 1/\epsilon^2\rceil$ we would have $\eta_{m+1} < \eta_m - \epsilon^2\norm{v}^2$, so for $j=\lceil 1/\epsilon^2\rceil$,
\[\eta_j \leq \eta_{j-1} - \epsilon^2\norm{v}^2 < \eta_{j-2} - 2\epsilon^2\norm{v}^2 < \ldots < \norm{v} -j\epsilon^2\norm{v}^2\]
\[= \norm{v} - \lceil 1/\epsilon^2\rceil \epsilon^2\norm{v}^2 \leq 0,\]
which is a contradiction to the fact that $\eta_j\geq 0$. Hence,  there is  $m\leq \lceil 1/\epsilon^2\rceil$ and there are  $\boldsymbol{\gamma}= \{\gamma_1,\ldots,\gamma_m\} \in \mathbb{R}^m$ and 
$\mathbf{v}=\{v_1,\ldots,v_m\}\in\mathcal{K}_1\times \ldots \times \mathcal{K}_m$ with
\[\sum_{i=1}^{m}\gamma_i v^{(i)}_j \geq 0\]
for every $j\in[J]$
such that
\[\norm{v-\sum_{i=1}^m\gamma_iv_i} \leq \eta_{m+1}+\epsilon^2\norm{v}^2.\]
Denote $v^*=\sum_{i=1}^m\gamma_iv_i$ and let $w\in\mathcal{K}_{m+1}$ with $w_j\geq0$ for every $j\in[J]$. By the definition of $\eta_{m+1}$, we have for every $t>0$,
\[\norm{v-(v^*+tw)}^2 \geq \norm{v-v^*}-\epsilon^2\norm{v}^2.\]
This is equivalent to
\[\norm{w}^2t^2 - 2\ip{w}{v-v^*}t + \epsilon^2\norm{v}^2\geq 0.\]
The discriminant of this quadratic polynomial is
\[4\ip{w}{v-v^*}^2 - 4\norm{w}^2\epsilon^2\norm{v}^2\]
and it must be non-positive to satisfy the inequality, which proves the lemma.

Or, since $t\geq 0$, this may be satisfied under a less strict condition? 

The aboev can be positive also if there are roots, but all roots have to be negative. So let's see when the roots are negtive.

Roots:
\[\frac{2\ip{w}{v-v^*} \pm \sqrt{4\ip{w}{v-v^*}^2 - 4\norm{w}^2\epsilon^2\norm{v}^2}}{2\norm{w}^2}.\]

When are both roots negative?

If the discriminant is positive, and $\ip{w}{v-v^*}$ is negative, then both roots must be negative.

%We know that $w$ is nonegative. We know that $v$ is nonegative. 

%Both roots are negative if $\ip{w}{v-v^*}<0$ and
%\[4\ip{w}{v-v^*}^2 - 4\norm{w}^2\epsilon^2\norm{v}^2< 4\ip{w}{v-v^*}^2.\]
%or
%\[\]

So actually the lemma says that either
\[\ip{w}{v-v^*}^2 \leq  \norm{w}^2\epsilon^2\norm{v}^2\]
or
\[\ip{w}{v-v^*}  < 0\]
In this second case it also needs to satisfy
\[\ip{w}{v-v^*}^2 \geq \norm{w}^2\epsilon^2\norm{v}^2\]
or
\[-\ip{w}{v-v^*} \geq \norm{w}\epsilon\norm{v}\]

This means that $w$ captures more places where $v^*$ is smaller than $v$ than larger. 

By the construction, this is true for every $tw$ such that
\[v^*+tw\geq0.\]




$ $

I need to make it then finer.

Instead of taking $w>0$, I take $w$ such that
\[v^*+w\geq 0.\]
A special case is $w\geq 0$.

Now, can be have $\ip{w}{v-v^*}  < 0$?
We know that $v\geq 0$. So $w$ would have to capture places where $v^*>v$. 

%Do I then need in the analysis to have
%\[\sum \gamma_j v_j <\]





\[\]

\[\]

Should we also restrict to $v-v^*\geq 0$ nonegative?

But then $w$ will be problematic...  


\end{proof}




\[\]


\[\]

In (\ref{sec:expressive_lim_inr}) we show that there is no possibility for inner product recreation to approximate a bipartite SBM. 
In this section we develop a reconstruction technique for such graphs and which will generalize BigClam. Furthermore, in (\ref{sec:IeClam_universality}) we show that this method can recreate any SBM, and find an upper bound on the number of communities it requires to do so.

\paragraph{The expressive limit of BigClam \label{sec:expressive_lim_inr}}

\textcolor{red}{Give the bipartite example. Then say generally that off-diagonal heavy SBMs cannot be approximated by BigClam.}
\begin{definition}{(IPRM)} 

Define an inner product reconstruction model (IPRM) as a community affiliation model $M(V, F(\cdot), W(\cdot))$ \textcolor{red}{how is this definition? F is a function of vertex and adjacency W is a function F if it's good we'll write it in the intro} \textcolor{cyan}{it's defined to be a community affiliation model in which the reconstruction is a function of the inner product}  for which the reconstruction function satisfies $W(F_{u_1},F_{u_2}) = \Phi(F_{u_1}^\top F_{u_2})$ where $\Phi$ is surjective. Since $\Phi$ is surjective and continuous it is also monotonous (and ascending).
    
\end{definition}
\begin{definition}{(Consistent CAM)}
\label{def:consistent}

We say that a CAM is consistent if  nodes of the same connectivity have the same features:  $\forall u_1, u_2 \in V : (\forall v \in V | v \neq u_1, u_2 : E_{u_1,v} = E_{u_2, v} ) \implies  F_{u_1} = F_{u_2}$.
    
\end{definition}
\ref{def:consistent} should make intuitive sense as we would expect a similar group of nodes to have a similar community affiliation.
\begin{claim}
    
An IPRM \textbf{cannot} be consistent community affiliation model. 
\end{claim}
We consider a simple counter example of a weighted bipartite graph and show that the structure cannot be approximated with an IPEM while giving the same group affiliation to nodes in the same group.

Consider a bipartite undirected graph $G = (V, \Omega)$ with $N$ nodes in which there are two groups of nodes $S$ and $K$.  For $\forall k,k_1,k_2 \in K$ and $\forall s, s_!, s_2 \in S$ define $\Omega_{KK} = \Omega_{k_1, k_2}$ , $ \Omega_{SS} = \Omega_{s_1,s_2}$ and $\Omega_{t, s} = \Omega_{TS}$.  Also, Assume that $\Omega_{KS} > \Omega_{SS} = \Omega_{KK}$.  (For extra simplicity, one can imagine $|S| = |K| = 2$ and $w_{s,k} = 1$, $\Omega_{k_1, k_2} = \Omega_{s_1, s_2} = 0$).
\textcolor{cyan}{i have a better proof that doesn't assume monotonicity}
Let us consider a consistent IPEM reconstruction for this graph denoted by $M$.  Since there are only two groups of nodes and two weights there are also only two types of features (due to consistency): $F_K$ and $F_S$.  We know that $\Phi({F_S}^\top  F_S) = \Omega_{SS}$, $\Phi({F_K}^\top  F_K)=\Omega_{KS}$ and $\Phi({F_S}^\top  F_K) = \Omega_{SK}$. Therefore, because of monotonicity we have $|F_S|^2 < {F_S}^\top  F_K$  which in turn implies that $F_S < F_K$ and thus:
\[W(|F_S|^2) = \Omega_{SS} = \Omega{KK} = W(|F_K|^2) > W(|F_S|^2)\] Where the last inequality is implied by monotonicity and leads to the contradiction.



\paragraph{Equivalence to the Graphon Object} 
We can make the equivalence of the priored conditional probability to the graphon object by describing the generation process as follows:

Assume that we have a function $\Phi$ that is a diffeomorphism from a gaussian measure space $(\Omega, \mathcal{F}, \mathcal{N}(0,1))$, and denote $P(E|F_u, F_v) = \beta(F_u, F_v)= 1-e^{-F_u B F_v}$
To sample a graph we may do the following (in case the latent probabilit is uniform in $[0,1]$): 
\begin{enumerate}
    \item Sample vectors $\{x_i\}_i \sim \mathcal{N}$
    %\item Transform the features with the inverse $\Phi^{-1}$.
    \item  connect the nodes with probability: $W(x_u,x_v)= \beta(\Phi^{-1}(x_u),\Phi^{-1}(x_v))$ 
\end{enumerate}
This retraces the graphon generation process as described by [cite]

\subsection{Geometric Interpretation}
Let us assume that we have one $s$ feature and one $t$ feature. We can think of  the two features as if they were one complex feature $f_u = s + it = re^{i\theta}$. For which the inner product is $\mathcal{R}(f_uf_v) = r_ur_v\cos(\theta_1 + \theta_2)$,  

\begin{figure}[h]
\centering
\begin{tikzpicture}
    \begin{axis}[
        axis lines=middle,
        xlabel={$t$},
        ylabel={$is$},
        xmin=-1, xmax=1,
        ymin=-1, ymax=1,
        xtick=\empty, % Remove x-axis ticks
        ytick=\empty, % Remove y-axis ticks
        xticklabels={,,}, % Ensure no x-axis tick labels
        yticklabels={,,}, % Ensure no y-axis tick labels
        xlabel style={at={(ticklabel* cs:1)},anchor=north west},
        ylabel style={at={(ticklabel* cs:1)},anchor=south west}
    ]
    % Add grid
    \addplot [domain=0:10, samples=2, dashed] {x}; % 45 degrees line
    \addplot [domain=0:10, samples=2, dashed] {-x}; % -45 degrees line
    % Add vectors
    % Vector r1 at 30 degrees
    \draw[->, thick, red] (axis cs:0,0) -- (axis cs:{cos(30)},{sin(30)}) node[anchor=south west] {$\mathbf{r_1}$};
    % Vector r2 at -20 degrees
    \draw[->, thick, blue] (axis cs:0,0) -- (axis cs:{cos(-20)},{sin(-20)}) node[anchor=north east] {$\mathbf{r_2}$};

    % Add angles

    
    % Angle for r1
    \draw (axis cs:0.5,0) arc [start angle=0, end angle=30, radius=0.5] node[pos=0.5, anchor=south] {$\theta_1$};
    % Angle for r2
    \draw (axis cs:0.5,0) arc [start angle=0, end angle=-20, radius=0.5] node[pos=0.5, anchor=north] {$\theta_2$};
        \end{axis}
\end{tikzpicture}
\caption{Inclusive and Exclusive communities as Complex Numbers}
\end{figure}

We see that as long as the features stay inside the "space cone", inner products are positive and probabilities make sense. Also, see that if features have a same sign angle, they tend to have a lower inner product because of the cosine. 

\subsubsection{Lorenz Invariance Symmetry}
Notice that the norm is the Lorenz scalar and is invariant under Lorenz boosts for different velocities in $\mathbb{R}$.  This implies invariance under hyperbolic rotation:

\[
\begin{pmatrix}
x' \\
t'
\end{pmatrix}
=
\begin{pmatrix}
\cosh(\zeta) & -\sinh(\zeta) \\
-\sinh(\zeta) & \cosh(\zeta)
\end{pmatrix}
\begin{pmatrix}
x \\
t
\end{pmatrix}
\]






\paragraph{The universality of IeClam}\label{sec:IeClam_universality}

We hereby prove the universality of the suggested method.
\begin{theorem}[Universality]
\label{thm:universality}
For every $\epsilon>0$ there exist constants $K_s(\epsilon),K_t(\epsilon)$ s.t. for every graph $G$ with $N$ nodes,  there exists a PieClam model with $N$ nodes $P=(p_{ij})_{ij \in [N]\times[N]}$, with $K_s(\epsilon)$ inclusive and $K_t(\epsilon)$ exclusive communities such that $||G-P||_{{\Box}}<\epsilon$ 
\end{theorem}
\begin{proof}[Proof]
It is known that there is an SBM $W$ with $J(\epsilon) < N$ blocks for which $||G-W||_\Box < \epsilon$,  so it will sufice to prove that any SBM can be approximated with a PieClam for which  $K_s$ and $K_t$ depend only on $J$ .

If we have a number of communities that is the number of distinct edges, we can represent each edge by two communities as follows:

Let $w_{ij}$ denote the connection weight between blocks $i$ and $j$. 
\begin{lemma}
For every SBM $W$ with $J$ nodes, there exists a PieClam model $P$ with $N=J$ nodes and $K_s(J)\leq \frac{J(J+1)}{2}$, $K_t(J)\leq \frac{J(J+1)}{2}$ such that: $||W-P||_\Box = 0$ 

We will first prove the lemma for a single bipartite block:
\begin{claim}

An SBM $W$ with only one off diagonal part $w_{ij} \neq 0$ can be approximated with the following feature setup:
\begin{align}
S^i_{kJ + l} = S^j_{kJ + l}=&\sqrt{\frac{1}{2}\log(1-w_{k,l})}  \\
& T^i_{kJ + l} = \sqrt{\frac{1}{2}\log(1-w_{k,l})} \hspace{1cm} \\ & \hspace{1cm} T^j_{kJ + l} = -\sqrt{\frac{1}{2}\log(1-w_{k,l})}
\end{align}
need to somehow remove all of the nodes lJ+k as they will have the same inner product
\end{claim}

[Proof of Claim]
Calculating the inner product 
\begin{claim}
    An SBM $W$ with only one diagonal community can be approximated with the community:
    \begin{align}
S^i_{kJ + k} = S^j_{kJ + k}=&\sqrt{\frac{1}{2}\log(1-w_{k,k})}  \\
T^i_{kJ + k} = 0 &; \hspace{1cm} T^j_{kJ + k} = 0
\end{align}

we make $J$ communities like this that represent the connection of a node to itself.
\end{claim}

\begin{proof}[Proof of lemma]

Denote the edges of  $W$ as $w_{k,l}$.  and $w_{ii} =w_{jj} = 0$.
set: 

And the demands are met. 
Since there are only $\frac{J^2 + J}{2}$ possible edge probabilities the proof  of the lemma is finished.
\end{proof}
\end{lemma}
Since We found a PieClam representation of an arbitrary SBM, theorem \ref{thm:universality} is proved.

\end{proof}



\[\]

\[\]




\subsection{Correct construction starts here}


\textcolor{red}{Cite the regularity lemma, and motivate the proof. (if there is no space, we will move this to the appendix)} %\textcolor{cyan}{that is what i didn't understand - the motivation. we want to say that under some conditions our model is a global approximator? }


\begin{definition}
  The \emph{cut norm} of a matrix $\mathbf{X}\in\mathbb{R}^{N\times N}$ is defined to be
    \begin{equation}
    \label{eq:GraphCutNorm0}  \|\mathbf{X}\|_{\square}:=\frac{1}{N^2}\sup_{\mathcal{U},\mathcal{V}\subset[N]}\Big|\sum_{i\in \mathcal{U}}\sum_{j\in \mathcal{V}}x_{i,j}\Big|.
\end{equation}  
\end{definition}
The cut norm can be thought of  as the largest valued rectangle that can be produced in an adjacency matrix using a permutation of rows and columns, normalized by the number of dyads. 

\textcolor{red}{Explain that there is a probabilistic interpretation of cut norm.} \textcolor{cyan}{I can't think of one... the densest subgraph? the biggest cut?}

%\textcolor{cyan}{dan thoughts}
% For a matrix of dyadic probabilities $P(\cdot,\cdot)$ as in our case, the cut norm signifies the sets.... for a matrix of probabilities all of the entries are positive so U and V are just [N]... the interesting things happen in the cut distance.... 
%\textcolor{cyan}{dan thoughts}
The measure of accuracy we use is defined to be
%\[E(A,P):=\]
% \[\frac{1}{N^2}\sup_{\mathcal{U},\mathcal{V}\subset[N]}\Big|\log\Big(\prod_{i\in \mathcal{U}}\prod_{j\in \mathcal{V}}
%\frac{1-P(F_u,F_v)}{1-a_{u,v}}\Big)\Big|,\]
\textcolor{cyan}{what do all of these letters mean? what are the commented formulas?} \textcolor{red}{no need to read commented text.}
\[E(A,P):=\]
\[\inf_{0\leq d<1} \Bigg( d+ \frac{1}{N^2}\sup_{\mathcal{U},\mathcal{V}\subset[N]}\Big|\log\Big(\prod_{i\in \mathcal{U}}\prod_{j\in \mathcal{V}}
\frac{1-P(F_u,F_v)}{1-(1-d)a_{u,v}}\Big)\Big|\Bigg).\]
Cut-KL divergence or cut-relative entropy.

\textcolor{cyan}{if this distance measure is smaller than $\epsilon$ then necessarily $d$ is smaller than $\epsilon$, which means that 
$0<\|\log(1-P) - \log(1-(1-d)A)\|_\square < \epsilon - d$
this is the closest we can get to an expression like:
$1-A < 1-P < (1-A)e^{\epsilon}$ 
}
The second term is the cut distance $\| X-Q\|_{\square}$ between the matrix $X$ with entries $x_{u,v}=-\log(1-(1-d)a_{u,v})$ and the matrix $Q$
 with entries $q_{u,v}=-\log(1-P(F_u,F_v))$. Namely, the cut distance between the log likelihoods of non-edges. The parameter $d$ makes the $\{0,1\}$-valued adjacency weights $a_{u,v}$ amenable to modelling via exponents. Note that the probabilistic model $1-\exp(-F_uBF_v)$ cannot attain probability 1 since $F_uBF_v$ is real. Hence, the goal of $d$ is to regularize the probability of the edges of the graph $A$, where higher regularization is penalized via the additive term $d$ in the definition of $E(A,P)$.

 The distance $E$ is related to the Kullback-Leibler (KL) divergence. Namely, for each choice of $U,V\subset[N]$, 
 \[\log\Big(\prod_{i\in \mathcal{U}}\prod_{j\in \mathcal{V}}
\frac{1-P(F_u,F_v)}{1-(1-d)a_{u,v}}\Big)\]
is an un-normalized (for density functions not summing to 1) KL divergence between the non-edge probability of the graph and the non-edge probability of the IeClam graph over the edges between  $U$ and $V$, where edges are drawn ``unnormalized-uniformly'' with probability $1/N^2$. The unnormalized uniform distribution discourages the error $E$ from choosing small blocks for maximizing (**),. Note that normalized uniform distributions would lead $E$ to choose  small blocks, which do not reflect meaningful empirical estimates of the statistics  of the edges (the edge densities of small blocks would not be interpretable as expected number of edges). Hence, $E(A,P)$ is interpreted as the maximal KL divergence between $A$ and $P$ over all blocks (up to the best regularizer $d$ plus regularization panelty).
 


We do not use the above as the optimization target in practice, since finding the supremum is NP hard. Instead, we replace it by the log likelihood described in (**). %the $L_2$ norm between the log likelihoods. This is justifies by the constructive weak regularity lemma, which asserts that minimizing the Frobenius error (which is the L2 norm over the entries of the matrix) leads to a small error in cut norm.

%This analysis also explains the approximation theoretic grounds of BigClam.

Next we define a general decoder that estimates edges by operating on pairs of points in a code space.

\begin{definition}
   A pairwaise decoder over the code spaces $\mathbb{R}^M$ is a mapping $D_M:\mathbb{R}^{2M}\rightarrow[0,1]$. Given $N$ points in the code space $\mathbf{z}=\{z_n\in\mathbb{R}^M\}_{n=1}^N$, the decoded graph $G_N(\mathbf{z})$ is the weighted graph with adjacency matrix 
   \[\mathbf{D}_M(\mathbf{z})=\big(D_M(z_n,z_k)\big)_{n,k=1}^N.\]
\end{definition}

The pairwise decoder of BigClam and IeClam involves a construction with an exponential, which cannot attain probability 1.


%\textbf{We don't need this. For now I keep it, since we may use it in the proof:}

%Clam models cannot predict edges with weight 1 due to the exponential construction. Since adjacency matrixes of unweighted graphs do have weights $1$, we have to limit the weight is a weaker sense. We propose the following assumption on the densities of edge weights, rather than individual edge weights. .

%\begin{definition}
 %   Let $R:(0,1)\rightarrow \mathbb{R}_+$ such that $R_{\epsilon}$ goes to infinity as $\epsilon$ goes to zero. 
%
%An adjacency matrix $\mathbf{A}\in[0,1]^{N\times N}$ is called \emph{weight-regular} (with respect to $R$) if for every $\epsilon>0$ there exists a matrix $A_{R_{\epsilon}}\in[0,1-\exp(-R_{\epsilon})]^{N\times N}$ such that
%\[\|\mathbf{A}_{R_{\epsilon}}-\mathbf{A}\|_{\square} < \epsilon.\]
%\end{definition}

%This bound is trivially correct  for a fixed $N$. Indeed, one can replace all the edge weights above $1-\exp(-R_{\epsilon})$ by $1-\exp(-R_{\epsilon})$, and get an error
%\[\|\mathbf{A}_{R_{\epsilon}}-\mathbf{A}\|_{\square} \leq \|\mathbf{A}_{R_{\epsilon}}-\mathbf{A}\|_1 \leq \exp(-R_{\epsilon}).\]

%So just take $\exp(-R_{\epsilon})=\epsilon$, or
%\[R_{\epsilon}=-\log(\epsilon).\]
%Then, the intersecting community regularity lemma will require $K=\frac{-18\log(\epsilon)}{4\epsilon^2}$, which is good asymptotics. 

%\[\]

\[\]

%\textbf{Continuing:}

The motivation in the following definition is that we would like to uniformly choose the dimension of the code space and the decoder, such that every graph can be approximated by decoding some points in the code space. Namely, we would like one \emph{universal} decoder that works for all graphs, as opposed to choosing the dimension of the code space depending on the graph.

\begin{definition}
    A family of code spaces $\mathbb{R}^M$ and corresponding pairwise decoders $D_M:\mathbb{R}^{2M}\rightarrow[0,1]$, parameterized by $M\in\mathbb{N}$, is called  \emph{universal} if for every $\epsilon>0$ there is $M\in\mathbb{N}$ (that depends just on $\epsilon$) such that for every $N\in\mathbb{N}$ and every graph  with adjacency matrix $\mathbf{A}$ and $N$ nodes there are $N$ points in the code space $\{z_n\in\mathbb{R}^M\}_{n=1}^N$ such that 
    \[E\big(\mathbf{A},\mathbf{D}_M(\mathbf{z})\big)<\epsilon\]
\end{definition}

\textcolor{red}{After this, show that BigClam is not universal. For full bipartite graph, you need at least as many dimensions as half the number of nodes, so $M$ cannot depend just on $\epsilon$.}

\textcolor{red}{Also show that inner product and norm distance decoder is not universal.}

%The following is a slight reformulation of \ref{weakReg}. It is a direct result of Lemma 3.1 in \citep{Szemeredi_analyst}, for the special case where the graphon in induced by a graph.
%\begin{theorem}[Weak Regularity Lemma \cite{weakReg}]
%\label{Theorem:WRL}
%    For every every graph $G=([N],\mathcal{E})$ with adjacency matrix $\mathbf{A}\in\mathbb{R}^{N\times N}$ and every $\epsilon>0$  , there is a partition $\mathcal{P}=\{\mathcal{N}_1,\ldots,\mathcal{N}_k\}$ of $[N]$ into $k\leq 2^{2/\epsilon^2}$ classes, and graph with block adjacency matrix $\mathbf{B}\in\mathbb{R}^{N\times N}$, with blocks $\{\mathcal{N}_i\times \mathcal{N}_j\}_{i,j=1}^k$, such that $\| \mathbf{A}- \mathbf{B} \|_{\square} \leq \epsilon$. 
%\end{theorem}

\[\]

For IeClam, the decoder is
\[P(F_u,F_v) = 1-\exp\big({-F_u^\top BF_v}\big)\]
Denote the decoded matrix by $P$.

\[\]



\begin{theorem}
    Let $\epsilon>0$. For $K=\frac{-18\log(\epsilon)}{\epsilon^2}$, there is a community affiliation space with $K$ inclusive and $K$ exclusive communities, such that any adjacency matrix $\mathbf{A}$ can be approximated by a decoded IeClam matrix $P$ with error $E(A,P)<\epsilon$.
\end{theorem}

As a corollary, IeClam is universal.

$ $

\textbf{Proof of universality of IeClam}

$ $

Define ${\rm diag}(\mathbf{r})\in\mathbb{R}^{N\times N}$ as the diagonal matrix with diagonal elements $\mathbf{r}\in\mathbb{R}^N$.



\begin{definition}
    A (hard) intersecting community graph (ICG) with $N$ nodes and  $K$ communities is a matrix $\mathbf{C}$ of the following form. There exist $K$ column vectors $\mathbf{Q}=\big(\mathbf{q}_k\in\{0,1\}^N\big))_{k=1}^K\in \{0,1\}^{N\times K}$ and $k$ coefficients $\mathbf{r}=(r_k\in\mathbb{R})_{k=1}^K\in\mathbb{R}^K$ such that
    \[\mathbf{C}= \mathbf{Q}{\rm diag}(\mathbf{r})\mathbf{Q}^{\top }.\]
\end{definition}

The following is a special case of the weak regularity lemma from **. Up to the small modification of the adjacency matrix, allowing it to have values in $[0,R]$, it is well known as a step in the proof of the weak regularity lemma, see, .e.g  [Frieze, Fast. , Lovasz. Analyst]. For completness, we prove it in Appendix ***.

\begin{theorem}
%\label{thm:reg_sprase}
   % Let $1\leq p\leq2$. % and $D\in\sN$.
     Let $\mathbf{A}\in[0,R]^{N\times N}$ be an adjacency matrix of a graph with $N$ nodes. Let $\epsilon>0$. Denote $K=\frac{9R^2}{4\epsilon^2}$.         Then, there exists a hard ICG $\mathbf{C}$ with $K$ communities such that 
    \begin{equation}
        \label{eq:reg_sparse}
    \|\mathbf{A}-\mathbf{C}\|_{\square} \leq \epsilon.
    \end{equation}
\end{theorem}

\[\]

Let $\epsilon>0$. 
Let $\mathbf{A}$. Let $0<d\leq 1$.

Consider the matrix  $\mathbf{X}$ with entries
\[x_{u,v}=-\log(1-(1-d)a_{u,v}).\]
%First, we approximate $X$ by $X_{\epsilon}\in[0,1-\epsilon]$ such that 
%\[\|{X-X_{\epsilon}}\|_{\square} < \epsilon.\]
%This is done by truncating the values of $X$ at $1-\epsilon$ to obtain $X_{\epsilon}$, and bounding the cut norm by the L1 norm.


%\[\]

%Problem: the first approximation is good for $A$:
%We truncate $A$ to get $A_{\epsilon} \in [0,1-\epsilon]^{N\timesN}$, with error
%\[\| A-A_{\epsilon}\| < \epsilon.\]
%Then, the error is in the log space...

%\[-\log(1-a^{\epsilon}_{u,v})\]
%has values in $[0,-\log(\epsilon)]$.

%Then the ICG regularity lemma gives an approximation in the log space.

%Maybe we define the error directly via a formal triangle inequality:

%\[E(A,P):=\]
%\[\inf_{0\leq d<1} \Bigg( d+ \frac{1}{N^2}\sup_{\mathcal{U},\mathcal{V}\subset[N]}\Big|\log\Big(\prod_{i\in \mathcal{U}}\prod_{j\in \mathcal{V}}
%\frac{1-P(F_u,F_v)}{1-a_{u,v}(1-d)}\Big)\Big|\Bigg).\]

%This will make the proof work.

%\[\]


In the construction, we build regularity features $F$ and we want
\[1-\exp(-F_uBF_v^{\top })\approx (1-d)a_{u,v}.\]
So the regularity lemma should be against the matrix $\mathbf{X}$ with entries
\[x_{u,v}=-\log(1-(1-d)a_{u,v}).\]
Note that $-\log(1-(1-d)a_{u,v})$ is increasing in $(1-d)a_{u,v}$. For $(1-d)a_{u,v}=0$ the value of this function is 0, and for $(1-d)a_{u,v}=1-d$ it is $R=-\log(d)$. Choose $d=\epsilon$. For this specific choice of $d$, if we replace $d$ by $\epsilon$ in the definition of $E$ and omit the infimum, we get an upper bound of $E$.
 
%We will later bound:
%\[\| C-X \|_{\square} \leq \|C-X_{\epsilon}\|_{\square} + \|X_{\epsilon}-X\|_{\square}.\]






We now use the overlapping weak regularity lemma.


Usind the overlaping weak regularity lemma, we approximate $X$ by an ICG with 
\[K=\frac{-18\log(\epsilon)}{4\epsilon^2}\]
communities,
\[\mathbf{C}= \mathbf{Q}{\rm diag}(\mathbf{r})\mathbf{Q}^{\top },\]
such that
\[\|\mathbf{X}-\mathbf{C}\|_{\square} \leq \epsilon.\]

Let $\mathbf{r}_+={\rm ReLU}(\mathbf{r})$ and $\mathbf{r}_-={\rm ReLU}(-\mathbf{r})$. Denote
\[\mathbf{C}_+=\mathbf{Q}{\rm diag}(\sqrt{\mathbf{r}_+})\]
and
\[\mathbf{C}_-=\mathbf{Q}{\rm diag}(\sqrt{\mathbf{r}_-}).\]
Denote the rows of $\mathbf{C}_+$ by $\mathbf{S}_n$ and  the rows of $\mathbf{C}_+$ by $\mathbf{T}_n$, for $n=1,\ldots,N$. We then concatenate them to the feature $F_n$, with the first $K$ coordinates being $\mathbf{S}_n$ (the inclusive communities) and the last $K$ coordinates being $\mathbf{T}_n$ (the exclusive communities).

It is easy to see that $F_uBF_v$ is the $(u,v)$ enrty of $Q{\rm diag}(r)Q^{\top }$. This proves that
\[\|-\log(1-(1-\epsilon)\mathbf{A})+\log(1-P)\|_{\square}\]
\[=\|\mathbf{X}-\mathbf{C}\|_{\square} \leq \epsilon.\]
The claim now follows by adding $\epsilon$ to the above formula, and then replacing $\epsilon$ by $\epsilon/2$.

In more detail, our  metric is bounded as follows
\[E(A,P) \leq \]
\[\epsilon+\frac{1}{N^2}\sup_{\mathcal{U},\mathcal{V}\subset[N]}\Big|\log\Big(\prod_{i\in \mathcal{U}}\prod_{j\in \mathcal{V}}
\frac{1-P(F_u,F_v)}{1-(1-\epsilon)a_{u,v}}\Big)\Big|\]
\[ =
\epsilon+\frac{1}{N^2}\sup_{\mathcal{U},\mathcal{V}\subset[N]}\Big|\sum_{i\in \mathcal{U}}\sum_{j\in \mathcal{V}}\Big(-\log\big(1-(1-\epsilon)a_{u,v}\big) \]
\[+ \log\big(1-P(F_u,F_v)\big)\Big)\Big|\]
\[=\epsilon + \| C-X \|_{\square} = 2\epsilon. \]
We now replace $\epsilon$ by $\epsilon/2$.

\[\]

\[\]




\section{Experiments}

\subsection{Reconstructing the prior}

\subsection{Reconstructing SBMs}

\subsection{Anomaly detection}
We compare our models PieClam and IeClam with BigClam and the state of the art methods in anomaly detection. For the priored models we use three detection methods: Vanilla Star (VS), Prior (P), and Prior Star (PS). For the non priored methods we use only Vanilla Star.

The results are brought in the table below. The initials signify the detection method on the left and the model abbreviation on the right (e.g.: VS - BC means vanilla star detection on BigClam optimization).
The results are compared with the results obtained in \cite{qiao2024ggad} for the same datasets.

\paragraph{Hyperparameter Configuration}



\section{Appendix}


\subsection{Extended Details to BigClam}
\label{Extended Details to BigClam}


BigClam is a model where a simple random graph is decoded from a set of AFs $\{\mathbf{f}_n=(f_n^1,\ldots,f_n^C)\}_{n=1}^N$ in the AS. For each pair of nodes $n,m,\in[N]$, the probability of the event $\neg(n,m)$ (no edge between $n$ and $m$), given their amount of membership of the same community $c$ is defined to be
\[
%\label{eqn:poisson}
P(\neg (n,m)|f^c_n, f^c_m) = e^{-f^c_n f^c_m}.
\]
BigClam makes two assumptions of independence. First is that the membership in one community does not affect the membership in another community. Hence,
\[P(\neg (n,m))|\mathbf{f}_n, \mathbf{f}_m)=e^{-\mathbf{f}_n^{\top }\mathbf{f}_m},\]
\[P((n,m))|\mathbf{f}_n, \mathbf{f}_m)=1-e^{-\mathbf{f}_n^{\top }\mathbf{f}_m}.\]
Due to this formula, BigClam is a so called \emph{Bernoulli-Poisson model} \ron{CITE}.
The second assumption is that the event of an edge between two nodes is independent.  As a result, the probability of the entire graph $([N],E)$ conditioned on the AF is
\[\footnotesize
\begin{split}
   &  P(E|\mathbf{F}) =\\
   & \prod_{n \in [N]}{
        \sqrt{\prod_{m \in \mathcal{N}(n)}P( (n,m)|\mathbf{f}_n, \mathbf{f}_m)\prod_{m \notin \mathcal{N}(n)}{P(\neg (n,m))|\mathbf{f}_n, \mathbf{f}_m)}
        }}. \label{eqn:likelyhood}
\end{split}
\]
\textcolor{red}{Give number only to the second row. If you don't cite this formula, don't give it a label.}
Here, the square root is taken because the product considers each edge twice. 

Given a simple graph $([N],E)$, BigClam encodes the graph to the affiliation space by maximizing the log of likelihood with respect to the affiliation features $\mathbf{F}$
\begingroup
\footnotesize
\begin{align}
\label{eqn:naiveloss}
     l(\mathbf{F}) &= \frac{1}{2}\sumgraph \Big(\sumneighbors{}{\log(1-e^{-\mathbf{f}_n^{\top } \mathbf{f}_m})} - 
 \sumstrangers{\mathbf{f}_n^{\top } \mathbf{f}_m}\Big).
 \end{align}
 \endgroup
BigClam is optimized by gradient descent, with update
    \begin{align}
    \label{eqn:vanilla_optimization}
    F_u^{t+1} = F_u^\top  + \nabla_{F_u}l(F^\top _u) 
    \end{align}
    \textcolor{red}{fix notations.}
In order to implement the above interation with $O(|E|)$ operations at each step, instead of $O(N^2)$, the loss can be rearranged as
\begingroup
\footnotesize
\begin{align}
\label{eqn:comp_loss}
   &  l(\mathbf{F})=\\
   & \frac{1}{2}\bigg(\sumgraphu \Big(\sumneighbors{\log(e^{F_u^\top  F_v} - 1)} \hspace{3mm} - F_u^\top {\sumgraph{F_v}} + |F_u|^2\Big)
    \bigg)
\end{align}
\endgroup
The gradient of the loss is
\begin{align}
\label{eqn:BCupdate}
    % \nabla_{F_u}l = 
    % \sumneighbors{F_v \frac{e^{\inr[F_u]{F_v}}}{e^{\inr[F_u]{F_v}-1}}} - \sumgraph{F_v} + F_u \\
    \nabla_{F_u}l = 
    \sumneighbors{F_v \big(1-e^{-F_u^\top  F_v}
    \big)^{-1}} - \sumgraph{F_v} + F_u .
\end{align}
Since the global term needs to be calculated only once, the number of operations is  $O(E)$ instead of $O(N^2)$.  

We observe that the optimization process is a message passing scheme.
Looking at the dynamics of the optimization process, we see that every node is pushed in the direction of a weighted average of all of its neighbors, and pushed in the opposite direction by an average of all of the nodes. The sum of both forces tends to drive communities towards the axes.

\subsection{Normalizing Flows}


Normalizing flows is a deep learning algorithm which guarantees that the resulting approximated function is invertible, and provides the inverse.

The algorithm optimizes a reparameterization of $\mathbb{R}^d$, also referred to as "flow", denoted by $Z(F)$. The distribution in the $Z$ space is fixed to be some well known function, usually a gaussian. The transformation $Z$ is then trained to maximize the probability of the transformed training data. We denote $Z_\theta \equiv Z(\cdot | \theta)$ to be a parametric model with parameters $\theta$.  
From the constraint that the norm of the probability function is $1$ in either space we get: 
\begin{align}
\label{eqn:trans_prob}
    p_F(F) = p_Z(Z_\theta(F)) \big|\det(\partial_F Z_\theta(F)\big| 
\end{align}
Where  $\partial_F Z(F)$ is the Jacobian of the transformation with respect to the input.

A composition of flows is also a flow with a Jacobian that is a matrix product of the the two composed flows, so the formula for any composition of flows simply becomes:
\begin{align}
\label{eqn:trans_prob_composed}
    p_F(F) = p_Z(\circ_{i=0}^L Z^i_{\theta^i}(F)) \prod^L_{i=0}\big|\det(\partial_F Z^i_{\theta^i}(F)\big|  
\end{align}

In the framework of deep learning, a flow is created by an affine transformation whose parameters are outputs of a neural network.
%\subsubsection{Real NVP [\cite{dinh2016density}]}

The Real NVP (Real Valued Non-volume preserving \cite{dinh2016density}) is composed of a concatenation of so called coupling blocks. In a coupling block the input $Z$ is split into two parts: $Z_A$ and $Z_B$, and each part is treated differently: $Z_A$ is fed into a coupling network $\theta(\cdot) = (s(\cdot), t(\cdot))$ (e.g. an MLP). 
\begin{align}
{Z_A}^{i+1} &= {Z_A}^i \\
{Z_B}^{i+1} &= s({Z_A}^i){Z_B}^i + t({Z_A}^i)    
\end{align}
In between each of the blocks, a permutation is inserted to shuffle the elements of $Z^i$.

As long as the affine transformation is non zero, the coupling block is invertible and is therefore a "flow".



\textcolor{red}{Cite the regularity lemma, and motivate the proof. (if there is no space, we will move this to the appendix)} %\textcolor{cyan}{that is what i didn't understand - the motivation. we want to say that under some conditions our model is a global approximator? }

\subsection{Cut KL Regularity Proof}

\begin{definition}
  The \emph{cut norm} of a matrix $\mathbf{X}\in\mathbb{R}^{N\times N}$ is defined to be
    \begin{equation}
    \label{eq:GraphCutNorm0}  \|\mathbf{X}\|_{\square}:=\frac{1}{N^2}\sup_{\mathcal{U},\mathcal{V}\subset[N]}\Big|\sum_{i\in \mathcal{U}}\sum_{j\in \mathcal{V}}x_{i,j}\Big|.
\end{equation}  
\end{definition}
%The cut norm can be thought of  as the largest square that can be produced in an adjacency matrix.
\textcolor{red}{Explain that there is a probabilistic interpretation of cut norm.}
%\textcolor{cyan}{dan thoughts}
% For a matrix of dyadic probabilities $P(\cdot,\cdot)$ as in our case, the cut norm signifies the sets.... for a matrix of probabilities all of the entries are positive so U and V are just [N]... the interesting things happen in the cut distance.... 
%\textcolor{cyan}{dan thoughts}
The measure of accuracy we use is defined to be
%\[E(A,P):=\]
% \[\frac{1}{N^2}\sup_{\mathcal{U},\mathcal{V}\subset[N]}\Big|\log\Big(\prod_{i\in \mathcal{U}}\prod_{j\in \mathcal{V}}
%\frac{1-P(F_u,F_v)}{1-a_{u,v}}\Big)\Big|,\]
\[E(A,P):=\]
\[\inf_{0\leq d<1} \Bigg( d+ \frac{1}{N^2}\sup_{\mathcal{U},\mathcal{V}\subset[N]}\Big|\log\Big(\prod_{i\in \mathcal{U}}\prod_{j\in \mathcal{V}}
\frac{1-P(F_u,F_v)}{1-(1-d)a_{u,v}}\Big)\Big|\Bigg).\]
The second term is the cut distance $\| X-Q\|_{\square}$ between the matrix $X$ with entries $x_{u,v}=-\log(1-(1-d)a_{u,v})$ and the matrix $Q$
 with entries $q_{u,v}=-\log(1-P(F_u,F_v))$. Namely, the cut distance between the log likelihoods of non-edges. The parameter $d$ makes the $\{0,1\}$-valued adjacency weights $a_{u,v}$ amenable to modelling via exponents. Note that the probabilistic model $1-\exp(-F_uBF_v)$ cannot attain probability 1 since $F_uBF_v$ is real. Hence, the goal of $d$ is to regularize the probability of the edges of the graph $A$, where higher regularization is penalized via the additive term $d$ in the definition of $E(A,P)$.

 The distance $E$ is related to the Kullback-Leibler (KL) divergence. Namely, for each choice of $U,V\subset[N]$, 
 \[\log\Big(\prod_{i\in \mathcal{U}}\prod_{j\in \mathcal{V}}
\frac{1-P(F_u,F_v)}{1-(1-d)a_{u,v}}\Big)\]
is an un-normalized (for density functions not summing to 1) KL divergence between the non-edge probability of the graph and the non-edge probability of the IeClam graph over the edges between  $U$ and $V$, where edges are drawn ``unnormalized-uniformly'' with probability $1/N^2$. The unnormalized uniform distribution discourages the error $E$ from choosing small blocks for maximizing (**),. Note that normalized uniform distributions would lead $E$ to choose  small blocks, which do not reflect meaningful empirical estimates of the statistics  of the edges (the edge densities of small blocks would not be interpretable as expected number of edges). Hence, $E(A,P)$ is interpreted as the maximal KL divergence between $A$ and $P$ over all blocks (up to the best regularizer $d$ plus regularization panelty).

\ron{Say that we can also add $(1-b)$ to the top part when defining for general probability measures, but in our case this is unnecessary.}
 


We do not use the above as the optimization target in practice, since finding the supremum is NP hard. Instead, we replace it by the log likelihood described in (**). %the $L_2$ norm between the log likelihoods. This is justifies by the constructive weak regularity lemma, which asserts that minimizing the Frobenius error (which is the L2 norm over the entries of the matrix) leads to a small error in cut norm.

%This analysis also explains the approximation theoretic grounds of BigClam.

Next we define a general decoder that estimates edges by operating on pairs of points in a code space.

\begin{definition}
   A pairwaise decoder over the code spaces $\mathbb{R}^M$ is a mapping $D_M:\mathbb{R}^{2M}\rightarrow[0,1]$. Given $N$ points in the code space $\mathbf{z}=\{z_n\in\mathbb{R}^M\}_{n=1}^N$, the decoded graph $G_N(\mathbf{z})$ is the weighted graph with adjacency matrix 
   \[\mathbf{D}_M(\mathbf{z})=\big(D_M(z_n,z_k)\big)_{n,k=1}^N.\]
\end{definition}

The pairwise deconder of BigClam and IeClam involves a construction with an exponential, which cannot attain probability 1.


%\textbf{We don't need this. For now I keep it, since we may use it in the proof:}

%Clam models cannot predict edges with weight 1 due to the exponential construction. Since adjacency matrixes of unweighted graphs do have weights $1$, we have to limit the weight is a weaker sense. We propose the following assumption on the densities of edge weights, rather than individual edge weights. .

%\begin{definition}
 %   Let $R:(0,1)\rightarrow \mathbb{R}_+$ such that $R_{\epsilon}$ goes to infinity as $\epsilon$ goes to zero. 
%
%An adjacency matrix $\mathbf{A}\in[0,1]^{N\times N}$ is called \emph{weight-regular} (with respect to $R$) if for every $\epsilon>0$ there exists a matrix $A_{R_{\epsilon}}\in[0,1-\exp(-R_{\epsilon})]^{N\times N}$ such that
%\[\|\mathbf{A}_{R_{\epsilon}}-\mathbf{A}\|_{\square} < \epsilon.\]
%\end{definition}

%This bound is trivially correct  for a fixed $N$. Indeed, one can replace all the edge weights above $1-\exp(-R_{\epsilon})$ by $1-\exp(-R_{\epsilon})$, and get an error
%\[\|\mathbf{A}_{R_{\epsilon}}-\mathbf{A}\|_{\square} \leq \|\mathbf{A}_{R_{\epsilon}}-\mathbf{A}\|_1 \leq \exp(-R_{\epsilon}).\]

%So just take $\exp(-R_{\epsilon})=\epsilon$, or
%\[R_{\epsilon}=-\log(\epsilon).\]
%Then, the intersecting community regularity lemma will require $K=\frac{-18\log(\epsilon)}{4\epsilon^2}$, which is good asymptotics. 

%\[\]

\[\]

%\textbf{Continuing:}

The motivation in the following definition is that we would like to uniformly choose the dimension of the code space and the decoder, such that every graph can be approximated by decoding some points in the code space. Namely, we would like one \emph{universal} decoder that works for all graph, as opposed to choosing the dimension of the code space depending on the graph.

\begin{definition}
    A family of code spaces $\mathbb{R}^M$ and corresponding pairwaise decoders $D_M:\mathbb{R}^{2M}\rightarrow[0,1]$, parametrized by $M\in\mathbb{N}$, is called  \emph{universal} if for every $\epsilon>0$ there is $M\in\mathbb{N}$ (that depends just on $\epsilon$) such that for every $N\in\mathbb{N}$ and every graph  with adjacency matrix $\mathbf{A}$ and $N$ nodes there are $N$ points in the code space $\{z_n\in\mathbb{R}^M\}_{n=1}^N$ such that 
    \[E\big(\mathbf{A},\mathbf{D}_M(\mathbf{z})\big)<\epsilon\]
\end{definition}

\textcolor{red}{After this, show that BigClam is not universal. For full bipartite graph, you need at least as many dimensions as half the number of nodes, so $M$ cannot depend just on $\epsilon$.}

%The following is a slight reformulation of \ref{weakReg}. It is a direct result of Lemma 3.1 in \citep{Szemeredi_analyst}, for the special case where the graphon in induced by a graph.
%\begin{theorem}[Weak Regularity Lemma \cite{weakReg}]
%\label{Theorem:WRL}
%    For every every graph $G=([N],\mathcal{E})$ with adjacency matrix $\mathbf{A}\in\mathbb{R}^{N\times N}$ and every $\epsilon>0$  , there is a partition $\mathcal{P}=\{\mathcal{N}_1,\ldots,\mathcal{N}_k\}$ of $[N]$ into $k\leq 2^{2/\epsilon^2}$ classes, and graph with block adjacency matrix $\mathbf{B}\in\mathbb{R}^{N\times N}$, with blocks $\{\mathcal{N}_i\times \mathcal{N}_j\}_{i,j=1}^k$, such that $\| \mathbf{A}- \mathbf{B} \|_{\square} \leq \epsilon$. 
%\end{theorem}

\[\]

For IeClam, the decoder is
\[P(F_u,F_v) = 1-\exp\big({-F_u^\top BF_v}\big)\]
Denote the decoded matrix by $P$.

\[\]



\begin{theorem}
    Let $\epsilon>0$. For $K=\frac{-18\log(\epsilon)}{\epsilon^2}$, there is a community affiliation space with $K$ inclusive and $K$ exclusive communities, such that any adjacency matrix $\mathbf{A}$ can be approximated by a decoded IeClam matrix $P$ with error $E(A,P)<\epsilon$.
\end{theorem}

As a corollary, IeClam is universal.

$ $

\textbf{Proof of universality of IeClam}

$ $

Define ${\rm diag}(\mathbf{r})\in\mathbb{R}^{N\times N}$ as the diagonal matrix with diagonal elements $\mathbf{r}\in\mathbb{R}^N$.



\begin{definition}
    A (hard) intersecting community graph (ICG) with $N$ nodes and  $K$ communities is a matrix $\mathbf{C}$ of the following form. There exist $K$ column vectors $\mathbf{Q}=\big(\mathbf{q}_k\in\{0,1\}^N\big))_{k=1}^K\in \{0,1\}^{N\times K}$ and $k$ coefficients $\mathbf{r}=(r_k\in\mathbb{R})_{k=1}^K\in\mathbb{R}^K$ such that
    \[\mathbf{C}= \mathbf{Q}{\rm diag}(\mathbf{r})\mathbf{Q}^{\top }.\]
\end{definition}

The following is a special case of the weak regularity lemma from **. Up to the small modification of the adjacency matrix, allowing it to have values in $[0,R]$, it is well known as a step in the proof of the weak regularity lemma, see, .e.g  [Frieze, Fast. , Lovasz. Analyst]. For completness, we prove it in Appendix ***.

\begin{theorem}
%\label{thm:reg_sprase}
   % Let $1\leq p\leq2$. % and $D\in\sN$.
     Let $\mathbf{A}\in[0,R]^{N\times N}$ be an adjacency matrix of a graph with $N$ nodes. Let $\epsilon>0$. Denote $K=\frac{9R^2}{4\epsilon^2}$.         Then, there exists a hard ICG $\mathbf{C}$ with $K$ communities such that 
    \begin{equation}
        \label{eq:reg_sparse}
    \|\mathbf{A}-\mathbf{C}\|_{\square} \leq \epsilon.
    \end{equation}
\end{theorem}

\[\]

Let $\epsilon>0$. 
Let $\mathbf{A}$. Let $0<d\leq 1$.

Consider the matrix  $\mathbf{X}$ with entries
\[x_{u,v}=-\log(1-(1-d)a_{u,v}).\]



In the construction, we build regularity features $F$ and we want
\[1-\exp(-F_uBF_v^{\top })\approx (1-d)a_{u,v}.\]
So the regularity lemma should be against the matrix $\mathbf{X}$ with entries
\[x_{u,v}=-\log(1-(1-d)a_{u,v}).\]
Note that $-\log(1-(1-d)a_{u,v})$ is increasing in $(1-d)a_{u,v}$. For $(1-d)a_{u,v}=0$ the value of this function is 0, and for $(1-d)a_{u,v}=1-d$ it is $R=-\log(d)$. Choose $d=\epsilon$. For this specific choice of $d$, if we replace $d$ by $\epsilon$ in the definition of $E$ and omit the infimum, we get an upper bound of $E$.
 

We now use the overlapping weak regularity lemma.


Usind the overlaping weak regularity lemma, we approximate $X$ by an ICG with 
\[K=\frac{-18\log(\epsilon)}{4\epsilon^2}\]
communities,
\[\mathbf{C}= \mathbf{Q}{\rm diag}(\mathbf{r})\mathbf{Q}^{\top },\]
such that
\[\|\mathbf{X}-\mathbf{C}\|_{\square} \leq \epsilon.\]

Let $\mathbf{r}_+={\rm ReLU}(\mathbf{r})$ and $\mathbf{r}_-={\rm ReLU}(-\mathbf{r})$. Denote
\[\mathbf{C}_+=\mathbf{Q}{\rm diag}(\sqrt{\mathbf{r}_+})\]
and
\[\mathbf{C}_-=\mathbf{Q}{\rm diag}(\sqrt{\mathbf{r}_-}).\]
Denote the rows of $\mathbf{C}_+$ by $\mathbf{S}_n$ and  the rows of $\mathbf{C}_+$ by $\mathbf{T}_n$, for $n=1,\ldots,N$. We then concatenate them to the feature $F_n$, with the first $K$ coordinates being $\mathbf{S}_n$ (the inclusive communities) and the last $K$ coordinates being $\mathbf{T}_n$ (the exclusive communities).

It is easy to see that $F_uBF_v$ is the $(u,v)$ enrty of $Q{\rm diag}(r)Q^{\top }$. This proves that
\[\|-\log(1-(1-\epsilon)\mathbf{A})+\log(1-P)\|_{\square}\]
\[=\|\mathbf{X}-\mathbf{C}\|_{\square} \leq \epsilon.\]
The claim now follows by adding $\epsilon$ to the above formula, and then replacing $\epsilon$ by $\epsilon/2$.

In more detail, our  metric is bounded as follows
\[E(A,P) \leq \]
\[\epsilon+\frac{1}{N^2}\sup_{\mathcal{U},\mathcal{V}\subset[N]}\Big|\log\Big(\prod_{i\in \mathcal{U}}\prod_{j\in \mathcal{V}}
\frac{1-P(F_u,F_v)}{1-(1-\epsilon)a_{u,v}}\Big)\Big|\]
\[ =
\epsilon+\frac{1}{N^2}\sup_{\mathcal{U},\mathcal{V}\subset[N]}\Big|\sum_{i\in \mathcal{U}}\sum_{j\in \mathcal{V}}\Big(-\log\big(1-(1-\epsilon)a_{u,v}\big) \]
\[+ \log\big(1-P(F_u,F_v)\big)\Big)\Big|\]
\[=\epsilon + \| C-X \|_{\square} = 2\epsilon. \]
We now replace $\epsilon$ by $\epsilon/2$.

\[\]
An SBM is $Z C Z^\top $ where $Z$ is a community affiliation vector in the shape $N \times K$, typically one hot,  and $C$ is a matrix of dimensions $K\times K$. We prove that any SBM matrix C can be factorized using inclusive exclusive model.
\begin{enumerate}
 \item Consider an SBM $C$. 
\item     If $C \succ 0$ can be factorized $C = Q^\top  Q$ where Q is orthonormal. (homophilic model) 
\item If $C \prec 0$ it is still symmetric and therefore has SVD $C = Q^\top  S Q$ where S is diagonal.  so:
\tiny
\begin{align}
    C &= Q^\top  \begin{pmatrix}
        \sqrt{|s_{11}|} & 0 & \ldots \\
        0 & \sqrt{|s_{22}|} & \ldots \\
        \vdots & \vdots & \vdots
    \end{pmatrix} \begin{pmatrix}
        \text{sign}(s_{11}) & 0 & \ldots \\
        0 & \text{sign}(s_{22}) & \ldots \\
        0 & 0 & \ldots
    \end{pmatrix} \begin{pmatrix}
        \sqrt{|s_{11}|} & 0 & \ldots \\
        0 & \sqrt{|s_{22}|} & \ldots \\
        \vdots & \vdots & \vdots
    \end{pmatrix} Q \\
    &= F^\top  B F
\end{align}
\tiny
 \end{enumerate}
 $S$ elements can now be ordered according to magnitude and kept to principal components.
\[\]


\[\]

\subsection{Experiments}
\subsubsection{Ablation Study}
In this section we test the necessity of all of the components used in the optimization:
- Noise addition
- Amplitude halving
- prior 
- back and forth optimization

\paragraph{Real World Implementation}

\paragraph{Early stop on BP probability}
An algorithm that might work in the real world can connect early stop to the clam probabilities. We can run each of the fit functions until the clam probability saturates and do an early stop on that.

\paragraph{hyper choice}
1. try to find an iter count that works for all datasets.

2. if doesn't work try early stopping with link prediction (or log likelihood) loss decide for k parameters and the delta of improvement. either do hyperparameter optimization or see with naked eye. number of non improvements k, and delta of improvement.
when using link prediction, there are several options: - check the auc of all of the graph after several iterations. 

3. densification: The average degree of a graph is an indicator on whether to use densification or not.  

\clearpage

 a general guide for experienced \LaTeX{} users. If you do not know how to use \LaTeX{}, please obtain assistance locally. AAAI cannot provide you with support and the accompanying style files are \textbf{not} guaranteed to work. If the results you obtain are not in accordance with the specifications you received, you must correct your source file to achieve the correct result.

These instructions are generic. Consequently, they do not include specific dates, page charges, and so forth. Please consult your specific written conference instructions for details regarding your submission. Please review the entire document for specific instructions that might apply to your particular situation. All authors must comply with the following:

\begin{itemize}
\item You must use the 2025 AAAI Press \LaTeX{} style file and the aaai25.bst bibliography style files, which are located in the 2025 AAAI Author Kit (aaai25.sty, aaai25.bst).
\item You must complete, sign, and return by the deadline the AAAI copyright form (unless directed by AAAI Press to use the AAAI Distribution License instead).
\item You must read and format your paper source and PDF according to the formatting instructions for authors.
\item You must submit your electronic files and abstract using our electronic submission form \textbf{on time.}
\item You must pay any required page or formatting charges to AAAI Press so that they are received by the deadline.
\item You must check your paper before submitting it, ensuring that it compiles without error, and complies with the guidelines found in the AAAI Author Kit.
\end{itemize}

\section{Copyright}
All papers submitted for publication by AAAI Press must be accompanied by a valid signed copyright form. They must also contain the AAAI copyright notice at the bottom of the first page of the paper. There are no exceptions to these requirements. If you fail to provide us with a signed copyright form or disable the copyright notice, we will be unable to publish your paper. There are \textbf{no exceptions} to this policy. You will find a PDF version of the AAAI copyright form in the AAAI AuthorKit. Please see the specific instructions for your conference for submission details.

\section{Formatting Requirements in Brief}
We need source and PDF files that can be used in a variety of ways and can be output on a variety of devices. The design and appearance of the paper is strictly governed by the aaai style file (aaai25.sty).
\textbf{You must not make any changes to the aaai style file, nor use any commands, packages, style files, or macros within your own paper that alter that design, including, but not limited to spacing, floats, margins, fonts, font size, and appearance.} AAAI imposes requirements on your source and PDF files that must be followed. Most of these requirements are based on our efforts to standardize conference manuscript properties and layout. All papers submitted to AAAI for publication will be recompiled for standardization purposes. Consequently, every paper submission must comply with the following requirements:

\begin{itemize}
\item Your .tex file must compile in PDF\LaTeX{} --- (you may not include .ps or .eps figure files.)
\item All fonts must be embedded in the PDF file --- including your figures.
\item Modifications to the style file, whether directly or via commands in your document may not ever be made, most especially when made in an effort to avoid extra page charges or make your paper fit in a specific number of pages.
\item No type 3 fonts may be used (even in illustrations).
\item You may not alter the spacing above and below captions, figures, headings, and subheadings.
\item You may not alter the font sizes of text elements, footnotes, heading elements, captions, or title information (for references and mathematics, please see the limited exceptions provided herein).
\item You may not alter the line spacing of text.
\item Your title must follow Title Case capitalization rules (not sentence case).
\item \LaTeX{} documents must use the Times or Nimbus font package (you may not use Computer Modern for the text of your paper).
\item No \LaTeX{} 209 documents may be used or submitted.
\item Your source must not require use of fonts for non-Roman alphabets within the text itself. If your paper includes symbols in other languages (such as, but not limited to, Arabic, Chinese, Hebrew, Japanese, Thai, Russian and other Cyrillic languages), you must restrict their use to bit-mapped figures. Fonts that require non-English language support (CID and Identity-H) must be converted to outlines or 300 dpi bitmap or removed from the document (even if they are in a graphics file embedded in the document).
\item Two-column format in AAAI style is required for all papers.
\item The paper size for final submission must be US letter without exception.
\item The source file must exactly match the PDF.
\item The document margins may not be exceeded (no overfull boxes).
\item The number of pages and the file size must be as specified for your event.
\item No document may be password protected.
\item Neither the PDFs nor the source may contain any embedded links or bookmarks (no hyperref or navigator packages).
\item Your source and PDF must not have any page numbers, footers, or headers (no pagestyle commands).
\item Your PDF must be compatible with Acrobat 5 or higher.
\item Your \LaTeX{} source file (excluding references) must consist of a \textbf{single} file (use of the ``input" command is not allowed.
\item Your graphics must be sized appropriately outside of \LaTeX{} (do not use the ``clip" or ``trim'' command) .
\end{itemize}

If you do not follow these requirements, your paper will be returned to you to correct the deficiencies.

\section{What Files to Submit}
You must submit the following items to ensure that your paper is published:
\begin{itemize}
\item A fully-compliant PDF file.
\item Your \LaTeX{} source file submitted as a \textbf{single} .tex file (do not use the ``input" command to include sections of your paper --- every section must be in the single source file). (The only allowable exception is .bib file, which should be included separately).
\item The bibliography (.bib) file(s).
\item Your source must compile on our system, which includes only standard \LaTeX{} 2020 TeXLive support files.
\item Only the graphics files used in compiling paper.
\item The \LaTeX{}-generated files (e.g. .aux,  .bbl file, PDF, etc.).
\end{itemize}

Your \LaTeX{} source will be reviewed and recompiled on our system (if it does not compile, your paper will be returned to you. \textbf{Do not submit your source in multiple text files.} Your single \LaTeX{} source file must include all your text, your bibliography (formatted using aaai25.bst), and any custom macros.

Your files should work without any supporting files (other than the program itself) on any computer with a standard \LaTeX{} distribution.

\textbf{Do not send files that are not actually used in the paper.} Avoid including any files not needed for compiling your paper, including, for example, this instructions file, unused graphics files, style files, additional material sent for the purpose of the paper review, intermediate build files and so forth.
\textbf{Obsolete style files.} The commands for some common packages (such as some used for algorithms), may have changed. Please be certain that you are not compiling your paper using old or obsolete style files.

\textbf{Final Archive.} Place your source files in a single archive which should be compressed using .zip. The final file size may not exceed 10 MB.
Name your source file with the last (family) name of the first author, even if that is not you.


\section{Using \LaTeX{} to Format Your Paper}

The latest version of the AAAI style file is available on AAAI's website. Download this file and place it in the \TeX\ search path. Placing it in the same directory as the paper should also work. You must download the latest version of the complete AAAI Author Kit so that you will have the latest instruction set and style file.

\subsection{Document Preamble}

In the \LaTeX{} source for your paper, you \textbf{must} place the following lines as shown in the example in this subsection. This command set-up is for three authors. Add or subtract author and address lines as necessary, and uncomment the portions that apply to you. In most instances, this is all you need to do to format your paper in the Times font. The helvet package will cause Helvetica to be used for sans serif. These files are part of the PSNFSS2e package, which is freely available from many Internet sites (and is often part of a standard installation).

Leave the setcounter for section number depth commented out and set at 0 unless you want to add section numbers to your paper. If you do add section numbers, you must uncomment this line and change the number to 1 (for section numbers), or 2 (for section and subsection numbers). The style file will not work properly with numbering of subsubsections, so do not use a number higher than 2.

\subsubsection{The Following Must Appear in Your Preamble}
\begin{quote}
\begin{scriptsize}\begin{verbatim}
\documentclass[letterpaper]{article}
% DO NOT CHANGE THIS
\usepackage{aaai25} % DO NOT CHANGE THIS
\usepackage{times} % DO NOT CHANGE THIS
\usepackage{helvet} % DO NOT CHANGE THIS
\usepackage{courier} % DO NOT CHANGE THIS
\usepackage[hyphens]{url} % DO NOT CHANGE THIS
\usepackage{graphicx} % DO NOT CHANGE THIS
\urlstyle{rm} % DO NOT CHANGE THIS
\def\UrlFont{\rm} % DO NOT CHANGE THIS
\usepackage{graphicx}  % DO NOT CHANGE THIS
\usepackage{natbib}  % DO NOT CHANGE THIS
\usepackage{caption}  % DO NOT CHANGE THIS
\frenchspacing % DO NOT CHANGE THIS
\setlength{\pdfpagewidth}{8.5in} % DO NOT CHANGE THIS
\setlength{\pdfpageheight}{11in} % DO NOT CHANGE THIS
%
% Keep the \pdfinfo as shown here. There's no need
% for you to add the /Title and /Author tags.
\pdfinfo{
/TemplateVersion (2025.1)
}
\end{verbatim}\end{scriptsize}
\end{quote}

\subsection{Preparing Your Paper}

After the preamble above, you should prepare your paper as follows:
\begin{quote}
\begin{scriptsize}\begin{verbatim}
\begin{document}
\maketitle
\begin{abstract}
%...
\end{abstract}\end{verbatim}\end{scriptsize}
\end{quote}

\noindent If you want to add links to the paper's code, dataset(s), and extended version or similar this is the place to add them, within a \emph{links} environment:
\begin{quote}%
\begin{scriptsize}\begin{verbatim}
\begin{links}
  \link{Code}{https://aaai.org/example/guidelines}
  \link{Datasets}{https://aaai.org/example/datasets}
  \link{Extended version}{https://aaai.org/example}
\end{links}\end{verbatim}\end{scriptsize}
\end{quote}

\noindent You should then continue with the body of your paper. Your paper must conclude with the references, which should be inserted as follows:
\begin{quote}
\begin{scriptsize}\begin{verbatim}
% References and End of Paper
% These lines must be placed at the end of your paper
\bibliography{Bibliography-File}
\end{document}
\end{verbatim}\end{scriptsize}
\end{quote}


\begin{quote}
\begin{scriptsize}\begin{verbatim}
\begin{document}\\
\maketitle\\
...\\
\bibliography{Bibliography-File}\\
\end{document}\\
\end{verbatim}\end{scriptsize}
\end{quote}

\subsection{Commands and Packages That May Not Be Used}
\begin{table*}[t]
\centering

\begin{tabular}{l|l|l|l}
\textbackslash abovecaption &
\textbackslash abovedisplay &
\textbackslash addevensidemargin &
\textbackslash addsidemargin \\
\textbackslash addtolength &
\textbackslash baselinestretch &
\textbackslash belowcaption &
\textbackslash belowdisplay \\
\textbackslash break &
\textbackslash clearpage &
\textbackslash clip &
\textbackslash columnsep \\
\textbackslash float &
\textbackslash input &
\textbackslash input &
\textbackslash linespread \\
\textbackslash newpage &
\textbackslash pagebreak &
\textbackslash renewcommand &
\textbackslash setlength \\
\textbackslash text height &
\textbackslash tiny &
\textbackslash top margin &
\textbackslash trim \\
\textbackslash vskip\{- &
\textbackslash vspace\{- \\
\end{tabular}
%}
\caption{Commands that must not be used}
\label{table1}
\end{table*}

\begin{table}[t]
\centering
%\resizebox{.95\columnwidth}{!}{
\begin{tabular}{l|l|l|l}
    authblk & babel & cjk & dvips \\
    epsf & epsfig & euler & float \\
    fullpage & geometry & graphics & hyperref \\
    layout & linespread & lmodern & maltepaper \\
    navigator & pdfcomment & pgfplots & psfig \\
    pstricks & t1enc & titlesec & tocbind \\
    ulem
\end{tabular}
\caption{LaTeX style packages that must not be used.}
\label{table2}
\end{table}

There are a number of packages, commands, scripts, and macros that are incompatable with aaai25.sty. The common ones are listed in tables \ref{table1} and \ref{table2}. Generally, if a command, package, script, or macro alters floats, margins, fonts, sizing, linespacing, or the presentation of the references and citations, it is unacceptable. Note that negative vskip and vspace may not be used except in certain rare occurances, and may never be used around tables, figures, captions, sections, subsections, subsubsections, or references.


\subsection{Page Breaks}
For your final camera ready copy, you must not use any page break commands. References must flow directly after the text without breaks. Note that some conferences require references to be on a separate page during the review process. AAAI Press, however, does not require this condition for the final paper.


\subsection{Paper Size, Margins, and Column Width}
Papers must be formatted to print in two-column format on 8.5 x 11 inch US letter-sized paper. The margins must be exactly as follows:
\begin{itemize}
\item Top margin: .75 inches
\item Left margin: .75 inches
\item Right margin: .75 inches
\item Bottom margin: 1.25 inches
\end{itemize}


The default paper size in most installations of \LaTeX{} is A4. However, because we require that your electronic paper be formatted in US letter size, the preamble we have provided includes commands that alter the default to US letter size. Please note that using any other package to alter page size (such as, but not limited to the Geometry package) will result in your final paper being returned to you for correction.


\subsubsection{Column Width and Margins.}
To ensure maximum readability, your paper must include two columns. Each column should be 3.3 inches wide (slightly more than 3.25 inches), with a .375 inch (.952 cm) gutter of white space between the two columns. The aaai25.sty file will automatically create these columns for you.

\subsection{Overlength Papers}
If your paper is too long and you resort to formatting tricks to make it fit, it is quite likely that it will be returned to you. The best way to retain readability if the paper is overlength is to cut text, figures, or tables. There are a few acceptable ways to reduce paper size that don't affect readability. First, turn on \textbackslash frenchspacing, which will reduce the space after periods. Next, move all your figures and tables to the top of the page. Consider removing less important portions of a figure. If you use \textbackslash centering instead of \textbackslash begin\{center\} in your figure environment, you can also buy some space. For mathematical environments, you may reduce fontsize {\bf but not below 6.5 point}.


Commands that alter page layout are forbidden. These include \textbackslash columnsep,  \textbackslash float, \textbackslash topmargin, \textbackslash topskip, \textbackslash textheight, \textbackslash textwidth, \textbackslash oddsidemargin, and \textbackslash evensizemargin (this list is not exhaustive). If you alter page layout, you will be required to pay the page fee. Other commands that are questionable and may cause your paper to be rejected include \textbackslash parindent, and \textbackslash parskip. Commands that alter the space between sections are forbidden. The title sec package is not allowed. Regardless of the above, if your paper is obviously ``squeezed" it is not going to to be accepted. Options for reducing the length of a paper include reducing the size of your graphics, cutting text, or paying the extra page charge (if it is offered).


\subsection{Type Font and Size}
Your paper must be formatted in Times Roman or Nimbus. We will not accept papers formatted using Computer Modern or Palatino or some other font as the text or heading typeface. Sans serif, when used, should be Courier. Use Symbol or Lucida or Computer Modern for \textit{mathematics only. }

Do not use type 3 fonts for any portion of your paper, including graphics. Type 3 bitmapped fonts are designed for fixed resolution printers. Most print at 300 dpi even if the printer resolution is 1200 dpi or higher. They also often cause high resolution imagesetter devices to crash. Consequently, AAAI will not accept electronic files containing obsolete type 3 fonts. Files containing those fonts (even in graphics) will be rejected. (Authors using blackboard symbols must avoid packages that use type 3 fonts.)

Fortunately, there are effective workarounds that will prevent your file from embedding type 3 bitmapped fonts. The easiest workaround is to use the required times, helvet, and courier packages with \LaTeX{}2e. (Note that papers formatted in this way will still use Computer Modern for the mathematics. To make the math look good, you'll either have to use Symbol or Lucida, or you will need to install type 1 Computer Modern fonts --- for more on these fonts, see the section ``Obtaining Type 1 Computer Modern.")

If you are unsure if your paper contains type 3 fonts, view the PDF in Acrobat Reader. The Properties/Fonts window will display the font name, font type, and encoding properties of all the fonts in the document. If you are unsure if your graphics contain type 3 fonts (and they are PostScript or encapsulated PostScript documents), create PDF versions of them, and consult the properties window in Acrobat Reader.

The default size for your type must be ten-point with twelve-point leading (line spacing). Start all pages (except the first) directly under the top margin. (See the next section for instructions on formatting the title page.) Indent ten points when beginning a new paragraph, unless the paragraph begins directly below a heading or subheading.


\subsubsection{Obtaining Type 1 Computer Modern for \LaTeX{}.}

If you use Computer Modern for the mathematics in your paper (you cannot use it for the text) you may need to download type 1 Computer fonts. They are available without charge from the American Mathematical Society:
http://www.ams.org/tex/type1-fonts.html.

\subsubsection{Nonroman Fonts.}
If your paper includes symbols in other languages (such as, but not limited to, Arabic, Chinese, Hebrew, Japanese, Thai, Russian and other Cyrillic languages), you must restrict their use to bit-mapped figures.

\subsection{Title and Authors}
Your title must appear centered over both text columns in sixteen-point bold type (twenty-four point leading). The title must be written in Title Case according to the Chicago Manual of Style rules. The rules are a bit involved, but in general verbs (including short verbs like be, is, using, and go), nouns, adverbs, adjectives, and pronouns should be capitalized, (including both words in hyphenated terms), while articles, conjunctions, and prepositions are lower case unless they directly follow a colon or long dash. You can use the online tool \url{https://titlecaseconverter.com/} to double-check the proper capitalization (select the "Chicago" style and mark the "Show explanations" checkbox).

Author's names should appear below the title of the paper, centered in twelve-point type (with fifteen point leading), along with affiliation(s) and complete address(es) (including electronic mail address if available) in nine-point roman type (the twelve point leading). You should begin the two-column format when you come to the abstract.

\subsubsection{Formatting Author Information.}
Author information has to be set according to the following specification depending if you have one or more than one affiliation. You may not use a table nor may you employ the \textbackslash authorblk.sty package. For one or several authors from the same institution, please separate them with commas and write all affiliation directly below (one affiliation per line) using the macros \textbackslash author and \textbackslash affiliations:

\begin{quote}\begin{scriptsize}\begin{verbatim}
\author{
    Author 1, ..., Author n\\
}
\affiliations {
    Address line\\
    ... \\
    Address line\\
}
\end{verbatim}\end{scriptsize}\end{quote}


\noindent For authors from different institutions, use \textbackslash textsuperscript \{\textbackslash rm x \} to match authors and affiliations. Notice that there should not be any spaces between the author name and the superscript (and the comma should come after the superscripts).

\begin{quote}\begin{scriptsize}\begin{verbatim}
\author{
    AuthorOne\equalcontrib\textsuperscript{\rm 1,\rm2},
    AuthorTwo\equalcontrib\textsuperscript{\rm 2},
    AuthorThree\textsuperscript{\rm 3},\\
    AuthorFour\textsuperscript{\rm 4},
    AuthorFive\textsuperscript{\rm 5}}
}
\affiliations {
    \textsuperscript{\rm 1}AffiliationOne,\\
    \textsuperscript{\rm 2}AffiliationTwo,\\
    \textsuperscript{\rm 3}AffiliationThree,\\
    \textsuperscript{\rm 4}AffiliationFour,\\
    \textsuperscript{\rm 5}AffiliationFive\\
    \{email, email\}@affiliation.com,
    email@affiliation.com,
    email@affiliation.com,
    email@affiliation.com
}
\end{verbatim}\end{scriptsize}\end{quote}

You can indicate that some authors contributed equally using the \textbackslash equalcontrib command. This will add a marker after the author names and a footnote on the first page.

Note that you may want to  break the author list for better visualization. You can achieve this using a simple line break (\textbackslash  \textbackslash).

\subsection{\LaTeX{} Copyright Notice}
The copyright notice automatically appears if you use aaai25.sty. It has been hardcoded and may not be disabled.

\subsection{Credits}
Any credits to a sponsoring agency should appear in the acknowledgments section, unless the agency requires different placement. If it is necessary to include this information on the front page, use
\textbackslash thanks in either the \textbackslash author or \textbackslash title commands.
For example:
\begin{quote}
\begin{small}
\textbackslash title\{Very Important Results in AI\textbackslash thanks\{This work is
 supported by everybody.\}\}
\end{small}
\end{quote}
Multiple \textbackslash thanks commands can be given. Each will result in a separate footnote indication in the author or title with the corresponding text at the botton of the first column of the document. Note that the \textbackslash thanks command is fragile. You will need to use \textbackslash protect.

Please do not include \textbackslash pubnote commands in your document.

\subsection{Abstract}
Follow the example commands in this document for creation of your abstract. The command \textbackslash begin\{abstract\} will automatically indent the text block. Please do not indent it further. {Do not include references in your abstract!}

\subsection{Page Numbers}

Do not print any page numbers on your paper. The use of \textbackslash pagestyle is forbidden.

\subsection{Text}
The main body of the paper must be formatted in black, ten-point Times Roman with twelve-point leading (line spacing). You may not reduce font size or the linespacing. Commands that alter font size or line spacing (including, but not limited to baselinestretch, baselineshift, linespread, and others) are expressly forbidden. In addition, you may not use color in the text.

\subsection{Citations}
Citations within the text should include the author's last name and year, for example (Newell 1980). Append lower-case letters to the year in cases of ambiguity. Multiple authors should be treated as follows: (Feigenbaum and Engelmore 1988) or (Ford, Hayes, and Glymour 1992). In the case of four or more authors, list only the first author, followed by et al. (Ford et al. 1997).

\subsection{Extracts}
Long quotations and extracts should be indented ten points from the left and right margins.

\begin{quote}
This is an example of an extract or quotation. Note the indent on both sides. Quotation marks are not necessary if you offset the text in a block like this, and properly identify and cite the quotation in the text.

\end{quote}

\subsection{Footnotes}
Use footnotes judiciously, taking into account that they interrupt the reading of the text. When required, they should be consecutively numbered throughout with superscript Arabic numbers. Footnotes should appear at the bottom of the page, separated from the text by a blank line space and a thin, half-point rule.

\subsection{Headings and Sections}
When necessary, headings should be used to separate major sections of your paper. Remember, you are writing a short paper, not a lengthy book! An overabundance of headings will tend to make your paper look more like an outline than a paper. The aaai25.sty package will create headings for you. Do not alter their size nor their spacing above or below.

\subsubsection{Section Numbers.}
The use of section numbers in AAAI Press papers is optional. To use section numbers in \LaTeX{}, uncomment the setcounter line in your document preamble and change the 0 to a 1. Section numbers should not be used in short poster papers and/or extended abstracts.

\subsubsection{Section Headings.}
Sections should be arranged and headed as follows:
\begin{enumerate}
\item Main content sections
\item Appendices (optional)
\item Ethical Statement (optional, unnumbered)
\item Acknowledgements (optional, unnumbered)
\item References (unnumbered)
\end{enumerate}

\subsubsection{Appendices.}
Any appendices must appear after the main content. If your main sections are numbered, appendix sections must use letters instead of arabic numerals. In \LaTeX{} you can use the \texttt{\textbackslash appendix} command to achieve this effect and then use \texttt{\textbackslash section\{Heading\}} normally for your appendix sections.

\subsubsection{Ethical Statement.}
You can write a statement about the potential ethical impact of your work, including its broad societal implications, both positive and negative. If included, such statement must be written in an unnumbered section titled \emph{Ethical Statement}.

\subsubsection{Acknowledgments.}
The acknowledgments section, if included, appears right before the references and is headed ``Acknowledgments". It must not be numbered even if other sections are (use \texttt{\textbackslash section*\{Acknowledgements\}} in \LaTeX{}). This section includes acknowledgments of help from associates and colleagues, credits to sponsoring agencies, financial support, and permission to publish. Please acknowledge other contributors, grant support, and so forth, in this section. Do not put acknowledgments in a footnote on the first page. If your grant agency requires acknowledgment of the grant on page 1, limit the footnote to the required statement, and put the remaining acknowledgments at the back. Please try to limit acknowledgments to no more than three sentences.

\subsubsection{References.}
The references section should be labeled ``References" and must appear at the very end of the paper (don't end the paper with references, and then put a figure by itself on the last page). A sample list of references is given later on in these instructions. Please use a consistent format for references. Poorly prepared or sloppy references reflect badly on the quality of your paper and your research. Please prepare complete and accurate citations.

\subsection{Illustrations and  Figures}

\begin{figure}[t]
\centering
\includegraphics[width=0.9\columnwidth]{figure1} % Reduce the figure size so that it is slightly narrower than the column. Don't use precise values for figure width.This setup will avoid overfull boxes.
\caption{Using the trim and clip commands produces fragile layers that can result in disasters (like this one from an actual paper) when the color space is corrected or the PDF combined with others for the final proceedings. Crop your figures properly in a graphics program -- not in LaTeX.}
\label{fig1}
\end{figure}

\begin{figure*}[t]
\centering
\includegraphics[width=0.8\textwidth]{figure2} % Reduce the figure size so that it is slightly narrower than the column.
\caption{Adjusting the bounding box instead of actually removing the unwanted data resulted multiple layers in this paper. It also needlessly increased the PDF size. In this case, the size of the unwanted layer doubled the paper's size, and produced the following surprising results in final production. Crop your figures properly in a graphics program. Don't just alter the bounding box.}
\label{fig2}
\end{figure*}

% Using the \centering command instead of \begin{center} ... \end{center} will save space
% Positioning your figure at the top of the page will save space and make the paper more readable
% Using 0.95\columnwidth in conjunction with the


Your paper must compile in PDF\LaTeX{}. Consequently, all your figures must be .jpg, .png, or .pdf. You may not use the .gif (the resolution is too low), .ps, or .eps file format for your figures.

Figures, drawings, tables, and photographs should be placed throughout the paper on the page (or the subsequent page) where they are first discussed. Do not group them together at the end of the paper. If placed at the top of the paper, illustrations may run across both columns. Figures must not invade the top, bottom, or side margin areas. Figures must be inserted using the \textbackslash usepackage\{graphicx\}. Number figures sequentially, for example, figure 1, and so on. Do not use minipage to group figures.

If you normally create your figures using pgfplots, please create the figures first, and then import them as pdfs with proper bounding boxes, as the bounding and trim boxes created by pfgplots are fragile and not valid.

When you include your figures, you must crop them \textbf{outside} of \LaTeX{}. The command \textbackslash includegraphics*[clip=true, viewport 0 0 10 10]{...} might result in a PDF that looks great, but the image is \textbf{not really cropped.} The full image can reappear (and obscure whatever it is overlapping) when page numbers are applied or color space is standardized. Figures \ref{fig1}, and \ref{fig2} display some unwanted results that often occur.

If your paper includes illustrations that are not compatible with PDF\TeX{} (such as .eps or .ps documents), you will need to convert them. The epstopdf package will usually work for eps files. You will need to convert your ps files to PDF in either case.

\subsubsection {Figure Captions.}The illustration number and caption must appear \textit{under} the illustration. Labels and other text with the actual illustration must be at least nine-point type. However, the font and size of figure captions must be 10 point roman. Do not make them smaller, bold, or italic. (Individual words may be italicized if the context requires differentiation.)

\subsection{Tables}

Tables should be presented in 10 point roman type. If necessary, they may be altered to 9 point type. You must not use \texttt{\textbackslash resizebox} or other commands that resize the entire table to make it smaller, because you can't control the final font size this way.
If your table is too large you can use \texttt{\textbackslash setlength\{\textbackslash tabcolsep\}\{1mm\}} to compress the columns a bit or you can adapt the content (e.g.: reduce the decimal precision when presenting numbers, use shortened column titles, make some column duble-line to get it narrower).

Tables that do not fit in a single column must be placed across double columns. If your table won't fit within the margins even when spanning both columns and using the above techniques, you must split it in two separate tables.

\subsubsection {Table Captions.} The number and caption for your table must appear \textit{under} (not above) the table.  Additionally, the font and size of table captions must be 10 point roman and must be placed beneath the figure. Do not make them smaller, bold, or italic. (Individual words may be italicized if the context requires differentiation.)



\subsubsection{Low-Resolution Bitmaps.}
You may not use low-resolution (such as 72 dpi) screen-dumps and GIF files---these files contain so few pixels that they are always blurry, and illegible when printed. If they are color, they will become an indecipherable mess when converted to black and white. This is always the case with gif files, which should never be used. The resolution of screen dumps can be increased by reducing the print size of the original file while retaining the same number of pixels. You can also enlarge files by manipulating them in software such as PhotoShop. Your figures should be 300 dpi when incorporated into your document.

\subsubsection{\LaTeX{} Overflow.}
\LaTeX{} users please beware: \LaTeX{} will sometimes put portions of the figure or table or an equation in the margin. If this happens, you need to make the figure or table span both columns. If absolutely necessary, you may reduce the figure, or reformat the equation, or reconfigure the table.{ \bf Check your log file!} You must fix any overflow into the margin (that means no overfull boxes in \LaTeX{}). \textbf{Nothing is permitted to intrude into the margin or gutter.}


\subsubsection{Using Color.}
Use of color is restricted to figures only. It must be WACG 2.0 compliant. (That is, the contrast ratio must be greater than 4.5:1 no matter the font size.) It must be CMYK, NOT RGB. It may never be used for any portion of the text of your paper. The archival version of your paper will be printed in black and white and grayscale. The web version must be readable by persons with disabilities. Consequently, because conversion to grayscale can cause undesirable effects (red changes to black, yellow can disappear, and so forth), we strongly suggest you avoid placing color figures in your document. If you do include color figures, you must (1) use the CMYK (not RGB) colorspace and (2) be mindful of readers who may happen to have trouble distinguishing colors. Your paper must be decipherable without using color for distinction.

\subsubsection{Drawings.}
We suggest you use computer drawing software (such as Adobe Illustrator or, (if unavoidable), the drawing tools in Microsoft Word) to create your illustrations. Do not use Microsoft Publisher. These illustrations will look best if all line widths are uniform (half- to two-point in size), and you do not create labels over shaded areas. Shading should be 133 lines per inch if possible. Use Times Roman or Helvetica for all figure call-outs. \textbf{Do not use hairline width lines} --- be sure that the stroke width of all lines is at least .5 pt. Zero point lines will print on a laser printer, but will completely disappear on the high-resolution devices used by our printers.

\subsubsection{Photographs and Images.}
Photographs and other images should be in grayscale (color photographs will not reproduce well; for example, red tones will reproduce as black, yellow may turn to white, and so forth) and set to a minimum of 300 dpi. Do not prescreen images.

\subsubsection{Resizing Graphics.}
Resize your graphics \textbf{before} you include them with LaTeX. You may \textbf{not} use trim or clip options as part of your \textbackslash includegraphics command. Resize the media box of your PDF using a graphics program instead.

\subsubsection{Fonts in Your Illustrations.}
You must embed all fonts in your graphics before including them in your LaTeX document.

\subsubsection{Algorithms.}
Algorithms and/or programs are a special kind of figures. Like all illustrations, they should appear floated to the top (preferably) or bottom of the page. However, their caption should appear in the header, left-justified and enclosed between horizontal lines, as shown in Algorithm~\ref{alg:algorithm}. The algorithm body should be terminated with another horizontal line. It is up to the authors to decide whether to show line numbers or not, how to format comments, etc.

In \LaTeX{} algorithms may be typeset using the {\tt algorithm} and {\tt algorithmic} packages, but you can also use one of the many other packages for the task.

\begin{algorithm}[tb]
\caption{Example algorithm}
\label{alg:algorithm}
\textbf{Input}: Your algorithm's input\\
\textbf{Parameter}: Optional list of parameters\\
\textbf{Output}: Your algorithm's output
\begin{algorithmic}[1] %[1] enables line numbers
\STATE Let $t=0$.
\WHILE{condition}
\STATE Do some action.
\IF {conditional}
\STATE Perform task A.
\ELSE
\STATE Perform task B.
\ENDIF
\ENDWHILE
\STATE \textbf{return} solution
\end{algorithmic}
\end{algorithm}

\subsubsection{Listings.}
Listings are much like algorithms and programs. They should also appear floated to the top (preferably) or bottom of the page. Listing captions should appear in the header, left-justified and enclosed between horizontal lines as shown in Listing~\ref{lst:listing}. Terminate the body with another horizontal line and avoid any background color. Line numbers, if included, must appear within the text column.

\begin{listing}[tb]%
\caption{Example listing {\tt quicksort.hs}}%
\label{lst:listing}%
\begin{lstlisting}[language=Haskell]
quicksort :: Ord a => [a] -> [a]
quicksort []     = []
quicksort (p:xs) = (quicksort lesser) ++ [p] ++ (quicksort greater)
	where
		lesser  = filter (< p) xs
		greater = filter (>= p) xs
\end{lstlisting}
\end{listing}

\subsection{References}
The AAAI style includes a set of definitions for use in formatting references with BibTeX. These definitions make the bibliography style fairly close to the ones  specified in the Reference Examples appendix below. To use these definitions, you also need the BibTeX style file ``aaai25.bst," available in the AAAI Author Kit on the AAAI web site. Then, at the end of your paper but before \textbackslash end{document}, you need to put the following lines:

\begin{quote}
\begin{small}
\textbackslash bibliography\{bibfile1,bibfile2,...\}
\end{small}
\end{quote}

Please note that the aaai25.sty class already sets the bibliographystyle for you, so you do not have to place any \textbackslash bibliographystyle command in the document yourselves. The aaai25.sty file is incompatible with the hyperref and navigator packages. If you use either, your references will be garbled and your paper will be returned to you.

References may be the same size as surrounding text.
However, in this section (only), you may reduce the size to {\em \textbackslash small} (9pt) if your paper exceeds the allowable number of pages. Making it any smaller than 9 point with 10 point linespacing, however, is not allowed.

The list of files in the \textbackslash bibliography command should be the names of your BibTeX source files (that is, the .bib files referenced in your paper).

The following commands are available for your use in citing references:
\begin{quote}
{\em \textbackslash cite:} Cites the given reference(s) with a full citation. This appears as ``(Author Year)'' for one reference, or ``(Author Year; Author Year)'' for multiple references.\smallskip\\
{\em \textbackslash shortcite:} Cites the given reference(s) with just the year. This appears as ``(Year)'' for one reference, or ``(Year; Year)'' for multiple references.\smallskip\\
{\em \textbackslash citeauthor:} Cites the given reference(s) with just the author name(s) and no parentheses.\smallskip\\
{\em \textbackslash citeyear:} Cites the given reference(s) with just the date(s) and no parentheses.
\end{quote}
You may also use any of the \emph{natbib} citation commands.


\section{Proofreading Your PDF}
Please check all the pages of your PDF file. The most commonly forgotten element is the acknowledgements --- especially the correct grant number. Authors also commonly forget to add the metadata to the source, use the wrong reference style file, or don't follow the capitalization rules or comma placement for their author-title information properly. A final common problem is text (expecially equations) that runs into the margin. You will need to fix these common errors before submitting your file.

\section{Improperly Formatted Files }
In the past, AAAI has corrected improperly formatted files submitted by the authors. Unfortunately, this has become an increasingly burdensome expense that we can no longer absorb). Consequently, if your file is improperly formatted, it will be returned to you for correction.

\section{Naming Your Electronic File}
We require that you name your \LaTeX{} source file with the last name (family name) of the first author so that it can easily be differentiated from other submissions. Complete file-naming instructions will be provided to you in the submission instructions.

\section{Submitting Your Electronic Files to AAAI}
Instructions on paper submittal will be provided to you in your acceptance letter.

\section{Inquiries}
If you have any questions about the preparation or submission of your paper as instructed in this document, please contact AAAI Press at the address given below. If you have technical questions about implementation of the aaai style file, please contact an expert at your site. We do not provide technical support for \LaTeX{} or any other software package. To avoid problems, please keep your paper simple, and do not incorporate complicated macros and style files.

\begin{quote}
\noindent AAAI Press\\
1101 Pennsylvania Ave, NW Suite 300\\
Washington, DC 20004 USA\\
\textit{Telephone:} 1-202-360-4062\\
\textit{E-mail:} See the submission instructions for your particular conference or event.
\end{quote}

\section{Additional Resources}
\LaTeX{} is a difficult program to master. If you've used that software, and this document didn't help or some items were not explained clearly, we recommend you read Michael Shell's excellent document (testflow doc.txt V1.0a 2002/08/13) about obtaining correct PS/PDF output on \LaTeX{} systems. (It was written for another purpose, but it has general application as well). It is available at www.ctan.org in the tex-archive.

\appendix
\section{Reference Examples}
\label{sec:reference_examples}

\nobibliography*
Formatted bibliographies should look like the following examples. You should use BibTeX to generate the references. Missing fields are unacceptable when compiling references, and usually indicate that you are using the wrong type of entry (BibTeX class).

\paragraph{Book with multiple authors~\nocite{em:86}} Use the \texttt{@book} class.\\[.2em]
\bibentry{em:86}.

\paragraph{Journal and magazine articles~\nocite{r:80, hcr:83}} Use the \texttt{@article} class.\\[.2em]
\bibentry{r:80}.\\[.2em]
\bibentry{hcr:83}.

\paragraph{Proceedings paper published by a society, press or publisher~\nocite{c:83, c:84}} Use the \texttt{@inproceedings} class. You may abbreviate the \emph{booktitle} field, but make sure that the conference edition is clear.\\[.2em]
\bibentry{c:84}.\\[.2em]
\bibentry{c:83}.

\paragraph{University technical report~\nocite{r:86}} Use the \texttt{@techreport} class.\\[.2em]
\bibentry{r:86}.

\paragraph{Dissertation or thesis~\nocite{c:79}} Use the \texttt{@phdthesis} class.\\[.2em]
\bibentry{c:79}.

\paragraph{Forthcoming publication~\nocite{c:21}} Use the \texttt{@misc} class with a \texttt{note="Forthcoming"} annotation.
\begin{quote}
\begin{footnotesize}
\begin{verbatim}
@misc(key,
  [...]
  note="Forthcoming",
)
\end{verbatim}
\end{footnotesize}
\end{quote}
\bibentry{c:21}.

\paragraph{ArXiv paper~\nocite{c:22}} Fetch the BibTeX entry from the "Export Bibtex Citation" link in the arXiv website. Notice it uses the \texttt{@misc} class instead of the \texttt{@article} one, and that it includes the \texttt{eprint} and \texttt{archivePrefix} keys.
\begin{quote}
\begin{footnotesize}
\begin{verbatim}
@misc(key,
  [...]
  eprint="xxxx.yyyy",
  archivePrefix="arXiv",
)
\end{verbatim}
\end{footnotesize}
\end{quote}
\bibentry{c:22}.

\paragraph{Website or online resource~\nocite{c:23}} Use the \texttt{@misc} class. Add the url in the \texttt{howpublished} field and the date of access in the \texttt{note} field:
\begin{quote}
\begin{footnotesize}
\begin{verbatim}
@misc(key,
  [...]
  howpublished="\url{http://...}",
  note="Accessed: YYYY-mm-dd",
)
\end{verbatim}
\end{footnotesize}
\end{quote}
\bibentry{c:23}.

For the most up to date version of the AAAI reference style, please consult the \textit{AI Magazine} Author Guidelines at \url{https://aaai.org/ojs/index.php/aimagazine/about/submissions#authorGuidelines}

\section{Acknowledgments}
AAAI is especially grateful to Peter Patel Schneider for his work in implementing the original aaai.sty file, liberally using the ideas of other style hackers, including Barbara Beeton. We also acknowledge with thanks the work of George Ferguson for his guide to using the style and BibTeX files --- which has been incorporated into this document --- and Hans Guesgen, who provided several timely modifications, as well as the many others who have, from time to time, sent in suggestions on improvements to the AAAI style. We are especially grateful to Francisco Cruz, Marc Pujol-Gonzalez, and Mico Loretan for the improvements to the Bib\TeX{} and \LaTeX{} files made in 2020.

The preparation of the \LaTeX{} and Bib\TeX{} files that implement these instructions was supported by Schlumberger Palo Alto Research, AT\&T Bell Laboratories, Morgan Kaufmann Publishers, The Live Oak Press, LLC, and AAAI Press. Bibliography style changes were added by Sunil Issar. \verb+\+pubnote was added by J. Scott Penberthy. George Ferguson added support for printing the AAAI copyright slug. Additional changes to aaai25.sty and aaai25.bst have been made by Francisco Cruz, Marc Pujol-Gonzalez, and Mico Loretan.

\bigskip
\noindent Thank you for reading these instructions carefully. We look forward to receiving your electronic files!



\end{document}




\section{Ron's proof notes for nonuniversality of BigClam}


$f_1,\ldots,f_N$ should be close to orthogonal for pairwise distance close to zero. But we use cut norm.

The matrix $X$ in the first part is $f_n^\top f_m$. The target matrix $Y$ in this part is $0$. Cut norm is bounded by cut norm in this subgraph. So cut norm of $\{f_n^\top f_m\}_{n,m=1}^N$. Since this is non-negative (the AS is $\mathbb{R}_+^C$), the cut norm is just the $L_1$ norm. Namely,
\[\norm{X-Y}_{\square} = \frac{1}{N^2}\sum_{n,m=1}^N | f_n^\top f_m |.\]
Wince we now work in $L_1$, in the following we flatten the set of inner products to a sequence of $[N^2]$ points.

To make this small, we need the average inner product between two $f_n,f_m$ to be less than $\epsilon$. A necessary condition is that there are $N$ pairs of indices, WLOG the first ones $[N]\subset [N^2]$ (up to reordering), such that  $f_n^\top f_m \leq N\epsilon$ for every $(n,m)\in[N]$ (we formally define $[x]$ of a real positive $x$ the $[n]$ where $n$ is the floor of $x$). 
Otherwise, for every set of $N$ indices (of pairs), one inner product  is larger than $N\epsilon$. We can partition $N^2$ into $N$ such disjoint sets, and obtain $N$ inner products larger than $N\epsilon$, so the total $L_1$ norm of $\{f_n^\top f_m\}_{n,m=1}^N$ is larger than $\epsilon$.

The only way to achieve this (the existence of such $N$ inner products), for small enough $\epsilon$, is to have at least $C \geq \sqrt{N}$ axes if $\epsilon$ is small enough. Indeed, if $f_1,\ldots f_{\sqrt{N}}\in \mathbb{R}_+^{C<\sqrt{N}}$ have inner product less than $N\epsilon$, and if $N<\frac{1}{\sqrt{\epsilon}}$, so the error is less than $\sqrt{\epsilon}$, and $\sqrt{\epsilon}< ...$ then 

[Here, for every $\epsilon$, choose $N$ not too big, and we need - actually the dimension $D$ would depend on $\epsilon$, which is not good. I need an analysis that allows $N$ to be big independently of $\epsilon$. So, I need the infinity norm to lower bound the 1 norm. But this would be a bound $N^2\epsilon$ on individual inner products, which is not helpful.]

[MAke the above construction with different powers of $N$? There must be a set of $N^r$ inner products such that all of them are less than $N^{r}\epsilon$. Otherwise, partition the inner products to disjoint $N^{2-r}$ sets, and in each one there is an inner product larger than $N^{r}\epsilon$, so the total normalized L1 sum is larger than $N^{-2}N^{2-r} N^r\epsilon =\epsilon$.]

[Short: I write $e:=\epsilon$.]

[Given $\epsilon$ and $N$, we choose $r$ such that $N^{r}\epsilon=0.5$, namely, $r=\log(1/2\epsilon)/\log(N)$. For $r>0$ we do need a constraint on the size of $N$ - it has to satisfy $r=\log(1/e-N)>0$ so $1/e-N>1$ so $N<1/e-1$. So the choice of $D$ which depends on $N$, can actually be written as depending on $e$.]

[Maybe write the L1 norm as expected value of inner product, this will help in showing it cannot be small for bounded from below norm and bounded dimension and high $N$?]

Expected inner product of uniformly sampled points in the positive quadrant:
\[E(ip) = \int_{[0,1]^D}\int_{[0,1]^D} x^\top y dxdy\]
\[=D\int_0^1xdx\int_0^1ydy = D/4.\]
We may want to define inner product noramlzied by $D$.

The problem is that these are not indpependent, they are given from a set times set. See what I did in the paper with Ben in this situation.

Suppose I showed that for uniformly indepndently sampled points in the quadrant, the expedted value of the sum of inner products divided by $N^2$ is $D/4$. 

By a law of large numbers, or central limit/ Monte Carlo, given $N$ sample features, the empirical expected inner product (average, namely L1 normalized by $1/N^2$)  is close to $D/4$.

But is there a probability distribtion in which it can be significantly smaller than $D/4$? My intuition is that $D/4$ is the optimal one. 


Actually, this is bilinear, an linear WRT one side. We can write (This works since all is non-negative!)
\[L1 = N^{-2}\sum_{n,m} q_n^T q_m = \sum_m(\sum_{n} q_n^T)q_m\]
\[=v^Tv\]
where $v=N^{-1}\sum_n q_n$.

This makes things easy. the average inner product is just the norm of the average feature.

So the cut distance of one side is just the norm square of the average.

((But if we exclude product with ones self? This can get to zero in $D$ large enough.
\[\sum_{n}\sum_{m\neq n} q_n^T q_m\]
\[v^Tv - \sum_n q_n^Tq_n.\]
For now ignore this, but we can deal with it.
))

Now we want to get a small $v^Tv$. So the critical point is to show that $\|q_n\|>r$ must be bounded from below.
If we can show it, then 
\[v^Tv = (N^{-1}\sum_n q_n^T)(N^{-1}\sum_n q_n) >\]
What is the smallest it can be? 

If $D>N$ the smallest is if each $q_n$ is on a different axis. Then we get $v^Tv>r^2/N^2>r^2/D^2$. Otherwise, can I show that regardless of the dimension, we have $v^Tv>r^2/D^2$? This will not help, since then we can choose $D$ with respect to $\epsilon$.
We need to get a bound the depends on $N$, and a norm that increases in $N$. 

In the limit as $N$ increases, what is the distribution in the sphere positive quadrant that maximizes expected ...

Suppose we shrank each $q_n$ to have norm $r$. This can only reduce the norm of the sum

\[\| N^{-1}\sum q_n \|^2 > \]

Geometrically, any average point of point on the sphere must be in the convex hall of the sphere. 
The point in the convex hall of the sphere of radius $r$ which as lost norm is $(1,1,1,\ldots,1)r/D$. It has norm $r^2/D^{1/2}$. But, as I saied above, this does not help (we can choose $D$ as a function of $\epsilon$ and not $N$). The only way to get to this point is to put all $q_n$ orthogonal, but $N>>D$ so it is impossible.

\textbf{Actually, this is achievable. Just put an even number of points on each axis. The average will be $(1,1,1,\ldots,1)r/D$.}

So I can make one side have small $L_1$ norm, even if $r=1$. 
But the smallest possible $L_1$ norm will be  $D^{-1/2}$. 
So this does not lead to a contradiction.

So what about the bipartite part? This will make things fail?

Suppose that we fixed $D$. Many vectors from part 1 are close to orthogonal. So there is no way to make each vector for side 2 have high inner product with all of them.

Let's analyze the cut norm of the between sides.

The target is $\{0,a^2C\}$ valued. In the between sides it is constant $a^2C$.
Suppose we limit ourselves to features bounded by $a$. So cut distance of this part is just L1 norm, which is the sum without absolute value
\[a^2C-N^{-2}\sum_{n,m} q_ny_m\]
where $q$ and $y$ are the two sides. This is just
\[a^2C-q^Ty\]
where 
\[q=N^{-2}\sum_n q_n, \quad y=N^{-2}\sum_n y_n.\]
(correct the whole analysis to have $(4N)^{-2}$ since the cut norm is on both sides together..) (also then try not to bound from above the features.)

Show that cut norm is bounded from above by $1/4$ times the three cut norms norms of side 1, side 2, and twice the bipartite side. And each one can is equal to an L1 norm of entris of constant sign.
So I want to bound the sum of these $L_1$ norms from below.

So our problem is to minimize
\[R=q^Tq+y^Ty +a^2D - 2q^Ty \]
where $q,y$ are in the positive quadrant.
Is it possible to make this small? No, since
\[0<(q-y)^T(q-y) = q^Tq+y^Ty - 2q^Ty \]
so
\[R>a^2D.\]

{\Large Make a proof out of the above}

This is weird, since the problem is solvable if $D$ is not uniform? This is the thing - with self-loops the problem is never solvable!

The problem with this analysis is that we limit the size of the feature, and make the target maximal.

\textbf{Can I show that being larger than the value of the target is not helpful for lowering cut distance?}

\textbf{Yes:}

{\Large \textbf{Bound the cut norm from below by the sum without absolute value.} Now, we get the exact same minimization target, which cannot be minimized!}



[[In appendix - give the solution if $D=N^2$ for BigClam with no self-loops.]]



I need a refined analysis.



\[\]

This is not enough, I also need to assume that all of the features have norms bounded from below. Otherwise small inner product can be achived by small norms. So, why is it true that they have norms bounded from below?

If their norms are smaller than something, then the norms of the othere side must be very very big, and we simply do the current analysis on the other side.

So there is a lower bound, which can be taken $\sqrt{a}$. Since the inner product with the other side must be $a$ (Maybe we need $a^e$ for some $e>0$ - keep it this general, since we work with cut norm...).

So suppose each $\|f_n\|$ is bounded from below by $a^e$.
In this case, the inner product in the quadrant is related to  the distance by
\[\|\frac{f_1}{\|f_1\|}-\frac{f_2}{\|f_2\|}\|^2 = 2 -2 f_1^T f_2/\|f_1\|^2 \|f_2\|^2,\]
so
\[f_1^T f_2 =   2- \|f_1\|^2 \|f_2\|^2\|\frac{f_1}{\|f_1\|}-\frac{f_2}{\|f_2\|}\|^2 \]
\[\geq 2 - A^2\|\frac{f_1}{\|f_1\|}-\frac{f_2}{\|f_2\|}\|^2\]
where we also need to bound the features from above by some $A$ that depends on $a$.

Now, there is no way to put many points distant from each other in the sphere in the quadrant. But we do have curse of dimensionality.... So the asymptotics would not look nice?


Should the proof be easier? I want to show that there is no way to put $N$ points, bounded from below by some $a^e$, with error less than some $\delta$, if $N$ is chosen large enough.

Connet to ray packing. Ray packing in only for angles, but we know that the norm is bounded from below by $a^e$, so the inner product is bounded from below by the angle between the says (up to some cosine or something) times $a^{2e}$.


use the fact that the best way to put n nodes in n plus 1 or leas dimensions is eatf and theinner product is something known. negative 1 divided by n. does not help. I need aresult on the positive quadrant. Ray packing in positive quadrant.

just convert to the problem of ball packing on the sphere viaminimal norm and use the fact the the positivequadrant of the sphere is compact.

\[\]

\[\]

the best way to is a equiangular

For each choice, bound the cut norm from beforw by the following choice of $U,V$. Choose them 


\[\]






\[\]


\[\]


\[\]

\[\]